% ============================================================
%  R. Nielsen, "Evangelietroen og Theologien. Tolv Forelæsninger,
%  holdte ved Universitetet i Kjøbenhavn i Vinteren 1849--50."
%  Kjøbenhavn: C. A. Reitzels Forlag, 1850. VIII + 174 pp., Fraktur.
%
%  TRANSLATION (English) --- Hans Halvorson
%  Translated from transcription.tex (Danish), which is COMPLETE and
%  image-verified, pp. 1--174 + Forord III--VIII + Indhold.
%  Method: ../../../TRANSLATION-PLAYBOOK.md.
%
%  STATUS: in progress. Batch markers below, in reading order.
%
%  ---- TRANSLATION CONVENTIONS FOR THIS BOOK ----
%
%  The book turns on a small set of terms that Nielsen uses with technical
%  weight. These are fixed here and must be held constant throughout; the
%  argument is unreadable if they drift.
%
%    Evangelietroen        = the faith of the Gospel  (in the running title and
%                            the thesis restatements: "Gospel faith")
%    Troen                 = faith
%    Viden / Videnskaben   = knowledge / science  (Videnskab is Wissenschaft:
%                            "science" in the broad sense, incl. theology's own
%                            claim to be one --- never "natural science")
%    ueensartet med        = heterogeneous with   (THE thesis word; not
%                            "unlike", "different from", "incommensurable")
%    eensartet med         = homogeneous with
%    Speculation(en)       = speculation; speculativ = speculative
%    halvspeculativ        = half-speculative      (Nielsen's own coinage,
%                            13 occurrences --- see the near-miss note below)
%    halvphilosophisk      = half-philosophical
%    Existens              = existence  (Kierkegaardian; keep the noun)
%    Forargelse            = offense    (skandalon; not "annoyance")
%    Troeserkjendelse      = faith-knowledge / the knowledge of faith
%    Tilegnelse            = appropriation
%    Inderlighed           = inwardness
%    Bestaaende (det)      = the established order
%    Opvakte (vore)        = our awakened ones (the revivalist party)
%    at stride FOR Troen   = to strive FOR faith
%    at strides OM Troen   = to dispute ABOUT faith
%       --- the Forord's whole hinge. Keep the preposition contrast audible
%           in English every time it recurs.
%
%  - Quotes. Danish „...“ -> English curly doubles ``...''.
%  - Emphasis. \emph{} in the Danish is letterspacing (Sperrsatz) in the
%    print; carry every one over as \emph{}.
%  - Latin/foreign phrases in \textit{} stay \textit{} and stay untranslated
%    (reservatio mentalis, propter iniquitatem temporum, in republica
%    litterarum, fides religiosa, fides historica, in specie). German likewise
%    ("ein Aber dabei").
%  - Footnotes: \renewcommand{\thefootnote}{*)} matches the print, as in the
%    transcription. 13 footnotes in the book.
%  - Page markers "% --- p. N ---" are carried over at the same joints, so the
%    English can be collated against the Danish and the scan page by page.
%  - „en Anden“ (Forord pp. VII--VIII) is Kierkegaard, unnamed. Render as
%    ``another'' and DO NOT gloss --- the anonymity is the point.
%  - The Forord's opening block quotation is Kierkegaard/Anti-Climacus,
%    _Sygdommen til Døden_; Nielsen names only "a Christian psychologist who
%    knows ``The Sickness unto Death''". Englished here from Nielsen's Danish,
%    not swapped for an existing English Kierkegaard translation.
%  - Printer's defects are transcribed as printed in the Danish and logged in
%    % comments there. THE ENGLISH REPRODUCES THE SENSE, NOT THE DEFECT, in
%    every case and without exception: an unclosed quotation is closed where
%    the sense ends; a missing full stop is supplied; a misspelling
%    („Refomatorerne“, „lil“ for „til“, „medele“ for „meddele“) is spelled
%    correctly. A % note at each site records what the print actually has.
%    (Danish quote balance runs +2 across the book on purpose; the English is
%    expected to balance at 0 in every batch.)
%    Rationale: the Danish IS the diplomatic witness and carries every defect.
%    An English typo invented to "match" a Danish one is not a witness to
%    anything --- there is no 1850 English setting for it to be faithful to.
%    Two sites were drafted the other way and corrected; do not reintroduce.
%
%  ---- TWO TRAPS RECORDED IN RESUME-NOTES, REPEATED HERE ----
%  1. NEVER global-replace "halvspeculative"/"half-speculative" --- it is
%     Nielsen's own term in 13 places. Scope every heading fix to its block.
%  2. Sandbox compile: strip textalpha and map Greek to a placeholder, else
%     you get ~22 bogus "Unicode character not set up" errors. That is the
%     test rig, not the document. Never delete Greek from this file.
% ============================================================
\documentclass[12pt,a4paper]{book}
\usepackage[utf8]{inputenc}
\usepackage[T1]{fontenc}
\usepackage{libertinus}
\usepackage{libertinust1math}
\usepackage{textalpha}   % Greek typed directly
\usepackage{graphicx}
\usepackage[english]{babel}
\usepackage[protrusion=true,expansion=true]{microtype}
\usepackage{setspace}
\usepackage[margin=1.2in]{geometry}
\usepackage{enumitem}
\usepackage{fancyhdr}
\setlength{\headheight}{28pt}
\addtolength{\topmargin}{-13.5pt}

\renewcommand{\thefootnote}{*)}

\usepackage{hyperref}

\hypersetup{
  pdftitle={Gospel Faith and Theology},
  pdfauthor={R. Nielsen},
  hidelinks
}

\pagestyle{fancy}
\fancyhf{}
\fancyhead[LE,RO]{\thepage}
\fancyhead[RE]{\textit{Gospel Faith and Theology}}
\fancyhead[LO]{\textit{\leftmark}}

\setstretch{1.3}

\title{\textbf{Gospel Faith and Theology}\\[6pt]
  \large Twelve Lectures, Delivered at the University\\
  of Copenhagen in the Winter of 1849--50}
\author{R.\ Nielsen\\[4pt]
  \small Translated from the Danish by Hans Halvorson}
\date{Copenhagen: C.\ A.\ Reitzel, 1850}

\begin{document}
\maketitle
\thispagestyle{empty}
\newpage

\tableofcontents
\newpage

%% ============================================================
%%  FORORD / PREFACE (printed pp. III--VIII = PDF 8--13)
%% ============================================================
\chapter*{Preface}
\addcontentsline{toc}{chapter}{Preface}
\label{ch:forord}

% --- p. III ---
``If one assumes that all the many pastors here and abroad who deliver and
write sermons are believing Christians, how is it then to be explained that
one never hears or reads the prayer that lies so near at hand, especially in
our times: God in heaven, I thank Thee that Thou hast not required of a human
being that he should comprehend Christianity; for if it were required, then I
should be the most wretched of all. The more I seek to comprehend it, the more
incomprehensible it appears to me, the more I discover the possibility of
offense. Therefore I thank Thee that Thou requirest faith alone, and I pray
that Thou wilt increase it in me further. This prayer would be entirely correct
as orthodoxy, and, granted that it were true in the one praying, it would at
the same time be correct irony over the whole of speculation. But is faith to
be found upon earth!''

These words of a Christian psychologist who knows ``The Sickness unto Death''
I use as a preface to the present writing, partly because they have gripped me
deeply, partly also because they may best help me to a distinction without
which my lectures, in spite of all protests and deprecating assurances, will
hardly escape misinterpretation --- I mean the distinction between striving
for faith and disputing about faith.
% NB: no emphasis here in the print --- the Danish letterspaces neither
% preposition, and the contrast carries in English on "for"/"about" alone.
% Two \emph{} were added here in drafting and removed; do not re-add.
% (But see the open item: spacing.py flags a RUN "Men strides" on p. IV,
% i.e. „Men at strides om Troen er dog kun menneskeligt“ may be letterspaced
% and is not yet marked in the Danish. If that pass adds it there, mirror it
% in paragraph 6 here.)

To strive for faith is divine; to dispute about faith is human, indeed in part
worldly. The reader will thus already be able to discern from this why the
above words are not placed as a motto, but are simply set up as a preface: an
ingenious motto awakens an expectation; but this preface stands just here in
order to turn an expectation away.

However much I have striven in these lectures to avoid everything that could
in any way have the appearance of a merely personal quarrel between man and
man, I have nevertheless not been able to avoid coming so close to the
established order that the reader might, from his point of view, have a kind of
right to construe something personal in it. However much I have striven in
these lectures from
% --- p. IV ---
first to last to make it comprehensible that Christianity does not will to be
comprehended, I have nevertheless not been quite able to avoid the appearance
of setting a concept against a concept, a theory against a theory. For to
explain why one cannot explain is after all always to entrust oneself to an
explanation; to give grounds why no grounds can be given is, after all,
nevertheless to engage in a grounding.

The present lectures must therefore, in accordance with their design and
execution, naturally put up with being reckoned among that sort of attempt in
which one does not so much strive for faith as, properly speaking, dispute
about faith.

But to dispute about faith is, after all, only human; to strive for faith ---
that is in truth divine. The latter must on no account be confused with the
former; in order to forestall such a confusion, the distinction ought to be
brought out as definitely as possible.

As I now sat pondering this distinction and sought a telling expression, those
words came to my mind: ``\emph{I thank Thee, God in heaven, that Thou hast not
required of a human being that he should comprehend Christianity!}''; and the
words stirred in my memory and formed for me an image, a speaking image, the
image of a Lutheran pastor holding prayer in his congregation.

Since I then wished to present, as in an example, what difference there is
between disputing with grounds against grounds and forgetting all finite
grounds in a single decisive \emph{personal} ground; between disputing with
explanations against explanations and keeping all explanations in a single
all-explaining thought; in short, what difference there is between disputing
about faith and striving for faith: it then seemed to me that I could choose no
better example than the example of this praying pastor, who assuredly will not
dispute with anyone about faith, but will strive in all inwardness --- and yet
with all the world --- for faith.

That the striver's prayer must appear to an outside observer as a ``correct
irony over the whole of speculation'' is comprehensible; it is, after all, no
church prayer for ``Christian science,'' no prayer-day prayer for ``the
critical introduction'': the irony is then not in the words either, but in the
deed, not for the one praying, but for the reflecting observer. The prayer
itself is a simple and to that extent also ``unscientific'' prayer in faith for
faith; but the
% --- p. V ---
reflecting onlooker hears it at once, where the words of the one praying as it
were involuntarily betray a secret acquaintance with ``the whole of
speculation.'' How else should the pastor know within himself that he would be
``the most wretched of all'' if it were required of him that he should
comprehend Christianity; how should just this pastor, among so many
ratiocinating, criticizing, aestheticizing pastors who otherwise carry the
Church's ``dogmatic'' speculations without anyone's being able to notice on
them that the burden weighs; how --- I ask --- should just this one be able to
feel so deeply that on speculation's terms he would have to become --- ``the
most wretched!'' --- if he had not himself, and at his own risk, also once
allowed himself to look into ``the Concept,'' and then looked so deeply that
with horror he saw --- ``the possibility of offense.''

If, then, there really were everywhere by all pastors and in all congregations
such striving in faith for faith; if those days returned when the question of
the blessed life was the race's highest question; if the Church's great jubilee
year were now so near that it rang in our ears, as though we already heard the
peal of the bells that are to ring it in: yes, then our awakened ones could
easily enough dismiss a merely ``theoretical school-quarrel,'' and refer to
life's ``ordeal by fire'' those who --- ``instead of acting accordingly''! ---
only write about it, insofar as they write about faith. ``But'' --- adds the
Christian psychologist --- ``is faith to be found upon earth!''

If, however, in order to be rid of the ``school-quarrel'' together with its
investigations, one wishes to give the matter a more sober turn and declare
that the dispute of principle here announced does not after all signify much,
since the victory --- if we are finally to be quite candid --- has at bottom
already been decided beforehand in another place: then from my present
standpoint I can of course have not the least objection, but can only rejoice
that with so fine a candor one is at last disposed to give truth the honor.

The struggle in which the principle of faith can, with the most intensive power
of reflection, make war on dogmatizing speculation has assuredly been fought
out in a quite other and higher productivity than the one that can be offered
in these discourses. The victory that under the conditions at hand can be won
over the age's aesthetic, metaphysical, systematic, world-historical
``delusions of sense'' has --- so far as I can judge --- really already been
won by another, and that in the deepest silence, and that so comprehensively,
so completely, and on so extensive a
% --- p. VI ---
scale, that he who won such a victory might surely with right say to himself:
``we have more than conquered.''

All the same, there is here still, I suppose --- so far as I can see --- ``ein
Aber dabei.''

For with wars and victories things stand somewhat otherwise in the spiritual
than in the sensuous world. If a man in the world of the senses has been
bodily slain, anyone can easily convince himself that the fallen one really is
dead. But in the world of spirit it is otherwise. In the world of spirit a man
can sometimes be struck so mortally, and yet so deftly, that not even he
himself (much less so many others) rightly notices that anything has actually
happened to him. This will be the case precisely in proportion to the greatness
of the superiority before which the vanquished has fallen. In the bodily sense
it may indeed be ridiculous enough to treat a living man as a ghost, even
though rumor has falsely reported him dead; since sensuous appearance is of
course the most striking proof that the person is still alive. But in the world
of spirit it is otherwise. In the world of spirit, outward appearance and
phenomenal well-being stand, not so very seldom, in an --- alas! --- entirely
inverted relation to the individual's inner health and essential vital force;
here, then, what experience will also teach us holds above all: that appearance
deceives.

Since now the deceptive appearance is undeniably truth's most dangerous enemy,
it becomes in the world of ideas sometimes necessary that one not merely strive
for truth and right as in soundless silence, but that one also, with a somewhat
louder (and just for that reason somewhat coarser) kind of weapon, dispute at
the same time about who it actually is that has conquered in the struggle. If
one neglects this last too much, it might easily happen that the higher ideal
plenitude of power were --- at least for a time --- bungled, and a confusion
arose in the invisible republic, a confusion far more dangerous to the highest
and holiest interests than the confusion that arises in the visible when a
besieged fortress lacks a commandant.

So I conclude from this that it might not yet be so entirely without
significance to offer a dispute about the principle of faith-knowledge, and I
therefore beg the well-disposed reader to be content with my contributions.

How precarious my task nevertheless shows itself to be, everyone must surely
admit who thinks the matter over a little more closely.

% --- p. VII ---
If there really is to be a dispute here about who it is that has conquered in
the struggle, how near at hand does not the misinterpretation lie that the
alleged interest in the conception of a principle may after all serve only as
one of those pretexts with which a selfish passion would embellish a
presumptuous vanity.

Should this misinterpretation (which, however, I do not assume) prove able to
force its way through and establish itself as the authentic interpretation:
then I must ask that the accusation be advanced so openly and so entirely
without all \textit{reservatio mentalis} that it can expressly fall upon me
alone.

I put forward this request not in order to affect a heroism I do not possess,
but because I regard it as my duty to attempt everything that can in any way
contribute to keeping this cause as clean as possible.

That a person who has hitherto lived peaceably, quietly, contentedly, as one
who had enough in what was his own, should suddenly rush out and begin to make
a noise; that a man who is not otherwise known to have overstepped the bounds
of propriety should suddenly appear in motley dress, and in borrowed plumes,
and as one who is about to triumph; that an otherwise sober author, to whom the
authorities of literature have hitherto not denied a certain cultivation,
composure, and thoughtfulness, should all at once take it into his head to
comport himself as though he had either been stung by a fixed idea or else were
so far possessed by the demon of vanity that, for want of being something
himself, he had to fashion himself after another: that such a thing can happen,
how sad, how lamentable! that such a thing really does happen, how sad, how
lamentable for the one to whom it happens; but --- when it really does happen
--- it does not after all corrupt the cause, otherwise serious, that such a
prostituted person imagines himself to be championing.

If, on the other hand, an author with a somewhat sounder sense set about
treating a matter of faith, and yet nevertheless (be it from presumption or
from clumsiness) got the higher so entangled in the lower, the religious in the
worldly, the matter of principle in the personal, that he not only vexed the
established order by his loud-voiced objections (which in the old forms might
perhaps otherwise still preserve much that is both noble and good), but even
attempted to draw into this relative squabble, in itself utterly subordinate
for the religious, ``another'' as well --- ``another'' who is otherwise never
present where trumpets are blown, ``another'' who otherwise declines every
outward advancement, ``another'' who otherwise by no means covets to be
distinguished in wars, if only he
% --- p. VIII ---
may be unnoticed in peace, ``another'' who otherwise so urgently prays that he
as an irrelevance may always be forgotten, and that the thought, only the
thought, be remembered: then such an unpleasant proceeding ought assuredly to
be rejected from both sides with energetic protest, and every word in the
unwarranted attempt annulled, as the court annuls an abusive word.

This and much else I have now discussed and weighed as best I could; but with
all this I must still (at least provisionally) confess myself to the opinion
that one ought still --- \textit{propter iniquitatem temporum}! still --- as
matters stand \textit{in republica litterarum}! still ought to make an attempt
and dispute about faith.

Should it at length come to the point that one did not even find it worth the
trouble to dispute about faith, would that then also be a good sign? Would it
not rather be a sign that one despised this human struggle because at bottom
one did not respect the divine one either; a sign that one found it still less
worth the trouble to strive for faith: a sign that faith was no longer to be
found upon earth?

The author of the present writing candidly confesses that he does not yet
rightly understand, either in word or in deed, how to strive for faith as one
ought; but in the hope that the divine may one day fall to his lot, if only he
does not, in all too sensuous security, all too much neglect the human, he has
wished --- if not indeed for the instruction of many, then at least for his own
strengthening --- to express himself in these lectures concerning the manner in
which, at the present time, one --- disputes about faith.

\bigskip
\noindent Copenhagen, 17 February 1850.
\hfill R.\ Nielsen.

%% ============================================================
%%  FIRST LECTURE (pp. 1--13 = PDF 14--26)
%% ============================================================
\chapter*{First Lecture}
\addcontentsline{toc}{chapter}{First Lecture. Introductory Considerations}
\label{ch:forel1}
\markboth{First Lecture}{}

\begin{center}\itshape
Introductory considerations. Christianity is higher than science,
and the faith of the Gospel is heterogeneous with theology.
\end{center}

% --- p. 1 ---
It is the relation between the faith of the Gospel and theology that is made
the object of these lectures. Instead of saying ``the faith of the Gospel,'' we
may at once say ``Christianity''; on the other hand, we cannot without further
ado make ``theology'' equivalent to ``science.'' The question of the relation
of the faith of the Gospel to theology thus points provisionally to another and
more general question, namely that of Christianity's relation to science.

As is well known, this relation is conceived very differently at different
times and in the different schools, now as a relation of conflict, now again as
one of reconciliation.

In the conflict between Christianity and science there are some who go so far
as to demand, in the name of science, that Christianity be extirpated, while
others, on the contrary, in their zeal to advance Christianity, preach crusades
against free science.

With the feeling that there cannot possibly be conflict in the essence of
truth, and that the kingdom of truth cannot be divided against itself, men have
then again labored to reconcile the contending
% --- p. 2 ---
parties and to bring about a reconciliation by little by little making
Christianity scientific and --- what comes to the same thing --- science
Christian.

Which of these endeavors are we now to grant our approval? That is not so easy
to say. The whole matter seems so involved that we cannot approve any single
one of the attempts indicated here without approving them all, nor reject any
single one without rejecting them all. For if those who want the conflict are
right, then those who want the reconciliation must be right too. To invite to
reconciliation is here the same as to presuppose conflict and to incite to
conflict; and conversely, to proclaim conflict is here the same as to invite to
reconciliation.

This obscure speech will perhaps clear up and become intelligible when I now
set forth my thesis, the two-sided fundamental thought about which these
discourses are to move, the thesis, namely, that \emph{Christianity is higher
than science, and the faith of the Gospel heterogeneous with theology.}

The first member of this thesis is hypothetical; the last is categorical.

That Christianity is higher than science cannot be established by any
universally valid proof whatsoever. A universally valid proof of this would
have to be a scientific one; but a scientific proof that Christianity is higher
than science is a contradiction.

In hypothetical form, on the other hand, the proposition can be expressed thus:
``if Christianity is what it says itself to be, a revealed truth, then this
truth must belong to a sphere higher than the one to which science can raise
itself''; and conversely: ``if science's sphere is the highest sphere of the
highest truth, then Christianity is lost, and faith abolished.''

That science's sphere must necessarily be the highest sphere of the knowledge
of truth, science cannot prove;
% --- p. 3 ---
for in undertaking to conduct the proof it begins precisely by presupposing the
very thing it meant to prove, and does not get outside its own circle.

Since science, then, can neither prove nor disprove Christianity's truth, we
may with every reason combine our two hypothetical propositions in the
categorical assertion that Christianity in any case falls outside science, or
in other words, that Christianity is heterogeneous with science.

Should this assertion --- it is as yet, of course, only an assertion --- prove
defensible, then three things follow at once:

a) that the vehement conflict between Christianity and science must rest on a
misunderstanding; for between heterogeneous powers there can, properly
speaking, be no direct conflict. Science as science cannot possibly will to
offend Christianity, any more than Christianity can will to hamper or affront
science, if it is true that they move in wholly different spheres.

b) that the reconciliation aimed at between Christianity and science must
derive from a great misunderstanding. If the conflict is an illusion, the
reconciliation must be an illusion too; for a real reconciliation presupposes a
real conflict.

c) that the only condition under which science can come to offend and affront
Christianity is that it draw Christianity into its own sphere and make it
homogeneous with itself; and that theology is consequently, among all the
sciences, the only science that can in real earnest offend and affront
Christianity.

However offensive this last assertion may perhaps be to many, it follows with
unmistakable consistency from the thesis set up. Our investigation thus fastens
on the single question whether Christianity really
% --- p. 4 ---
is heterogeneous with every science, by whatever name it may be called, or not.

Christianity proclaims itself, like Judaism, to be a supernatural revelation.
God speaks, the miracle happens: with that revelation begins. This beginning is
of an unscientific or, if you will, of a supra-scientific nature, and can never
become an object of rational grounding.

What Jehovah says and does in the Old Testament is incommensurable with human
understanding, and cannot be scientifically grounded; what the Son, ``the
Only-begotten,'' says and does in the New Testament is incommensurable with
human understanding, and cannot be scientifically grounded. Moses and the
prophets, Christ and the apostles are unscientific, that is to say:
supra-scientific figures. None of them comes in the name of science, but they
all come forward, and speak and act, in the name of the revealed God.

The miracle is in its essence above, and therefore in its appearance against,
the law of nature, and cannot become an object of scientific treatment. Moses
strikes the rock, and the spring gushes forth: this fact natural science cannot
use, it does not concern it at all. Christ, ``the Son of Man,'' who goes to
meet his suffering and crucifixion, says in the prayer that he had glory with
the Father before the world was; this speech philosophy cannot fathom, though
it may well see and understand that, if the speech is true, then it must be a
supra-scientific, a divine speech.

If revelation is in its principle heterogeneous with science, how then should
they come into conflict with one another?

Revelation meets with natural science. Natural science now assures us that if
miracles have happened, then they must have happened against the order of
nature; but the very same thing is taught by revelation as well. Let natural
science then investigate things and discover the laws; let it convince itself
in continued inquiry that in
% --- p. 5 ---
nature's kingdom there is no leap, in the system of laws no place for the
supernatural facts; let it go further still, let it with ingenious sagacity
gather the manifold analogies of its rich experience in order to make these
facts as unreasonable, as improbable, as possible: it will nevertheless still
not come to offend Christianity; for Christianity by no means asserts that its
facts should accord with the analogies of sensuous experience. Natural science
may quite certainly assure us that Christianity's fact cannot be explained from
the connection of natural forces; provided only that it does not assure us of
more (and more it is surely not entitled, as science, to assure us of): then it
does not offend Christianity. It may even, if it finds itself called upon to do
so, take one step further and say that the phenomena of revelation are so
contrary to the laws of nature that they appear even absurd in relation to
natural phenomena; provided only that it does not say more (and more it is
surely not entitled, as science, to say): it still does not offend
Christianity; on the contrary! precisely in this respect there is here the very
best understanding between natural science and Christianity, since they both,
each in its own measure, at bottom say the very same thing.

Revelation next meets with philosophy, and the horizon widens. Natural science
and Christianity can still more easily preserve their heterogeneity; only as it
were in passing does natural science graze the precincts of the supernatural,
and it feels seldom tempted to press deeper in. It is otherwise with
philosophy, which, in testing all things, makes even Christianity itself an
object of scientific testing. But is this then permitted? Of course! So liberal
is Christianity that it does not even refuse to become an object of
philosophy's severest investigation (were it to refuse this, it would offend
science). So liberal is Christianity that philosophy
% --- p. 6 ---
is thus allowed to do with it what it will, to declare it comprehensible if it
can justify that before science, to reject it as incomprehensible if it can
justify that before science, but --- be it noted --- on one condition: that
philosophy with all this does not at the same time take it into its head to
want to make itself homogeneous with Christianity.

If this condition is inviolably kept, then there is surely no danger, either
for philosophy or for Christianity. If philosophy declares that it has now at
last comprehended, or is at any rate well on the way to comprehending,
Christianity, there is still no danger, provided it is said only in the name of
science. Because Christianity, \emph{scientifically understood}, now and then
shows itself to be comprehensible (or at any rate nearly comprehensible), it
can for that very reason, \emph{Christianly understood}, be just as fully what
it is and what it has always been: wondrous, inconceivable, for us sinful human
beings incomprehensible. And because philosophy now and then flares up and
wants on science's behalf to remonstrate with the cultivated world that
Christianity has now at last shown itself in its true light, as
incomprehensible, absurd, and meaningless, there is here nevertheless no real
danger, for here at bottom there is no real conflict. No, far from it!
Philosophy is now, without perhaps in the heat of it rightly noticing so
itself, just now on the way to the very best understanding with Christianity.
Christianity has, after all, never promised that it would ever let itself be
comprehended (or at any rate nearly comprehended) by philosophy. Christianity
is so far from being anxious about this philosophical discovery that on the
contrary it congratulates science on the discovery, since it is so glad to see
that philosophy, which seeks the truth, should at last come to the knowledge of
truth, i.e. to the knowledge that it is heterogeneous with Christianity.

So it is strange enough with this relation between science and Christianity.
They quarrel, are reconciled,
% --- p. 7 ---
and quarrel again. According to the prevailing way of seeing things, one is
much troubled about the conflict and ponders deeply on the reconciliation; and
the conflict would be still harder, and the reconciliation awaited still more
ardently, if the world had not already grown so clever that, with the feeling
that at bottom not much comes of either the one or the other, it had already
begun to grow a little embarrassed by the whole thing.

Here the matter stands otherwise. Instead of being troubled about the conflict
and hoping for the reconciliation, as if it were self-evident that both the
conflict and the reconciliation had true validity and exceeding significance,
we are now on the contrary so placed that we cannot rightly see what all this
fuss actually leads to, since the whole seems to rest on a misunderstanding.

There must, however, be a ground for this misunderstanding; there must be a
medium in which the heterogeneous powers meet one another, and engage one
another in such a way that they at least acquire the appearance of being
homogeneous: where now is this ground? which is this medium? --- the existing
human being!

In human existence, i.e. in the existing human being, heterogeneous magnitudes
are bound into unity. The sensuous and the spiritual, the temporal and the
eternal, the natural and the supernatural meet here. It is the human being's
general task to relate himself differently to what is different. But when
revelation enters in, the difference and the opposition become so strong that
the human being, even though he feels within himself that truth is only one and
can only be one, is nevertheless led into so deep a dualism of truth that in
this existence he must relate himself in entirely different ways to two truths
so heterogeneous in existence as: \emph{the truth of revelation and ordinary
rational truth.}

From this it is explained how the conflict between Christianity and science
arises within Christendom,

% The Danish reproduces the print's comma here (verified at 600 dpi) together
% with the paragraph break that printed p. 8 opens with --- the printing has
% both at this joint. Both are reproduced in the English too: the comma stands
% and the paragraph breaks. It reads as a colon-like pause introducing the
% elaboration that follows.
% --- p. 8 ---
The Christian becomes a naturalist; the naturalist is a Christian. Both
heterogeneous determinations are combined in one consciousness. The task is
that the individual concerned shall now understand how to relate himself
Christianly to Christianity and scientifically to science, so that he does not
derive conclusions for natural science from Christianity, or conversely draw
consequences for Christianity from natural science. But since it is after all
the same human being, the same consciousness, that is here to unite the
heterogeneous in existential unity: how near at hand does not the temptation
then lie to make the heterogeneous magnitudes a little homogeneous, to treat
science Christianly and Christianity scientifically, to awaken the discord by
striving after the reconciliation.

The Christian becomes a philosopher; the philosopher is a Christian. What
matters now is to relate oneself Christianly to Christianity and scientifically
to science. But since Christianity proclaims itself as the truth, and since
science, \textit{in specie} philosophy, investigates and tests the truth: how
near at hand does it not then lie to make Christianity's religious knowledge
homogeneous with the philosophical knowledge of Christianity, to reconcile
faith with knowledge, and precisely thereby to awaken discord between faith and
knowledge!

Natural science and philosophy are, however, with respect to principle and
point of departure still so different from Christianity that one easily becomes
aware of the heterogeneity, provided the conflict is only carried on
consistently. But there is yet another science, a science far more intimately
akin to Christianity, a science that with a certain self-regard calls itself
the Christian one, a science that claims to have principle and presuppositions
in common with Christianity, a science that will absolutely defend Christianity
in order to defend itself along with Christianity, a science that presses so
close up against Christianity that it is as though they were to live and die
together;
% --- p. 9 ---
this science is theology: is Christian theology, then, also heterogeneous with
Christianity? --- Certainly! Theology as science is assuredly just as
heterogeneous with Christianity as any other science. That the relation is here
more ambiguous is explained quite simply by the fact that theology is itself an
ambiguity.

How matters stand more closely with this ambiguity is not difficult to
discover.

Faith cannot be without knowledge. The believer must know what it is that he
believes, must be able to give an account, to himself and to others, of the
content of his faith and of the communion to which he confesses. Insofar as
theology now sets itself exclusively the task of finding out what the Christian
is, of developing and explaining what it is that the Church believes and the
believer confesses, to that extent theology as theology is doubtless
homogeneous with the objective doctrine of Christianity; but this objective
doctrine of Christianity is therefore by no means homogeneous with Christianity
itself. Insofar as the theologian relates himself Christianly to Christianity,
he does not relate himself scientifically; insofar as he relates himself
scientifically to Christianity, he does not relate himself Christianly.

The old orthodox theology, which in all points would be conformable to the
Church's doctrine, suffered from a great ambiguity. As the comprehensive,
complete, and exact presentation of the Church's doctrine, this theology was at
bottom nothing other than an articulated confession of the content of faith,
and as such not yet science. It was at most only the scientific expression and
form of faith-knowledge. But scarcely had it begun in form to take on the
stamp of science, and thus taken the first step toward a reconciliation of
Christianity and science, before it also found itself placed in a hard
self-contradiction. The content was Christian, the form scientific; this was
the contradiction between
% --- p. 10 ---
form and content, a contradiction under which orthodox theology at length
succumbed.

Orthodox theology presents every doctrinal proposition thetically, insofar as
it presents it in the name of faith; this is the Christian element. But then it
also treats every doctrinal proposition hypothetically, insofar as it treats it
in the name of learning, investigation, and research; this, then, is the
scientific element. Accordingly, every doctrinal proposition comes to be
regarded from a double point of view. Christianly understood it belongs to the
Church, is exalted above science, and defies every scientific attack.
Scientifically understood it stems from theological research, theological
discussion, and theological testing, and is thus forfeit to science, and must
as such become the object of a new discussion, a new research, a new testing.
Under the continued discussion, research now stands against research, testing
against testing. If the research is to signify anything for Christianity, it
must conform to the Church; if it is to count for anything in science, it must
be free and independent of the Church. Theological testing must therefore at
one and the same time be both free and unfree, both outwardly dependent and yet
outwardly independent; that is a contradiction. The theological doctrinal
proposition shows itself here as an ambiguous coin: on the one side, the sign
of the Cross! on the other, Minerva's owl! Here it does not even help to say:
render unto Caesar what is Caesar's; for this ambiguous coin belongs with equal
right both to God and to Caesar.

Theology, which in the beginning would be one with the Church's doctrine, then
gradually began, under free research, free discussion, and free testing, to
free itself from its original relation of dependence and to transform itself
from orthodox into heterodox: the more heterodox it became, the more
scientific, and conversely: the more scientific, the more heterodox. Only
% --- p. 11 ---
in the last century has theology grasped itself as an independent science; only
in the nineteenth century is it clearly seen that theology is heterogeneous
with Christianity.

But Christianity is homogeneous with the faith of the Gospel.

Christianity is only for the believer; theology is only for the knower. The
faith of the Gospel lives in the Church; theology lives in science. If they
will go on relating themselves as homogeneous, they must work against each
other, and will then soon work out their own mutual dissolution.

The faith of the Gospel lives in the Church. What, then, is the Church, and
where is it? We answer with the Augsburg Confession: ``\emph{The Church is the
communion of saints, in which the Gospel is rightly preached and the sacraments
rightly administered}'' (\textit{Ecclesia est congregatio sanctorum, in qua
evangelium recte docetur et recte administrantur sacramenta}).

But theology, what is it? where is it? Theology is the science that undertakes
to subject the Church's Gospel to a historical-critical testing; the science
that as time and occasion allow permits itself to reason for and against the
Lord's sacraments; the ``ecclesial'' science that may not indeed venture to
make itself heard in the church, but may all the more freely make itself heard
in the school.

The heterogeneity is now conspicuous.

The Church is that communion of faith in which the Gospel is \emph{rightly
preached}. But the Gospel is rightly preached only when it is preached as God's
Word. God's Word is an infallible, definite, and authoritative Word. Therefore
the free preaching of the Word in the congregation always goes back to the
revealed original Word; the evangelical sermon must have its text: Holy
Scripture supplies the text. The biblical text is read aloud for the
congregation's devotion; devotion is unscientific; devotion knows no criticism:
the text read out is in every way infallible and trustworthy, as it ought to
be. ``So many are the words,'' says the pastor, and then he adds a prayer. Thus
the evangelical
% --- p. 12 ---
sermon does not begin with ratiocinating criticism, but with prayer and
invocation, and only thus can the Gospel be rightly preached.

It is otherwise in theology; for theology is, after all, a science. Science, as
science, cannot begin with prayer and invocation. If, then, the Church's Gospel
is to be tested in science, it must be treated in a quite different way. And so
it is. In science ``the critical introduction'' sits with the Gospel before it
and interrogates it as to whether it is genuine or perhaps spurious; and the
longer one asks, and the more one interrogates, the greater the suspicion one
acquires that this Gospel is perhaps after all --- spurious. Then
``hermeneutics'' comes forward and will in the same Gospel scientifically
separate the spirit from the letter, the idea from the outer garb, the truth
from the representation; and the longer one sifts, the more acutely one
separates, the less there comes to be of God's Word and the more of the human
word in the biblical text.

If any conclusion is to be drawn from this, it must surely be the conclusion
that what the Church puts into faith, theology converts into doubt.

The Church is that communion of the Spirit in which the \emph{sacraments are
rightly administered}. But the sacraments are rightly administered only under
the one decisive presupposition that the word of institution is an infallible,
divine, miracle-working Word. The Church confesses that Baptism is a bath of
regeneration by the Holy Spirit, and the Lord's Supper the communion of
Christ's body and blood for the forgiveness of sins to those who believe. The
Church builds on the word of institution, but does not at all apply itself to
comprehending the sacrament's mystery. As the Word sounds over Baptism, so it
shall be believed --- against understanding; as the Word sounds in the Supper,
so it shall be believed --- against understanding. The miracle, the wondrous
efficacy of the means of grace, is not to be explained at all, not comprehended:
the appropriation is only in faith.

% --- p. 13 ---
It is otherwise in theology; for theology is, after all, science. As objective
science, theology must of course make the sacraments themselves into critical
objects, and then treats the means of grace as means for heightening the
scientific luster of idea-knowledge. Then it becomes doubtful again, highly
doubtful, how one is to explain to oneself the union of the supernatural with
the natural. The difficulty is that one cannot explain it rightly, and yet
would so heartily like to have it explained. The longer one goes on explaining,
the more doubtful the explanation becomes.

If any conclusion is to be drawn from this, it must surely be the conclusion
that what the Church puts into faith, theology converts into doubt.

That theology may in its own way be right in engaging in scientific
investigations, and therewith also right in engaging with doubt, shall not be
denied; but what must be denied in the most definite terms is that theology
should have the right to confound its own task with the Church's. It is no
accident that the Church lets everything be decided in faith; neither is it any
accident that theology dissolves everything in doubt. The Church tears itself
down when it doubts its own authority and gives up faith; theology denies
itself when by an authoritative pronouncement it breaks off its investigations
in order to arrest doubt.

As doubt cannot be nourished in the Church, so faith cannot be nourished in
theology. The Gospel in the Church confirms faith, as faith the Gospel; but
doubt in theology conflicts with the Gospel, as the Gospel with doubt.

If any conclusion is to be drawn from this, it must surely be the conclusion
that doubt is the Church's enemy, and \emph{the faith of the Gospel
heterogeneous with theology}.

%% ============================================================
%%  SECOND LECTURE (pp. 14--25 = PDF 27--38)
%% ============================================================
\chapter*{Second Lecture}
\addcontentsline{toc}{chapter}{Second Lecture}
\label{ch:forel2}
\markboth{Second Lecture}{}

\begin{center}\itshape
In its conception of the relation between historical and religious faith,
theology must separate what in religion is inseparable, and therefore the
faith of the Gospel is heterogeneous with theology.
\end{center}

% --- p. 14 ---
Positive revelation --- in Judaism and Christianity --- is a historical
revelation. As revelation it speaks to the inward human being and presupposes
subjective religious faith (\textit{fides religiosa}); as historical revelation
it works upon the natural human being's representation through a complex of
supernatural but nevertheless actual facts, and thus presupposes at the same
time a historical faith (\textit{fides historica}).

He who now relates himself Christianly to Christianity can and may not separate
the historical from the religious. In true and living faith both factors must
work together without, however, melting together. Since now historical faith,
considered by itself, is heterogeneous with subjective religious faith, while
the believer nevertheless unites these two heterogeneous determinations in the
unity of the faith of the Gospel: it is in consequence exceedingly difficult
for the ratiocinating understanding to understand the religious man.

% DEFECT (Danish p. 15): the Religious man's second speech is printed with ONE
% opening „ but TWO closing “ (after „tredie Dag“ and again after „derpaa“).
% The mark that is missing in the print is the OUTER opener: the single „ opens
% the creed recitation, the first “ closes it, and the second “ closes an outer
% speech never opened. The English supplies the outer opener and sets the creed
% as an inner quotation, so it balances. Danish stands as printed.
``The sacred facts do not become actual because I believe in them; but
conversely, because I know that they are actual, are historical, therefore I
believe in them. It has happened that
% --- p. 15 ---
God revealed himself to Abraham; it has happened that God in former times spoke
many times and in many ways to the fathers by the prophets; it has all
happened, as Scripture witnesses, and therefore I believe in it.'' Thus speaks
the religious man. --- ``\,`It has happened that Jesus Christ, God's
only-begotten Son, born of the Father from eternity, was in the fullness of
time conceived by the Holy Spirit and born of the Virgin Mary; it has happened
that he suffered under Pontius Pilate, that he was crucified and died and was
buried and descended into hell and rose again on the third day'; it has all
happened, as Scripture witnesses and the Church confesses; and since it has
assuredly happened, I believe in it.'' Thus speaks the religious man.

But for the ratiocinating understanding it is nevertheless very difficult to
understand this speech. It constantly looks as though the religious man were
straightforwardly building his faith on a straightforward historical certainty.
It looks as though revelation's facts had, in his opinion, so many outward,
particular, finite testimonies that no one could reject these facts without
rejecting all history. The religious man must then appear to the cultivated
(especially the theologically and scientifically cultivated) as an unscientific,
uncultivated, and simple person who has so much faith because he knows so little
of doubt. Many a one would then perhaps say: ``If only I knew that it had all
happened as it stands in Scripture; if only I knew that this one thing had
happened, that Christ, after he had really died and been buried, rose again on
the third day; if only I now knew that this one thing had really happened: yes,
then I should surely believe, and in faith become a new and better human being.
But who will vouch for us that these narratives are genuine, are reliable and in
every way trustworthy? The believer is so sure of his case because he does not
know the difficulties at all: he should just begin with the investigations, he
should know the objections, he should just try
% --- p. 16 ---
attending a theological course on the critical introduction, and he would then
soon enough come round to other thoughts.'' --- Thus judges the ratiocinating
understanding; and yet it is a misunderstanding.

The religious man can, if it must be so, very well acquaint himself with the
theological difficulties without taking harm to his faith. With ``the critical
introduction'' there is no danger at all. If the three men could go unharmed
through the fiery furnace, should the believer then not also be able to go
unharmed through ``the critical introduction''! If we then meet the religious
man as he steps out of criticism's labyrinth, and ask him how his conviction now
stands, he answers us that he has indeed now learned much that he did not know
before, but that he has thereby only convinced himself the further of what he
already knew beforehand: that human beings' doubtful thoughts must bow before
God's revealed thoughts, that human beings' works avail nothing against God's
works, and all human writings nothing against God's own Holy Scripture; and in
speaking thus, he speaks precisely out of his subjective religious faith
(\textit{fides religiosa}).

But this speech too is, when we consult the ratiocinating understanding, exposed
to misunderstanding. The believer who defies all in reliance on the revealed
Word now shows himself as one who is so sure of the historical because he at
bottom derives his historical faith from his religious faith. Many a one then
perhaps thinks in his own quiet mind: ``The believer is right. Nor does the long
learned introduction signify much, if only one could believe from the heart. If
only I could believe that Christ is risen, then he is assuredly also risen,
risen for me. But what use is it to hold to the Gospel of the resurrection when
this fact nevertheless stands as a distant and alien event that makes no
impression. If the event could become present to me, if I could get an immediate
impression of
% --- p. 17 ---
the risen Savior, like Mary Magdalene when she met him in the garden, or like
Thomas when he thrust his finger into the print of the nails and cried out: my
Lord and my God! But that is impossible; with the outward fact one cannot begin.
The historical is something past, and can no longer become present: one must
begin with subjective religious faith.''

Thus judges the reflecting observer, as he as it were inclines toward faith
without, however, being able to grasp faith. And yet this too is a
misunderstanding. One can in truth as little derive historical faith from
religious faith as conversely the latter from the former. When we one-sidedly
begin with subjective faith, the religious man holds back, throws himself to the
opposite side, and lays an infinite emphasis on the historical. ``It has all
happened as it stands written, and since it has happened, I believe in the
fact.'' --- If we then take this historical element in and for itself, and let
theological criticism shake the fact: then the religious man again strains
against it with all his might, throws himself to the other side, appeals to the
Spirit's inner testimony, and now holds the historical fast with infinite
passion in subjective, religious faith.

To want to derive religious faith from historical faith is a misunderstanding
that leads to a Catholic extreme; to abandon the historical and hold to
subjective religious faith is a modern Protestant extreme.

The Catholic Church would by a divine authoritative pronouncement forbid all
freer criticism, and comes into conflict with science. By its authoritative
pronouncement it shows precisely that it fears science, and is afraid that all
is not quite right with the historical. This fear awakens suspicion; this
suspicion undermines faith.

Modern Protestantism sets criticism free; criticism dissolves the historical,
and undermines faith. For if ``the critical
% --- p. 18 ---
introduction'' is to have the least influence on religious conviction; if from
criticism's results one may draw a single conclusion for or against faith, then
it is assuredly impossible for subjective faith, at least at that moment, to
believe religiously in Holy Scripture.

Whence now does this come? Doubtless from the fact that in Catholicism as in
Protestantism one has made \emph{Christianity homogeneous with science, the
faith of the Gospel homogeneous with theology}.

If Catholicism were conscious that Christianity and theology are heterogeneous,
it would not need at all either to forbid or to restrict free criticism. If
Protestantism fully saw that the faith of the Gospel is heterogeneous with
science, the Church would not need at all to heed the critical dissolution.

The delusion maintains itself by the fact that the historical and the
subjectively religious work together in the unity of faith, notwithstanding that
the historical, as historical, is nevertheless heterogeneous with the
subjectively religious.

Religious faith cannot draw a single breath without supporting itself on its
historical presupposition. This Catholicism has seen quite rightly, and
therefore it inculcates the historical. But now the Church's doctrine is at the
same time to be theology, and theology is to secure the broad historical
foundation (of Scripture and tradition) by historical testimonies and historical
learning, is by scientific research to raise objective, historical certainty to
absolute, incontrovertibly decisive certainty. That theology cannot do, it is
far beyond its powers; so one takes refuge in the ecclesial authoritative
pronouncement, and the Church tyrannizes over science.

In place of the ecclesial authoritative pronouncement, Protestantism sets free
subjective faith (\textit{fides religiosa}). The Protestant consciousness has
seen quite rightly that the historical is heterogeneous with the religious, that
everything at last depends on the Spirit, the
% --- p. 19 ---
living Spirit indwelling the congregation, and not on a greater or lesser sum of
historical, long since past events. Against this general principle we make no
general objection either; only it depends on how it is applied in practice.

The false application is that in the name of faith and on faith's behalf one
separates what faith cannot and may not separate. --- One takes the historical
by itself and makes it the object of a critical-scientific treatment. ``The
Bible is a collection of various writings of antiquity, written at various times
by various authors under various conditions.'' Thus ``the critical
introduction'' begins. Against this there is nothing to object, if ``the
critical introduction,'' in beginning its historical-critical course on these
``writings of antiquity,'' will then at once declare that the scientific
investigation signifies nothing, nothing whatever, for the faith of the Gospel
and the Church's doctrine. But that ``the critical introduction'' will assuredly
not declare. It rather lets it be understood that these investigations are
important, very important, not only --- what we readily grant --- for the sake
of learned science itself, but also important, indeed in all their particulars
and their whole prolixity highly important, for all the Church's future
teachers, for the pastors to be. ``The critical introduction'' will therefore
(mediately or immediately) have a voice in deciding what the religious
conception of Holy Scripture ought to be.

``The critical introduction'' thinks after this fashion: let us, with the
reservation of a certain general faith, first get the purely historical
critically settled, and then we can always afterward draw conclusions from it
for our religious conviction. But this is an entirely false presupposition. The
law is that historical faith point for point is to be maintained and borne by
inward, strong, subjective, religious faith; the law is that historical faith
perishes on the instant when it is abandoned by the religious. This law
% --- p. 20 ---
is violated by the critical introduction. The proof of this is that criticism's
objective, or rather half-objective and half-subjective, settlement runs out
into endless doubt.

Historical faith is in its essence relative, and goes only upon the probable;
but the merely probable can continually be converted and weakened into the
improbable. Subjective religious faith is the decisive, the unconditional, the
overreaching, which grasps the absolute. But in religious speech it nevertheless
looks as though both were equally absolute; for the religious man does not
separate them: only in reflection do they become separated.

The relation can therefore be expressed thus:

\emph{Subjective religious faith rests upon the historical as upon its
foundation, but itself bears the foundation upon which it rests.}

Having thus determined the relation in general, let us then go more closely into
the particular, and consider, for example, \emph{the theologians' dispute over
the Apostles' Creed.}

If only one would be content to relate oneself Christianly to Christianity and
scientifically to science, then the Church could at once, for its part, end the
dispute by vouching for the historical validity of its confession of faith; but
now one is (if for nothing else, then for theology's sake) also to relate
oneself scientifically to the Apostles' Creed: now the critical introduction
begins again. --- ``Can one then really believe so unconditionally in this
symbol as the Church assumes? Does it then come word for word from Christ
himself or from his apostles? Has it not rather formed itself gradually? Is not
one article taken up against the Gnostics, another against the Apollinarians? Is
``the descent into hell'' genuine? Let us compare the various copies and
recensions. Is not the recension found in Rufinus, after all, remarkably
divergent from other corresponding recensions?''

% --- p. 21 ---
Now a mass of difficulties arises, a mass of misgivings; now there is strife
here as for life and death (for here the strife is about the Church's very
foundation); now the theologians divide into two ranks; now it goes: with
grounds --- for and against, witnesses and testimonies --- for and against,
interpretations of these witnesses and these testimonies --- for and against,
interpretations again of these interpretations --- for and against. A scientific
vapor of opinions and pieces of reasoning envelops the contenders! It is an
involved strife; for theology carries on not only the strife, but at the same
time the history of the strife; whereupon strife can again arise about how the
history of the strife ought properly to be presented and judged: it goes on
into infinity!

What, then, is the Church now to do with all this strife? Is it to forbid it by
an authoritative pronouncement? That would be Catholic, and is besides in
Protestantism entirely impossible. Is the Church then perhaps to expect some
advantage from the strife's final settlement? No, by no means! If the Church at
any moment supports itself on a scientific defense, it must also put up with
suffering under a scientific attack. If the Church will share theology's
victories, it must also put up with sharing theology's defeats; but this
conflicts both with the Church and with Christianity. Rather, the Church shall
quite calmly rely on its own authority, let the strife go its scientific course,
and only see to it that science keeps to what is its own.

Scientific criticism is in no way able to deprive the Church of its confession
of faith. The objections, when pressed to the very furthest, cannot after all
carry it further than to establishing an improbability with respect to the
symbol's genuineness. When criticism has said: this confession of faith's
origin from Christ and the apostles is to me in the highest degree improbable:
then doubting, contradicting criticism has said the last word, and can say no
% --- p. 22 ---
more. But about this alleged improbability two things are to be noted:

a) The highest degree of improbability, which is otherwise decisive for every
other historical-critical testing (when \textit{fides historica} is taken in and
for itself), will in this case not suffice. In the communication of the
confession of faith the supernatural is, after all, at work. The confession of
faith has its origin from inspiration, and points to the miracle. But the
miracle itself belongs to the improbable. To reject the symbol on the ground of
an improbability is therefore a contradiction: the miracle will continually
weaken and disturb the critical attack.

b) The criticism of attack is besides weakened in another way. Those who defend
the symbol find just as many grounds and testimonies for their assumption as
those who attack it. So the strife is in an unceasing fluctuation of probability
and improbability; here there is no settlement, the settlement is religious:
faith or unbelief. The Church, which speaks the language of faith, therefore
does right in maintaining the confession without heeding scientific criticism.

If we now place the believer in the midst between criticism's doubt and the
Church's confession, how will he then relate himself? Differently, of course, to
what is different: Christianly to Christianity and scientifically to science.

The believer may well take part in the theologians' critical dispute over the
symbol, and reason for and against, like all the other critics (for insofar as
he relates himself scientifically to science, he is of course also a critic).
But wherein does he as a believer then differ from the theologian? The
theologian expects, after all, that at last in time some great
Christian-ecclesial result is to be brought out of the long critical
investigation; but that the believer does not expect at all. In relation to the
Church's earnestness, the whole prolix strife is to him only as a humanly
entertaining play of thought. The Christian
% --- p. 23 ---
settlement is in faith, not in theology. This will also emerge from a
consideration of the strife's inner character.

Those who attack the symbol bring forward, say, 3 grounds, 3 proofs against it
(whereby improbability provisionally gains the preponderance). Then the
defenders come forward with 4 grounds, 4 proofs for it (and at the same moment
probability rises). After some time has elapsed the attackers find out 7 grounds
and 7 proofs (improbability has the preponderance again), and so on.

The religious man, who sees and understands what comes of such proceedings, is
after all perfectly right when he says to himself: if I am to wait until this
dispute comes to an end, then I shall have to wait until doomsday; but if I am
to break off at the 20th ground because a counter-ground follows, or at the 21st
proof because a counter-proof follows, then I may just as well break off at the
first and only decisive ground, namely the one that lies in the Church's
testimony.

Who, then, is here the more intelligent, the believer or the theologian? I
should think the believer, and I venture besides to establish it by a
mathematical example. There are, as is well known, infinite fractions; to want
by continued division to exhaust such an infinite quotient in finite numbers is
proof of mathematical ignorance. But so naive is ``the critical introduction''
in all its acuteness that it supposes it can exhaust an infinite quotient in
finite numbers.

Yet there is still a difficulty remaining, not for faith, not for the Church,
but for ratiocinating reflection, for theology.

% DEFECT (Danish pp. 23--24): the quotation opened at „lad os kun indrømme is
% never closed in the 1850 printing --- it simply stops at „udeladt?“. The
% English closes it after ``left out?'', where the sense ends. Danish stands as
% printed.
``Let us grant,'' one says, ``let us just grant that the Apostles' Creed is
according to its content --- that is to say, according to its general content
--- both apostolic and fundamentally Christian; but then no one can deny either
that this divine original Word is expressed in human words, articulated in
finite articles and
% --- p. 24 ---
finite propositions. These articles and propositions can surely not, word for
word, be entirely unimprovable or unalterable. Every time the symbol is
translated from one language into another, the individual propositions and
expressions do after all always undergo some alteration. The symbol is besides
far from containing all the important chief propositions of faith, and an
article more or less can therefore not possibly be decisive for one's eternal
salvation. Or should a believing Christian not be able to be saved because he
was baptized on a confession of faith from which the descent into hell happened
to be left out?''

Thus reasons theological criticism. But however much prestige this reasoning may
procure for itself in science, it is nevertheless, in relation to the Church, an
entirely subtle and sophistical piece of reasoning.

When the criticizing theologian asks whether a single article or a single word
in the confession makes any difference one way or the other, we must
involuntarily think of the sophist who asks whether a single hair torn out makes
a man bald, or a single grain a heap. Instead of engaging in sophisms, the
Church ought for its part immediately to stop the endless reasoning, and lead
the principle-less investigation back to its principle. The question of
principle is not this, how many letters or how many propositions or how many
articles can be thought of as left out of the confession of faith without harm
to salvation; the question of principle is whether a critical separation of the
divine and the human in the Church's original Word can take place without danger
to faith. This question has already been answered with a No. In the original
Word of the confession, God's Word is not to be separated from the historical.
--- Thus runs the Church's language: ``As no one can eat of bread in general,
but of \emph{this} bread, \emph{this} consecrated bread, which we break; and as
no one can drink of a cup in general, but of \emph{this} cup, \emph{this} cup of
blessing, which we bless: so
% --- p. 25 ---
too can no one be baptized upon a Christian confession in general, or in a
religious communion in general, or by a minister of the Word in general; but
this is the original Word's decisive authority, that he who is received into the
bosom of the congregation must be baptized upon \emph{this} confession, in
\emph{this} communion, in \emph{this} church, by \emph{this} pastor, on
\emph{this} day. This is the original Word's decisive authority, that the eternal
here unites itself with the temporal, the divine with the human, in indissoluble
definiteness; here too it holds good: `what God hath joined together, let no man
put asunder.'\,'' Against this language of the Church, theological reasoning is
vain and powerless! What critic, pray, has authority to take away a single
article of the Church's confession? and what theologian, pray, has authority to
insert a single article into this confession? --- The faith of the Gospel joins
what theology separates, and why? \emph{Because the faith of the Gospel is
heterogeneous with theology.}
% DEFECT (Danish pp. 24--25): the outer quotation „Ligesom Ingen kan æde …
% (opened p. 24) is closed by the SAME “ that closes the inner „hvad Gud haver
% tilsammenføiet …“ --- the printer set only one mark, giving the Danish a net
% +1. The English closes both: inner with `…' and outer with ''.

%% ============================================================
%%  THIRD LECTURE (pp. 26--37)
%% ============================================================
\chapter*{Third Lecture}
\addcontentsline{toc}{chapter}{Third Lecture}
\label{ch:forel3}
\markboth{Third Lecture}{}

\begin{center}\itshape
In its conception of the confession of faith and of Holy Scripture,
scientific theology must consistently go on working against faith;
for the faith of the Gospel is heterogeneous with theology.
\end{center}

% --- p. 26 ---
To want at one and the same time, in one and the same instant, to relate oneself
believingly and yet scientifically to one's confession of faith is a
contradiction. Insofar as I am in the confession, I am outside science; insofar
as I am in science, I am outside the confession. To be great in the confession
is something quite other than to be great and distinguished in science. The
fisherman from Bethsaida was nothing at all in science; but the fisherman from
Bethsaida nevertheless became great in the confession when he bore witness to
the Master and said: \emph{Thou art Christ, the Son of the living God.} And
Christ, who hears his confession and knows what a confession means, does not
answer him either in the language of science, does not place himself in a
critical-scientific relation to him, but strengthens his faith, rejoices over
his confession and calls him blessed: ``\emph{Blessed art thou, Simon
Bar-Jona; for flesh and blood hath not revealed it unto thee, but my Father,
which is in heaven.}''
% DEFECT (Danish p. 26): the closing “ after „Himlene.“ has no opening „
% anywhere --- the printer set only the closing mark, giving the Danish a net
% -1. The English supplies the opener at Christ's words, where the sense
% begins, and balances. Danish stands as printed.

But --- it is many years since the first confession of Christ was heard in the
world; in these many, many years one has then gradually fallen into the habit of
relating oneself scientifically to Christianity and theologically-critically to
the confession.

% --- p. 27 ---
When, then, the critical introduction now takes the confession before it, holds
it up to science's light, and inspects it from every side, it acquires (besides
the misgivings already adduced) yet another scruple, which we must here heed all
the more since this scientific scruple at the same time seems to be a religious
scruple.

One reasons roughly thus:

``If the believer may on no account separate the confession's divine content
from the human word and letter, then he may not have any inkling either that
this human word can be subject to any alteration whatsoever. Let one judge the
critical investigation as one will: he who has once been touched by doubt cannot
arbitrarily break off and return to the first simplicity. If the religious man
once learns that there is more than one draft, more than one recension of the
Apostles' Symbol, and if he is nevertheless to be bound to the confession's word
and letter: then he must live in constant anxiety that he has not yet found the
right original text; who can then blame him for seeking by continued inquiry to
arrive at the greatest possible certainty about the original, for taking an
interest, say, in knowing whether ``the descent into hell'' is originally
genuine or a later addition. Suppose that this article were spurious, suppose
that Christ did not descend into hell, suppose that the representation of this
descent into the realm of the dead were only to be understood figuratively of
Christ's actual death; then the believer who clings to the confession's word
comes to believe an untruth, comes in his blind zeal to hold fast to the
apostolic precisely to accept an unapostolic representation.'' ---

This solicitude and concern of criticism for the purity of the Church's doctrine
is, however fair and pleasing it may look at first glance, nevertheless of a
very subtle and sophistical nature. From fear that something might possibly have
crept in here
% --- p. 28 ---
in the historical (in \textit{fides historica}) that was not purely and
genuinely historical; from fear that this something should then in turn corrupt
the Church and harm faith, theological criticism will benevolently seduce both
the Church and the believer into making a separation that annihilates faith,
namely a theoretical separation of the divine and the human in the confession's
original Word. The theologian will therefore help the believer! Yes, indeed he
will help him; but the help is somewhat ambiguous, the help is such as if he
would help him out of the water into the fire.

The danger that a spurious article should, despite the Church's vigilance, have
entered into the confession of faith is, religiously regarded, nothing beside
the danger of entering upon a ratiocinating separation of the inseparable, and
thereby losing faith.

In reliance on the word of the Church and of the Gospel, the believer goes with
his confession confidently through the world, into eternity, before that Judge
who knows all things and judges righteously. And when the human being's soul
then stands before that Judge who knows all, both the historical and the
non-historical and the supra-historical: then it will surely appear where the
danger is and wherein it consists. The mortal danger is then not that here, in
the clarity of the light, it should suddenly be discovered that the confession
had by a human error acquired one article too many. No, the danger is quite
another, the mortal danger is that the human being has not lived in faith in his
confession as he ought. If only he has lived in faith, striven for the
confession and striven the good strife: then the Judge will surely not reproach
him that there is one article too many in the confession with which he comes,
but will take him into grace and call him blessed, as he called his beloved
disciple blessed when he first heard his confession. But woe to the soul that
comes in doubt with a doubtful confession. ``The critical introduction'' cannot,
after all, help him when the Judge demands a reckoning. Or what use would it be
if this
% --- p. 29 ---
doubting soul now began to excuse itself and say: ``Lord, there were so many
editions and so many recensions of our ecclesial confession of faith down there,
so many mutually divergent and differing recensions, that at last I did not know
which I should choose; criticism then advised me to separate the divine from the
human, in Thy confession as in Thy Gospel. And so I followed criticism's advice,
and so I entered into the scientific investigations, and by my scientific
investigations I came further and further into doubt, and so my faith languished,
for at last I no longer knew --- alas! --- what I really ought to believe.'' Will
not then the Judge of conscience answer this doubter and say: ``Thou fool, why
didst thou not believe the Church's testimony, as it runs word for word? why
wouldst thou not understand that my Word is an inseparable unity of divine and
human word, as I myself am an inseparable unity of God and man?'' ---

If anyone will absolutely misunderstand this statement, let him simply give the
matter the turn as though I were hereby aiming to make science hated and to
outlaw criticism; for that is truly a great misunderstanding. So far is it from
being the case that in this attempt I have intended to abolish criticism, that
on the contrary I have wished to help theological science to a most needful and
entirely indispensable criticism: the criticism whereby one separates the real
danger from the imagined.

When Luther strove for his doctrine of the Lord's Supper, ratiocinating
criticism was already in full swing, and would with all its might make a divorce
between the divine and the human in the word of institution. But Luther strove
for his text, for he strove in faith. ``These words of institution'' --- he says
--- ``bring it about that the bread is his (Christ's) body, and the wine his
blood. He who now eats this bread, he eats Christ's true body; whoever drinks
this cup drinks Christ's true blood, be he
% --- p. 30 ---
worthy or unworthy. That one shall firmly believe; for the dear Christians shall
do God the honor of confessing that what God says, that he can also do, as St.
Paul writes of Abraham that he did (Rom. 4:21). Whoever will be a Christian, he
shall not do as our fanatics and mutineers, who trouble themselves about how it
could be that the bread is Christ's body and the wine Christ's blood. They will
measure and comprehend God with their reason, and because it does not accord
with their reason, they suppose God cannot do it either. But what is this, that
one has already troubled oneself about it so long? Though one tore oneself apart,
one would still not be able to comprehend God with human reason; for our Lord God
is not such a God as lets himself be measured, comprehended, and grasped by human
reason, and his works and words are not such works and words as should be subject
to human reason.''

``Let us suppose that our text and understanding were also uncertain and dark, as
it is not, just as theirs is (the Zwinglians' text and understanding); then you
still have the glorious defiant advantage that you can with a good conscience
stand on our text, and say thus: If I must and am to have an uncertain, dark text
and understanding, then I will rather have the one that is spoken by the divine
mouth itself than have the one that is spoken by a human mouth; and if I am to be
deceived, then I will rather be deceived by God, if that were possible, than by
human beings; for if God deceives me, he will surely answer for it and give me
redress. Human beings cannot give me redress, but lead me to hell. Such defiance
the fanatics cannot have; for they cannot say: I will rather stand on the text
that Zwingli and Oecolampadius utter in discord than on the text that Christ
himself utters in concord.''

``Therefore you can gladly say to Christ both in your death and at the last
judgment thus: My dear Lord Jesus Christ!
% --- p. 31 ---
There has arisen a quarrel about thy words in the Supper: some will have it that
they are to be understood otherwise than they run. Since they teach me nothing
certain, but only confuse me and make me uncertain, and can in no way either will
or prove their text: so I have abided by thy text, as the words run. If there is
anything dark in it, then thou hast willed to have it dark; for thou hast given
no other explanation of it, nor commanded one to be given. Likewise one finds in
no writing or language that is should mean \emph{signify}, or \emph{my body
sign of the body}.''
% Emphasis mirrors the Danish exactly: „at er skal hedde \emph{betyde}, eller
% \emph{mit Legeme Legemes Tegn}“ --- TWO emphasized spans, with "er"/"is" left
% unemphasized and the second span covering both terms. Drafted first as four
% spans (is / signify / my body / sign of the body) and corrected. The odd
% distribution is the print's, not an error here; do not "regularise" it.

``Should there now be a darkness in it, then thou wilt surely bear with me that I
do not hit upon it, just as thou didst bear with thine apostles that they did not
understand thee in many respects; as when thou didst proclaim thy suffering and
resurrection, and they nevertheless kept the words as they ran, and did not make
others of them. Just as thy dear mother also did not understand when thou saidst
to her (Luke 2:49): I must be about my Father's business, and she nevertheless
simply kept the words in her heart and did not make others of them. So I too have
abided by these words: \emph{This is my body, this is my blood}.''

What a contradiction it is to want at one and the same time to relate oneself
both Christianly and scientifically to Christianity becomes (if possible) still
clearer to us when we look to the theologians' involved dispute about Holy
Scripture. Let us first consider the dispute in its generality, in order the
better to be able to assess the principle that develops in its particulars.

That the Church has its Holy Scripture; that it regards this Scripture as a
Scripture breathed in by God, inspired, so that Scripture's human words are at
the same time God's Word, no one can deny who does not straightforwardly wish to
deny the existence of the Church and of Christianity.

The Church therefore has its Holy Scripture, and posits a difference of quality
between this holy, inspired Scripture and every other
% --- p. 32 ---
writing. Inspiration is Scripture's quality. This quality is inculcated with so
great emphasis by the Reformers that the validity
% The Danish here reads „Refomatorerne“ for „Reformatorerne“ --- printer's
% error, kept as printed there. The English spells it correctly. (Drafted
% first as "Refomers"; corrected, to keep ONE rule for defects across the
% whole translation --- see the note in the header and in RESUME-NOTES.)
of our evangelical Church's confessional writings is even made straightforwardly
dependent on Holy Scripture. The Augsburg Confession, which is after all regarded
as our Church's palladium, is built upon Scripture; it is therefore not to hold
in and for itself alone; but it is to hold \textit{quia et quatenus},
\emph{because} and \emph{insofar as} it accords with the Word revealed in Holy
Scripture.

The theologians, who at bottom wrangle about everything, dispute about everything
and doubt about everything, are of course still in much doubt and dispute about
how this \textit{quia} and \textit{quatenus}, this \emph{because} and this
\emph{insofar as}, is properly to be explained.

If we consider the relation from the standpoint of the Church and of the faith of
the Gospel, the explanation is easy to find out. When the evangelical Church
confesses of its confessional writings that they are to hold only \emph{insofar
as} (\textit{quatenus}) they agree with Holy Scripture: then it has by no means
been the Church's meaning to abandon these confessional writings to theology's
critical-scientific treatment. We know how it went with Melanchthon, who
attempted to alter the 10th article on the Lord's Supper. Melanchthon was the
most respected theologian of his time, indeed the man who had himself drawn up
the Confession; but the faith living in the Church set itself against the
theologian, and the theologian had to give way. The Church could not at any moment
doubt that the confession must be true and trustworthy (how, indeed, should a
church body be able to come forward with a confession doubtful by its own
admission); the Church declares that the Confession stands fast \emph{because} it
accords with Scripture. In this \textit{quia}, in this \emph{because}, lies
unconditional certainty; but in this \textit{quia}, \emph{because}, there is at
the same time expressed the thoroughgoing difference of quality that obtains
between Holy Scripture and the Confession. Scripture is inspired, the Confession
is not. This difference does not bear most closely on the content
% --- p. 33 ---
(for the essential content must at bottom be the same in the Confession as in
Scripture, since the Confession has its content from Scripture); but the
difference of quality points to the connection between the divine content and the
human word. In Scripture this connection is absolutely indissoluble, for Scripture
is holy and inspired; in the Confession, on the other hand, the connection is not
absolute: the same content can be expressed in other, perhaps more adequate, words
and propositions. Individual articles can be added or, in order to make the sense
still clearer, expressed in other words, should the Church at some time under other
conditions find itself called upon to do so. Thereby the expression
\textit{quatenus}, \emph{insofar as}, acquires its significance. Of the relation
between Scripture's word and Scripture's content, on the other hand, no
\emph{insofar as}, no \textit{quatenus}, holds. In Scripture itself the connection
of the divine and the human (content and form, meaning and expression) is in every
way so perfect, so absolute, so infallible and unalterable that no separation can
be thought: Scripture is holy and inspired; every word in the canonical Scripture
is a Word of God, breathed in by God.

But when this conception is made to hold in all its strictness, everyone surely
understands that it is a religious and not a scientific conception. The
contradiction, the glaring contradiction that lies in wanting at one and the same
time to relate oneself religiously and yet scientifically to Holy Scripture,
shines out from orthodox theology's own history.

There is in our Church's history a period that has rightly been called the
bibliolatrous period. Bibliolatry, idolatrous Bible-worship, is quite a telling
expression for the unblessed attempt to relate oneself at one and the same time,
in the same respect and in equal degree, both objectively-scientifically and
subjectively-religiously to Holy Scripture.

For the objective scientific observer the Bible is an
% --- p. 34 ---
outward object; for the religious man it is a holy object: for the scientifically
religious man the Bible is therefore an outward holy object, i.e. an object of
idolatrous worship.

So there sits the scientific but nevertheless orthodox Protestant, and looks
worshipfully at his Bible, just as the scientific but nevertheless orthodox
Catholic looks worshipfully up to his Pope. It is of no use at all that a
Protestant critic subjects the Pope to a scientific assessment and in this
assessment lays bare his faults and sins: the orthodox Catholic believes in the
Pope just as fully, the orthodox Catholic will be sure to defend the Pope's
objective infallibility with objective grounds; for the orthodox Catholic
worships the Pope. Just as little does it help that a Catholic critic --- to the
advantage of the one saving Church --- sets about reasoning against the Bible and
laying bare its human imperfections: the orthodox Protestant scientifically
rejects every scientific objection; the orthodox Protestant will be sure to go on
defending the Bible's objective infallibility with objective grounds; for the
orthodox Protestant worships the Bible.

In Protestantism's orthodox theology the objectively scientific theory of
inspiration was developed with bibliolatrous consistency. The Bible was not merely
inspired, but grammatically-linguistically inspired to the last dot and tittle.
The Holy Spirit had expressly dictated to the holy men every word, every letter
they were to write. The holy men were therefore also called ``God's scribes,''
``God's pens'' (\textit{notarii Dei, calami Dei}).

In the theory of inspiration, then, the faith of the Gospel and theology were to
have found their essential and indissoluble point of union; but behold! in the
theory of inspiration lay precisely the germ of their future separation.

In holding fast Scripture's quality in the
% --- p. 35 ---
unity of the divine and the human effected by inspiration, the faith of the
Gospel is at the same time willing to acknowledge the difference and to conceive
Scripture's human side as it factually shows itself. Faith does not at all aim at
comprehending the relation, but thanks God that it can appropriate the content:
faith relates itself subjectively to the inspired original Word.

But orthodox theology, which related itself objectively and scientifically, now
wanted at the same time to draw scientific conclusions from the religious
presupposition, and thereby brought about its own fall.

The religious unity of the divine and the human in Scripture's word learned
theology then in the first place conceived as a \textit{grammatical} unity. Thus
arose the dispute about the inspiration of the vowel points in the Old Testament
(the Buxtorfs --- Ldv. Capellus). In this dispute orthodoxy suffered a great
defeat.

The religious unity of Scripture's divine content and human form learned theology
next conceived as a linguistic-aesthetic unity. ``If the New Testament was
breathed in by the Holy Spirit, then the New Testament's Greek'' --- so the
theologians supposed --- ``must also be the purest and most perfect Greek.''
% The Danish has no full stop before the closing “ here, as set; the English
% supplies it.
Thus arose the dispute about the Greek style of the New Testament (Seb. Pfochen,
Siegm. Georgi --- Thom. Gataker, Salmasius, Sturz and others). In this dispute
orthodoxy suffered a great defeat.

With respect to the orthodox theory of inspiration we must therefore expressly
note: a) that it dissolved in its own consequences; b) that it weakened the
independence of the Church and of faith; c) that it paved the way for negative
criticism.

Once criticism had first broken through orthodoxy's theological rampart, it
gradually tore down the Church's remaining religious presuppositions, according as
it found that they stood in its way. Under the pretense that the historical facts
were at
% --- p. 36 ---
bottom inessential for faith, criticizing theology understood how by manifold
turns of phrase to fob off the general feeling of piety, and could then all the
more freely place itself in a purely scientific relation to Scripture.

The still moderate, half-scientific and half-consistent theology is a masked
freethinking; resolute freethinking is theology's scientific consequence.

In theology's dispute with the modern consciousness the theologians suffer day by
day one defeat after another, come deeper and deeper into the difficulties, deeper
and deeper into doubt.

Doubtless considered theology can still beat back many an attack that an
over-heated criticism may have directed against the genuineness of one or another
biblical writing; but what is thereby won? Can the theologian ever bring it to the
point where he dares to say: ``The genuineness of this disputed writing is settled,
now at last settled once and for all''? --- By no means! such an assertion is
unscientific, and conflicts straightforwardly with the conditions of historical
faith (\textit{fides historica}). The theologian brings it at the very most only to
a halt, a provisional halt: ``this writing's spuriousness is not yet settled!'' But
if the genuineness is not quite settled either, then it is after all subject to
doubt; but if it is subject to doubt, it is after all lost for faith; for faith
goes out where doubt comes in.

So now the freethinker and the theologian contend --- under the same
presupposition, under the same scientific principle: only the language and the mood
differ. The freethinker is resolute: ``The spuriousness is settled!'' The
theologian is irresolute: ``The spuriousness is not yet settled.'' The freethinker
has in the name of science exchanged faith for unbelief. The theologian has in
science lost faith, and is despite all scientific exertion no longer able to raise
his doubt to faith. The freethinker relates himself
% --- p. 37 ---
non-religiously to Scripture; the theologian can never bring it to the point where
he, as a scientific theologian, relates himself religiously.

If the theologian will hold to the Church, he must openly break with science; if he
will hold to science, he must openly break with the Church: this is the final
dilemma.

If the Church is to maintain its independence, it must guard itself in equal degree
against the theologian's defense and the freethinker's attack. The negative weapon
with which the Church defends itself is the fundamental thought \emph{that
Christianity is higher than science, and the faith of the Gospel heterogeneous with
theology.}

%% ============================================================
%%  FOURTH LECTURE (pp. 38--51)
%% ============================================================
\chapter*{Fourth Lecture}
\addcontentsline{toc}{chapter}{Fourth Lecture}
\label{ch:forel4}
\markboth{Fourth Lecture}{}

\begin{center}\itshape
For scientific theology the Bible is a historical document of antiquity
whose heterogeneous constituents must submit to ratiocinating criticism;
for the believing conception the Bible's wondrous content is exalted above
criticism, since it is impossible for the observer to relate himself
critically and ratiocinatingly to the miracle: the faith of the Gospel is
heterogeneous with theology.
\end{center}

% --- p. 38 ---
The Bible, the Church's Holy Scripture, has in many ways been misunderstood,
misused, and maltreated by human folly; but hardly has it ever before been so
misunderstood, misused, and maltreated in principle as by theological criticism
in recent times.

Many fanatics have misconstrued Scripture to their own perdition; but with all
that they have nevertheless related themselves to it as to a holy Scripture.
Many freethinkers have scoffed at Scripture --- to their own perdition; but with
all that they have nevertheless related themselves to it as to a holy Scripture.
For is it not precisely the mark of the holy that it tempts the fanatic and
provokes the mockery of unbelief? --- Why does the fanatic take refuge in the
Bible? Because from this ``holy inspired'' book he means to prove to us that he
himself is inspired. --- Why does the freethinker mock at the Bible? Because it
is now so great an offense to him that this old book with its ``many myths'' and
``holy fairy tales'' should still be held and count as inspired.

% --- p. 39 ---
In relation to the fanatic and the freethinker, then, the Bible still preserves
its original quality; but not so in relation to criticizing theology.

``The critical introduction'' will of course conduct itself neither as the
fanatic nor as the freethinker. ``The critical introduction'' is a punctilious
person who in learned caution has little by little taught itself to stand
impartially \textit{above all party}, and supposes it thereby preserves
religious equilibrium in scientific many-sidedness. ``On the one hand one
cannot'' --- so it runs --- ``indeed deny that there are found in the Bible many
spurious constituents, many mythical legends and superstitious narratives, which
must straightforwardly be referred to the popular fantasy and limited circle of
representations of that time. But on the other hand neither can one deny that
there are found in the Bible so many great and glorious thoughts, so many
exalted examples and elevating traits of character, that the book as a whole must
awaken the greatest reverence in every feeling mind, indeed that even the
scientific man must confess that this Book of books surely cannot, on the whole,
have been written without a special cooperation of God.'' --- Roughly thus reasons
``the critical introduction.''

But what sort of relation is it, then, in which ``the critical introduction''
thereby places itself to the Bible? Yes, it is a highly ambiguous relation, a
relation that is at one and the same time to be both religious and yet
scientific, and precisely for that reason is neither the one nor the other.

And what does the Bible then become in this ambiguous relation? It becomes an
ambiguous book, a book that both is inspired and yet is not inspired, a peculiar
book that is both the one and the other, and therefore at bottom neither the one
nor the other.

In this ambiguity, then, the Bible loses its quality; and therefore I say that
theological criticism's relation to the Church's Holy Scripture --- despite the
aberrations of fanaticism and the
% --- p. 40 ---
blunders of freethinking --- must be determined as the properly fundamental
misrelation, as the misrelation that comes nearest to untruth.

Holy Scripture is no objective scientific book, and one can as a consequence not
relate oneself objectively and scientifically to it either. What the Bible is
cannot be settled once and for all in and for itself; it must each time be
settled in and by the relation in which the reader places himself to the Bible.

\emph{As I myself am, so is my Bible, and as my Bible is, so am I myself.}

If I now relate myself believingly to the Bible, then the Bible is also for me a
book of revelations, a holy inspired Scripture that shows me what I am and how I
am, and casts light over my whole existence. My understanding may now doubtless
discover wondrous oppositions, indeed contradictions, in this holy book. The
heterogeneous is here put together into unity: the outwardly historical and the
inwardly psychical, the supra-worldly and the worldly, omnipotence and frailty,
the holy and the sinful are found here together in singular mixture and most
wondrous interaction. This may offend the unbelieving, but is nevertheless to
true edification for the believer. In the believer's life there is, after all,
something corresponding: here too the heterogeneous is bound into unity; here too
the outwardly historical comes to the inwardly psychical; here too the worldly
man is to relate himself to the supra-worldly, the frail to the almighty, the
sinful to the holy, he who is dust and ashes to Him who is from eternity and unto
eternity. In the believer's life too there are contradictions. Insofar as the
believer now learns to understand the contradictions that are in his own life, he
will also learn to understand the apparent contradictions that are in the holy
inspired Scripture; and conversely, insofar as he learns to understand the
contradictions in Scripture, he will with the Spirit's assistance also learn to
understand the contradictions in his own
% --- p. 41 ---
life. If those holy men, although they were God's elect, nevertheless had to go
on believing with hope against hope, how much more must this not hold of the
religious man, who may not so call himself God's elect! And when the religious
man sees that in this world one must believe with understanding against
understanding; when he learns to see that this law becomes stricter and stricter
in proportion as the life of faith becomes more developed and more variously
illuminated by reflection: should he then not also understand that those holy men
walked as the Scriptures witness only because they walked in faith and in hope,
because they walked as God's elect. The relation is therefore this: the more
inwardly the religious man exists in faith and understands faith's conditions, the
more inwardly does he recognize that the Church's Holy Scripture is an inspired
Scripture.

Even were the believer otherwise to become so distinguished in human learning and
insight that he could, if it came to that, criticize all the wise men of the world
and set himself up as their teacher: never will he bring it to the point of
criticizing Holy Scripture and setting himself up as its teacher; but the reverse
relation will increase more and more, namely that Holy Scripture in an ever higher
sense becomes his teacher, the comforter and disciplinarian who endures with him
unto the end.

\emph{As I myself am, so is my Bible, and as my Bible is, so am I myself.}

The believer can now provisionally abstract from his faith, enter into science,
and place himself in an objective, literary, critical-scientific relation to
% The Danish here reads „lil“ for „til“ --- printer's error (a plain Fraktur l
% for the t), kept as printed there. The English has the correct word.
the Bible. Against this there is doubtless, from the standpoint of cultivation,
nothing to object. If the historical-philological critic were forbidden to relate
himself critically and scientifically to Holy Scripture, then the philosopher
would also have to be forbidden to relate himself speculatively and scientifically
to Christianity. If biblical criticism were straightforwardly forbidden, then the
philosophy of religion would also have to be straightforwardly forbidden;
% --- p. 42 ---
but that by no means follows from the principle set up here. We inveigh as little
against a true scientific biblical criticism as against a true scientific
philosophy of religion. There is nothing to prevent Christianity from becoming an
object of dialecticizing speculation; but the speculative philosopher must then
know what it is he is doing. In placing himself in a dialectical-speculative
relation to Christianity, the philosopher is outside the relation of faith; for
the relation of faith is qualitatively different from the speculative relation of
knowledge. As the relation now changes its quality and becomes a quite different
relation, so Christianity too changes its quality and becomes something quite
different. For the believer, Christianity is God's wondrous provision for the
salvation and blessedness of sinners; for the speculating man, Christianity is a
magnificent phenomenon, a phenomenon in which new and wondrous ideas have shown
themselves, and the speculating man now seeks to fathom these ideas' essential,
metaphysical relation to the whole system of world-ideas in which we live and
breathe and think. --- As it goes with the philosophy of religion, so also with
biblical criticism. The religious relation of faith to Holy Scripture is
qualitatively different from the literary-critical relation. As the relation now
changes its quality and becomes a quite different relation, so the Bible too
changes its quality and becomes a quite different book. For the believer the Bible
is a holy inspired Scripture; for the literary critic it is neither more nor less
than a collection of writings of antiquity in the history of religion, a document
of antiquity which in principle is to be conceived from the same point of view
from which we conceive, say, the Indians' Veda, the Persians' Zend-Avesta, the
Muhammadans' Koran.

Here now a mass of tasks presents itself for learned science: to compare codices,
versions, and readings, to communicate archaeological and philological
% The Danish here reads „medele“ for „meddele“ --- printer's error (single d),
% kept as printed there. The English has the correct word.
information, and so on. That these
% --- p. 43 ---
tasks are of a purely scientific kind is seen from the fact that they can be
treated by anyone without regard to confession of faith. What our learned critics
and philologists have accomplished in this respect, what a Wetstein, a Griesbach,
a Lachmann, a Gesenius, a Winer have achieved in this direction, has its validity
and scientific warrant without regard to the standpoint of faith at which those
learned men have found themselves. Whether the critical philologist belongs to the
Reformed or the Lutheran Church, whether he is orthodox or a rationalist, indeed
whether he is at all a believer or a non-believer, makes here no difference
whatever to the matter. The learned critic does not, after all, inquire in the
quality of a believer, but in the quality of a knower; here the question is only
of talent and knowledge, not of religiosity. What the learned inquirer achieves is
to be tested and judged by the experts in the school, but not by the believers in
the congregation.

If it now stands fast that there is a double relation to the Bible, a religious
and a scientific one, if it stands fast that the Bible as a consequence changes
its quality with the changed relation: then it will surely also be necessary to
reflect a little more closely on this change of quality before one entrusts
oneself without further ado to the scientific guidance that is nowadays offered in
``the critical introduction.''

In accordance with the Bible's double quality, as inspired Scripture and as
literary document of antiquity, biblical criticism must necessarily come upon
tasks of the kind that lie on the boundary of the religious and the scientific,
and which we may therefore for the present call the mixed tasks.

To these we reckon: a) The settling of which books belong to the canon. This task
has a religious significance; for on it depends how many books one is to take as
inspired: the question of the canon thus becomes a question of faith. But
conceived from the historical-critical point of view
% --- p. 44 ---
it becomes at the same time a learned scientific question, namely that of the
books' genuineness (authenticity). Further: b) The investigation of the books'
trustworthiness, insofar as this trustworthiness can be built on the mutual
agreement of the accounts: the question of the harmony among the 4 Gospels is a
mixed question. Finally: c) The settling of all those points in which the holy
authors, with respect to the purely historical, diverge from the profane writers.

It is doubtless these mixed tasks that have led introduction-criticism to place
itself in that relation of ambiguity to Holy Scripture spoken of above, a relation
under which faith is neutralized with science. But in this relation the mixed
tasks cannot be solved; for if the mixed tasks are solved in ambiguity, the
solution must itself become ambiguous.

Before one attempts to solve these tasks, therefore, one must first of all be
agreed with oneself as to whether it is really under faith's or under science's
presupposition that the solution is to take place. For according as one chooses
the one or the other of the qualitative relations, the solution too will have an
entirely different outcome.

I. \emph{Let us now first choose the relation of faith.} The religious factor is
thus the decisive one, the historical varies. In order to see what lies in this,
we choose, say, the investigation of the four Gospels' genuineness.

That the historical factor (\textit{fides historica}) varies here everyone will
grant who knows a little of criticism. It is possible that our Greek Gospel of
Matthew was not written by Matthew, but by another author standing in nearer or
remoter contact with the evangelist himself. It is likewise possible that the
evangelist John did not himself edit the Gospel that bears his name, and so on.
Here there is now room for a multiplicity of investigations, hypotheses, and
explanations.

% --- p. 45 ---
The religious factor, on the other hand, is not varying, but constant. For
subjective faith (\textit{fides religiosa}) it stands unshakably fast that the
four Gospels (whoever may have written them) are inspired. This ecclesial
presupposition faith cannot give up without at the same time giving up its
principle.

If we now put both factors together, then the conditioned is determined by the
unconditioned, the relative by the absolute, and the task's solution can then be
expressed thus: the four inspired Gospels are written by four inspired authors;
these authors are named after Matthew, Mark, Luke, and John. The Greek
superscriptions \textit{Κατα Ματθαιον, Κατα Μαρκον} etc. fit here in the most
exact way.

The same procedure is to be applied when the question is of the four Gospels'
mutual agreement (\textit{harmonia evangeliorum}).

The historical factor varies. Between the four Gospels there is not only a
remarkable agreement, but at the same time remarkable divergences and apparent
contradictions (the genealogy of Jesus in Matthew, for example, criticism cannot
get to agree with the genealogy in Luke). One little word, one small alteration,
and the harmony would be there; but --- the little word is lacking, even the
happiest explanation is nevertheless only a subjective conjecture, and so
hypothesis alternates with hypothesis again.

The religious factor is constant. As the accounts stand, so shall they remain
standing with unalterable divine authority; here no reflection is made on any
disagreement: the one account is for the believer himself exactly as infallible as
the other.

If both factors are now put together in the unity of the faith of the Gospel, the
result can be expressed thus: in the four Gospels the divine content is presented
in such a human form that the essential agreement between the identical accounts
has at the same time taken on the stamp of apparent divergences in
% --- p. 46 ---
particular inessential points.\footnote{``The jurists say that a capital crime
swallows up all the lesser ones --- so it is with faith: its absurdity utterly
swallows up the petty. Disagreements that would otherwise have a disturbing effect
do not disturb here, and make no difference to the matter. On the other hand it
makes a great difference to the matter if someone by a calculus of pettiness
wishes to auction faith off to the highest bidder; for it makes so much difference
that he never comes to faith.'' Johannes Climacus.} But since the Gospels are
inspired, it is impossible for the believer to separate out the divine from the
human; both moments must therefore go on working together until the day comes that
clears up the harmony.

What has already been said of the confession of faith holds at the same time of
the Gospels, indeed of all the canonical books in the Church's Holy Scripture: the
danger is not that a single article or a single narrative should by a human error
have crept in; but the danger is that a critical separation of the inseparable is
made, whereby inspiration is abolished and faith disappears.

This solution, which is so satisfying for the faith of the Gospel, cannot possibly
be satisfying for theology; for theology is, after all, a science.

II. \emph{If we now choose the scientific relation,} and attempt from a purely
objective point of view to solve the mixed tasks, the matter at once takes another
turn. The religious factor then falls away, and the historical is forfeit to
criticism. Here there is now nothing fixed, nothing unconditioned: suspicion and
doubt know no bound. The explanation then ends with the result that not a single
Gospel is genuine, indeed that these Gospels first received their present form one
and a half or two hundred years after Christ.

This result, which has no significance whatever for subjective faith, is on the
other hand in a critical respect decisive for
% --- p. 47 ---
historical faith, and shows plainly enough how the historical factor evaporates
once it has stepped out of the right connection with the religious. But when
dissolving criticism, after it has by this procedure undermined the historical
foundation of faith in revelation, then supposes that it can at the same time
undermine the faith of the Gospel, this is --- as we have already
proved\footnote{P. 20. P. 24. Pp. 28--29.} --- a great misunderstanding. The
believer does not let himself be disturbed by the scientific declaration that the
biblical writings cannot stand before criticism's tribunal; he sees in this only a
confirmation of the truth that the Bible's content is a supra-scientific content, a
content to which he must relate himself religiously.

\emph{As I myself am, so is my Bible, and as my Bible is, so am I myself.}

If the Bible is inspired, then it must have its quality in a miracle; but if the
Bible has its quality in a miracle, then it is at the same time also settled that
in this quality it permits no scientific-critical conception.

A scientific-critical relation to a miracle is nonsense. This holds not only of
such an inward and psychical miracle as inspiration; but it holds at the same time
of the more outward and, so to speak, more conspicuous miracles.

The miracle presupposes an interaction of heterogeneous factors, of the
supernatural and the natural. If this interaction is a delusion, the miracle itself
must be a delusion; but if it is true, everyone must surely easily be able to see
how impossible it is to place oneself in an objective relation to an actual
miracle, to become a spectator at a miracle as at a physical event, or to relate
oneself scientifically and critically to a narrative of a miracle as to any other
historical account.

If we consider, for example, the miracles that seem to intervene most
% --- p. 48 ---
in the sensuous, the miracle by which Moses makes the spring gush forth in the
wilderness, and the miracle by which Christ feeds the hungry in the wilderness: is
not the reciprocal relation of the natural and the supernatural recognizable
enough?

Moses lifts up his staff and strikes the rock; that a spectator can see with
natural eyes, for it happens naturally. And the rock ``weeps,'' and the spring
wells up and streams, and the Israelites come to slake their thirst: that a
spectator can see with natural eyes, for it happens naturally. If the spectator
will not believe his own eyes, he can besides attain an immediate sensuous
certainty by himself bending down and drinking of the spring. But with all this the
spectator has nevertheless not come a hair's breadth nearer to the miracle itself.
There is in this event something that not even the most scientific observer would
be able to discover, and yet it is precisely this something on which the whole
conception essentially depends: it is the causal connection between the staff of
Moses and the gushing spring. The condition is natural; but the efficient cause,
what is it: natural or supernatural? If the spectator assumes the former, then the
event stands for his imagination still only as a mere \textit{mirabile}; if he
assumes the latter, it is transformed into a \textit{miraculum}.

Which of these two assumptions the observer is to prefer cannot be settled by
scientific criticism. The actual miracle can bring the spectator to a point where
he must say to himself: ``if I am to believe my senses here, I cannot believe my
understanding; if I am to believe my understanding, I cannot believe my senses.''
But further than to this dilemma the phenomenon cannot bring any spectator: the
objective miracle repels the observer from itself, drives him back into inwardness,
and leaves the decision to depend on the judging subjectivity.

Holy Scripture now tells us not only of the many miraculous works, but at the same
time describes for us the impressions these
% --- p. 49 ---
various miraculous works made on the various eyewitnesses. How various are these
impressions! Some believe unconditionally, and are saved by faith; but faith is in
subjectivity. Others are startled and terrified; but the utterance of terror is
conditioned by subjectivity. In the people's cry: ``a great prophet is risen up
among us!'' --- ``God hath in mercy visited his people!'' we perceive a hovering
between mere astonishment and inward faith; but this hovering takes place in
subjectivity. The crowd is curious, thoughtless, lost in the sensuous (thus the
many who came to Moses and drank of the spring in the wilderness; thus the many who
came to Christ and ate of the loaves in the wilderness); but this superficial
curiosity, this thoughtlessness, this animal torpor is a sorry defect in
subjectivity. Some are offended; from the passion of offense there then springs in
turn the explanation of offense (as when the Pharisees said of Christ: ``he casteth
out devils by Beelzebub''); but offense is in subjectivity.

An actual miracle is really enough to make one take leave of one's understanding.
But human beings take leave of their understanding in two ways: partly in such a
way that they straightforwardly go mad, partly conversely in such a way that they
give up their understanding for something higher and precisely thereby escape
madness. This is perhaps even, according to the Bible's meaning, the most exact
expression for the miracle's divine purpose: that human beings are thereby to be
brought to take leave of their understanding without, however, straightforwardly
losing their understanding.

The believer who believes the miracle really does take leave of his understanding,
insofar as he believes --- against understanding. That he nevertheless preserves
his sound understanding can be seen simply from the fact that, as regards merely
intelligible matters, he speaks and acts entirely sensibly. The offended man who is
offended at the miracle likewise takes leave of his understanding, insofar as in a
kind of passionate derangement he attempts to explain away the inexplicable; but
since the deranged
% --- p. 50 ---
explanation nevertheless always affords him a certain subjective relief, he
nevertheless preserves his understanding in all merely intelligible affairs, and
consequently does not go mad. The Pharisees' explanation of Christ's miracle was
doubtless a deranged explanation of its kind; but the offense gave itself vent in
the subjective outburst of subjective passion, and the offended kept their natural
understanding.\footnote{If the miracles of which the Bible tells us really repeated
themselves in our day, the theologians would --- in accordance with the principle
--- without doubt be among the first who stood in line to go straightforwardly mad.
This formation of the miracle's effect is entirely unknown in Holy Scripture.
Scripture knows only the subjective ones: \emph{faith and offense}, but madness it
does not know. That the Bible does not know such as went straightforwardly mad over
the miracle must surely be explained by the fact that it does not at all know such
as have attempted to relate themselves scientifically and critically to the
miracle, in order to get the miracle scientifically and critically justified.}

As it now stands with becoming an eyewitness to a miracle, so it stands also with
having to narrate and describe a miracle. The miracle can be grasped and
appropriated only in faith; it can be narrated, described, and rightly presented
only by such as stand on faith's highest step, i.e. by inspired authors.

In the narratives of miracles the two heterogeneous factors will then again show
themselves in their reciprocal relation; faith's subjective factor will again be
the absolute, the constant one; the historical factor the relative, the varying one.

When, therefore, two evangelists narrate the same miracle in a somewhat different
way, the divergence will not touch the miracle's religious kernel, but will touch
only the phenomenal attendant circumstances and psychological nuances.

As it stands with having to narrate a miracle,
% --- p. 51 ---
so must it surely stand also with having to read, conceive, and appropriate a
narrative of a miracle.

\emph{As I myself am, so is my Bible, and as my Bible is, so am I myself.}

It is this leading thought that I have striven to the best of my ability to hold
fast and develop in my writing \emph{Gospel Faith and the Modern Consciousness}.

There are only two pure, qualitative relations to Holy Scripture, the believer's
and the freethinker's.\footnote{A doubtless orthodox pastor and much esteemed
convention speaker has, as I see, been busy spreading the opinion that in my
writing ``Evangelietroen og den moderne Bevidsthed'' I am supposedly on the way to
establishing two mutually contradictory and yet equally warranted truths, and that
this would surely show itself if this view of mine were otherwise given occasion to
clarify itself. With so fleeting and, as it appears to me, so frivolous an
expression of opinion I can of course not further concern myself; I permit myself
only most respectfully to remark that if the esteemed speaker had not merely glanced
into the book's preface, but had really taken the time and read the book itself,
then he would without doubt have seen how here too, at the conclusion, it goes with
the cunning ``Caiaphas.'' The explanation that the World-Spirit (cf. VII. Jesus
Christ is the Savior of the world) at last gives of the modern dogma that
\emph{everything must happen naturally} is, even if not exactly very edifying, at
any rate sufficient to exhibit the consequences in which the freethinker's supposed
\emph{truth}, if it is otherwise given occasion to clarify itself, must at length
end.} The criticizing theologian's relation of faith is an unclear, ambiguous, and
impure relation. For that reason I have in the said writing --- however gladly I
should otherwise have wished it --- not been able to be so fortunate as to speak to
the theologians' liking; for it is truly my conviction that \emph{Christianity is
higher than science, and the faith of the Gospel heterogeneous with theology.}

%% ============================================================
%%  FIFTH LECTURE (pp. 52--65)
%% ============================================================
\chapter*{Fifth Lecture}
\addcontentsline{toc}{chapter}{Fifth Lecture}
\label{ch:forel5}
\markboth{Fifth Lecture}{}

\begin{center}\itshape
Theological interpretation of Scripture misconstrues the peculiar character
of revelation, and dissolves the content of faith into science's explanation;
religious interpretation of Scripture reduces science to a means for a
believing appropriation of the revealed Word: the faith of the Gospel is
heterogeneous with theology.
\end{center}

% --- p. 52 ---
If the Bible, the Church's Holy Scripture, is not to be lost in theology's
``critical introduction,'' the relation of faith must be posited unconditionally,
and inspiration maintained.

But even then, when the Bible in the quality of holy, inspired Scripture has come
safely through theological criticism; even then, when it has come so near to the
reader that no learned hypothesis can force itself in between the reading eye and
the text before it; even then, when the Bible lies before us as an open book in
which the believer reads the word of revelation; even then there is still a
difficulty to overcome, a difficulty that, far from being diminished, is on the
contrary increased by theology's theological assistance: this difficulty lies in
\emph{interpretation}.

What use is it that a believer reads the Bible, word for word, if he does not
rightly understand what it is that he is reading! To read without understanding is
meaningless; to read oneself into a misunderstanding is worse than not reading at
all: the Bible must therefore be construed, explained, and expounded, so that its
divine content
% --- p. 53 ---
can be duly understood and rightly applied: the Bible must be interpreted.

To interpret a human writing, a writing in which a superior human mind has laid
down a treasure of original thoughts, is already a great art; the task becomes
harder still when such thoughts and ideas are not systematically presented in
science's universally valid form, but take shape in the form of the imagination, in
a multiplicity of characters and situations, in the form of life, where thought is
colored by passion and language varies with mood.

If it is already a great art to enter duly into an original human writing and to
interpret such a writing as it ought to be interpreted: what then shall we say of
the task of entering duly into the holy inspired Scripture and interpreting this
Scripture as it ought to be interpreted?

If it must be granted that the Bible has a superhuman and yet human origin, a
superhuman and yet human content, then it also follows from this that the Bible's
right interpretation must be not merely a human art, but at the same time a
superhuman gift.

Insofar as biblical interpretation is now a human art, it must be capable of being
formed by science and built on scientific rules; but insofar as it is conditioned by
a superhuman gift, science does not suffice: everything depends on ``the Spirit's
testimony.''

If there is any point in Christian knowledge where faith and science touch one
another, it must surely be the interpretation of Holy Scripture.

It is not with the science of interpretation as with introduction-criticism. ``The
critical introduction'' with all its bother and quackery the believer can cast from
him at whatever moment he will; these artificial hypotheses and hypothetical arts
% --- p. 54 ---
at bottom do not concern him at all; the Bible is once and for all what it is and
as it is; we have no other Holy Scripture, so Luther's defiance holds: ``God's Word
they must let stand, and get no thanks for it.''

``The critical introduction'' has accordingly also only really been developed in the
last century; it is a branch of science that recent theology has chiefly cultivated
for its own scientific pastime. It is otherwise with hermeneutics and the guidance
science thereby offers the religious interpreter of Scripture. This guidance no one
may spurn who will make himself familiar with the Bible's human form and expression.

The faith that in this matter despises human guidance is no true and sound faith of
the Gospel, but a self-willed, obstinate, arbitrary faith that easily degenerates
into fanaticism. ``This one thing'' --- says Luther --- ``exalts holy Christendom
and the Church, that one understands Holy Scripture; this one thing depresses
Christendom and the Church, that one does not understand Holy Scripture . . . . All
other arts and crafts have their preceptors and masters, by whom one must let
oneself be taught according to order and law; --- only Holy Scripture and God's Word
must submit to everyone's arrogance, conceit, obstinacy, and presumption, must let
itself be mastered, twisted, and construed as each understands and wills after his
own head; hence so many fanatical parties, sects, and offenses, God have mercy.''

In order, then, to set a dam against subjective arbitrariness, one has --- and
rightly --- treated the art of interpretation scientifically; and theology has then
founded on this its ``theological hermeneutics.''

But as scientific hermeneutics has organized itself theologically, here too the
deception has crept in, the theological deception that the Christian is supposed to
be homogeneous with the scientific, the theological deception that by continued
scientific
% --- p. 55 ---
application of scientific rules one is supposed to penetrate deeper and ever deeper
into Scripture's divine content. A sorry error! for before one knows a word of it,
science has again gained power over Christianity, and is in full course of hollowing
out faith's kernel.

For it is characteristic of hermeneutics that it can set up such rules as must be
approved by all who will think about their Bible, such rules of interpretation as we
must from science's general point of view accord general validity, without its being
therefore in any way able to govern or determine the particular application that the
individual interpreter in each individual case can and ought to make of these rules.

\emph{Scripture} --- it is said --- \emph{shall interpret itself.} This rule of
interpretation was already pronounced at the Diet of Speyer in 1529, and is to that
extent a genuinely Protestant rule. But when the rule is set up as a scientific
presupposition, one can of course draw from it the most various, indeed entirely
contradictory, religious conclusions. The orthodox theologian, the heterodox, the
rationalist, the speculative theologian, the freethinker, they can all agree on the
Protestant rule; but in applying the rule they nevertheless arrive at quite
different results!

\emph{The particulars in Scripture} --- it is said --- \emph{must be explained from
Scripture's whole: the analogy of Scripture is the supreme rule for interpretation.}

A genuinely Protestant rule! but when the rule is set up as a scientific premise,
theology draws from it the most contradictory religious conclusions. The orthodox,
the rationalist, the speculative, the freethinker, they can all agree on the rule in
general, and even so each brings out his own particular view of Scripture's whole.

% --- p. 56 ---
This disagreement is not accidental or merely transitory, but deeply grounded in the
nature of the case. The conception of Scripture's whole is precisely conditioned by
the conception of the individual passages in their mutual connection; but the
conception of this connection is in turn conditioned by the emphasis one will each
time lay upon this or that individual passage. In one place Christ says: ``I and the
Father are one''; but in another place: ``the Father is greater than I.'' In some
places Scripture speaks of the Holy Spirit as of a personal, independent Spirit (in
the name of the Father, the Son, and the Holy Spirit); but in other places the
Spirit is called merely ``the power from on high,'' ``the power of the Highest.''
--- In some places an extraordinary weight is laid on human freedom: ``work out your
own salvation with fear and trembling''; but immediately thereupon human freedom and
human works are again reduced to a pure nothing, a nothing before God: ``it is God
who worketh in you both to will and to do of his good pleasure.''

These different statements stand sharply against one another; the one cannot be
derived from the other, cannot be limited by the other, cannot be straightforwardly
explained by the other; the understanding cannot mediate them, every attempt at
mediation has shown itself to be a conceit, a self-delusion.

Religious understanding is therefore in an absolute sense dependent on subjective
faith:

\emph{As my faith is, so is my interpretation of Scripture, and as my interpretation
of Scripture is, so is my faith.}

But subjective faith, which in reciprocal relation to the Gospel's Spirit yields the
highest and absolute principle of interpretation, is entirely heterogeneous with
``hermeneutics''' objective rules.

If the interpreters of Scripture will nevertheless make these heterogeneous factors
homogeneous by drawing religious conclusions from scientific
% --- p. 57 ---
premises, or conversely scientific conclusions from religious premises: then it
holds in truth, as Christ says, that what is revealed to babes is hidden from the
wise and the scribes.

We do not need to go far to find examples; they lie only too near at hand, and are
unhappily only too numerous. To name but one example, I will here recall the
theologian \textit{De Wette}.

The recently deceased theologian \textit{De Wette} has assuredly distinguished
merits both as a critic and as a scientific interpreter of Scripture. I should only
dishonor myself if I here allowed myself to dishonor his memory by reckoning him
among those who slight Christianity or treat Scripture frivolously. On the contrary!
\textit{De Wette} was doubtless, according to the present age's way of thinking, as
serious and upright as he was diligent and gifted an inquirer. But in a religious
sense --- I say it without reserve --- in a religious-ethical sense his concept of
revelation is entirely mistaken, and his theological inquiry, in relation to this,
built on a misunderstanding.

Or is it not a misunderstanding when \textit{De Wette} will deny the prophets that
gift of the Spirit which Christ and the apostles ascribe to them? is it not --- in a
religious sense --- a misunderstanding when he will have the prophetic predictions
grounded chiefly on vague wishes and indefinite hopes?%
% DEFECT (Danish p. 57): the opening „ before „Weil“ is never closed --- the
% footnote ends at „256.“ with no closing “. The English closes it. The German
% quotation is kept in German, as Nielsen printed it.
\footnote{``Weil die Idee in den Propheten durchaus vorherrscht, so sind ihre
Vorhersagungen zum Theil nur als Hoffnungen und Wünsche, . . . und sind fast immer
unbestimmt und schwebend.'' Lehrbuch der historisch-kritischen Einleitung in die
Bibel. 1. p.~256.}
Or is it not a misunderstanding when \textit{De Wette} allows himself to criticize
the prophet Ezekiel as one criticizes a mediocre author in modern literature? when
he, for
% --- p. 58 ---
example, censures this prophet's lack of depth and richness of spirit and of great
thoughts? when he laments that Ezekiel's prophetic style is too low and too diffuse
and sunk to a flat prose? when he dwells on the fact that the prophet does indeed by
allegorical and symbolic inventions wish to raise himself above the ordinary, but
usually lapses into the overloaded, the far-fetched, the confused?%
\footnote{Was bei ihm (Ezechiel) am meisten auffällt, ist seine levitische
Gesinnung, vermöge deren er auf heilige Gebräuche einen hohen Werth legt, und sich
selbst in Bildung von Idealen nicht darüber erheben kann. Daher auch sein Mangel an
Tiefe und Reichthum des Geistes und an großen Gedanken . . . . Die prophetische Rede
ist bei ihm zur niedrigen, weitschweifigen, matten Prosa herabgesunken; nur in
symbolischen und allegorischen Dichtungen will er sich über das Gemeine erheben,
verfällt aber gewöhnlich in das Ueberladene, Gesuchte und Verworrene. Op.\ cit.\
pp.~282---83.}
What scientific superiority, what critical self-confidence! Can a theologian then so
master a prophet?

It is a sorry misunderstanding, occasioned by the fact that one will relate oneself
scientifically-religiously to God's Word, and interpret the inspired authors like
other writers of genius. Interpretation transforms itself into criticism: one exalts
the one prophet and abases the other. And yet the concept of revelation is supposed
to remain inviolate in all this, and the prophets all in a certain manner regarded as
inspired. What a contradiction, what uncritical confusion!

If the prophets are not inspired, if hermeneutics and criticism can really vouch to
us that the Bible's inspiration really signifies nothing at all: why then do the
theologians not say so straight out, and give the freethinkers their due, as free
scientific thinkers? But if the Bible is inspired, if the prophets are together
inspired by God's Spirit: how then can theological criticism be so uncritical as to
apply the same measure to inspired authors as to uninspired ones?

% --- p. 59 ---
If the prophets are inspired: then the difference that with respect to style and
presentation may possibly be found between an Isaiah and an Ezekiel must, in relation
to inspiration itself, be an entirely inessential, a vanishing difference.

Theological science, which imagines that it speaks with earnestness and dignity of
the beautiful, the lofty, the deep, that is supposedly to be found in the genuine
Isaiah in contrast to the prophets of a lower manner of writing, offends Isaiah in a
religious sense just as surely by its supposed praise as it offends an Ezekiel or a
Daniel by its presumptuous censure.

The prophet does not at all have his quality in style's aesthetic excellences, but
the prophet has his quality in this, that he is inspired by the Lord's Spirit. One
shall believe the prophet's word and act accordingly; but one shall not reason about
his style or at once declare him spurious because one cannot scientifically come to
terms with his manner of writing.

In order finally to remove this theological misunderstanding with the confusion
flowing from it, one must in hermeneutics, as in introduction-criticism, sharply
separate the scientific conception from the believing appropriation, and therewith
also acknowledge that these conceptions are so heterogeneous that no direct
transition can be made from the one to the other.

When the interpreter relates himself merely scientifically to the Bible, the Bible is
to be treated philologically and historically like any other book of religion from
remote antiquity. This scientific relation of the interpreter to the Church's Holy
Scripture has been brought out in the most definite terms by L.\ I.\ Rückert in the
preface to his commentary on Paul's Epistle to the Romans, where it is said among
other things: ``The interpreter of the New Testament . . . has, as such, no
theological system at all and may have none, neither a dogmatic nor a system of
feeling; he is, insofar as he is an exegete, neither orthodox nor heterodox, neither
supranaturalist nor
% --- p. 60 ---
rationalist, nor pantheist, nor any --- ist whatsoever that there might otherwise be;
he is neither pious nor ungodly, neither moral nor immoral, neither tender-feeling
nor unfeeling; for he has merely the duty of finding out what his author says, in
order to hand this over as a pure result to the philosopher, the dogmatician, the
moralist, the ascetic, and so on. --- So as an exegete he may not have any other
interest either than this one, to understand the apostle rightly, to grasp his
thoughts purely without foreign admixture, and to lay them before the reader just as
pure and unmixed. . . . For him (--- the scientific exegete ---) it must be a matter
of indifference whether Paul speaks truth or falsehood, whether it is a moral or an
immoral spirit that runs through his letters, whether his teaching is salutary or
fundamentally pernicious; for he (the interpreter) is only to present what he (i.e.
the author) says, what spirit pervades him, what doctrine he propounds.''

Before Rückert (whose commentary on Romans appeared in Leipzig in 1831) several older
hermeneuticians (G.\ L.\ Bauer, Bretschneider, Griesbach and others) had already
inculcated the scientific freedom from presuppositions that ought to distinguish the
unprejudiced interpreter of the Bible. ``If one were'' --- it is said --- ``to assume
the New Testament beforehand to be revelation, and accordingly to accept only that
sense as the hermeneutically true one which seems worthy of a revelation, then one
would be moving in a circle; --- with equal right one would then expound every other
document of revelation, such as the Koran or the Zend books, in this way, and could
thereby establish its content as logically correct and true.'' ---

Thus the theologians themselves were, in their zeal for an unbiased scientific
interpretation of Scripture, well on the way to giving up the theological element in
theology. But since theology consists precisely in the contradiction of having to
make the heterogeneous homogeneous, since theology cannot give up the religious
without giving up
% --- p. 61 ---
its character, it followed of itself that that scientific consistency could not be
carried through either.

Instead of holding fast to the principle and letting the scientific clarify itself in
science, one of course began again with the usual pieces of reasoning about not going
too far, then tried again to cover oneself behind the consideration of faith and to
make some concessions to both sides, so that theology might after a fashion preserve
a theological coloring.

One can hardly blame theology for not feeling any particular urge to kill itself. If
theology cannot subsist except in halfness, inconsistency, and want of principle:
what wonder, then, that considered theologians regard consistency as an exaggeration
and the unconditional principle as a one-sidedness.

Thus it came about that many respected and pious theologians (among them Storr in
Tübingen, Seiler in Erlangen, Stäudlin in Göttingen, Stein and others) in a manner as
cautious as it was gentle endeavored to call faith's religious presupposition back to
life, and with great circumspection began to mediate between ``the historical and the
religious interest.''

The first scientific attempt to restore the alliance between hermeneutics and the
doctrine of faith is however really due to F.\ Lücke (now professor in Göttingen). As
a Privatdozent in Berlin, Lücke published a --- as we grant --- ``in many directions
stimulating youthful work (Grundriß der neutestamentlichen Hermeneutik u.\ ihrer
Geschichte, 1817)''; a closer development of the concept of theological hermeneutics
is found by the same author in ``Studien und Kritiken'' for the year 1830. The task
of biblical hermeneutics is then determined thus: ``\emph{to construe the hermeneutic
principles in such a way that the peculiarly theological moment can be united with
them in a really organic manner, and likewise to form and fix
% --- p. 62 ---
the theological moment in such a way that the general principles of exposition retain
their full validity.}''

What misdirection this otherwise well-meant endeavor labors under cannot well be
sufficiently set forth in this hour, but will show itself gradually as we advance
further in these discourses. Anyone who has any impression of what a scientific
principle of interpretation means need only go through the commentaries that are
chiefly laid out and executed according to the above principle --- I mean the
commentaries of Lücke himself, as well as of Tholuck and Olshausen.

There is in these commentaries a double defect, which affects partly the method,
partly the principle.

As regards the method, so far is it from being the case that the religious element is
here worked together into an ``organic'' unity with the scientific, that the religious
rather has a disturbing effect on the scientific and, conversely, the scientific on
the religious.

With all the prolixity of learning that distinguishes these half-edifying
commentaries, one soon notices that the philological-historical apparatus takes up
such a space that the religious-ethical development does not rightly come through, but
must usually be content with some general propositions and interspersed remarks.
Sometimes (especially in Olshausen) the reverse relation sets in, namely that the
religious explanations and dogmatic excursuses spread out with such fullness that the
philological and the historical suffer under it. The fundamental fault is that the
heterogeneous tasks neutralize one another, that the conception becomes ambiguous, and
the interest doubtful. This shows itself still more clearly when we look to the
principle itself.

In the said commentaries the religious and the scientific tasks generally cross one
another in such a way that the scientific principle continually weakens the religious.
This is seen not only in
% --- p. 63 ---
the treatment of the miracles, where the interpreter must now in so manifold ways come
to embarrassment with his half-scientific explanations (into what embarrassment has
Olshausen not come, for example, with his interpretations of the healings of the sick
in Galilee and his half-modern theory of the demons --- an embarrassment from which
Strauß has known how to draw the critical advantage); but it holds also of the
interpretation of Christ's words, especially of such as are expressly designed to
humble the natural understanding. To recall but a single passage, the hard saying runs
(John 6): ``Except ye eat the flesh of the Son of Man and drink his blood, ye have no
life in you. Whoso eateth my flesh and drinketh my blood hath eternal life, and I will
raise him up at the last day.''

Such passages of Scripture cannot be interpreted scientifically%
\footnote{Cf.\ Evangelietroen og den moderne Bevidsthed pp.\ 432---36.}; it is entirely
impossible. For the scientific view they lose their sting and thereby their force, and
transform themselves into merely figurative turns of speech. Instead of impressing on
himself the religious character of ``the hard saying'' and making faith's task
recognizable, the theological interpreter gets into a tangle of contradictory
explanations; and when he then feels that with all these explanations he can at bottom
explain nothing, he is at last overtaken by the doubt that the Gospel of John may
perhaps not even hold up.

In scientific interpretation the scientific relation holds, and therefore the
interpreter ought here also to be consistent enough to relate himself scientifically to
his Bible. It is this objectively-scientific relation that has been so clearly and
definitely expressed by L.\ I.\ Rückert.

In religious-ethical interpretation the religious relation holds, and for that reason
the interpreter ought here also to be consistent enough
% --- p. 64 ---
to relate himself religiously and ethically to the inspired Scripture. It is this
consistency that I have everywhere had the believer hold fast in ``Evangelietroen og
den moderne Bevidsthed.''

With this the spheres are separated: science and faith have now each its own principle
and each moves about its own center.

The eccentric circles can now for the rest touch one another, engage one another,
indeed illuminate one another mutually, if only one does not forget that the principles
are heterogeneous.

That there is a reciprocal relation between Christian and humane cultivation cannot be
denied; but humane cultivation must not therefore be confused with Christian.

Humane cultivation, which grows up out of its own principle without admixture of the
Christian, thrives, so to speak, by its own means, and reaches its ideal height in
science as science. At this stage we find the learned-scientific interpreter of the
Bible.

Christian cultivation will always presuppose the humane to a certain degree; but with
this presupposition it develops for the rest out of its own principle, and in such a
way that the same human being can stand at a relatively low stage of humane cultivation
and yet at a relatively very high stage of Christian cultivation. At this standpoint we
place the simple but unlearned interpreter of the Bible.

Finally, humane cultivation may have reached a certain scientific stage in a human being
who is so gripped by Christianity that he now reduces this cultivation to a means for a
principled and therefore ``unscientific'' development of the Christian life of faith. At
this standpoint we place the simple but nevertheless learned interpreter of the Bible.

Insofar as the learned interpreter explains Holy Scripture religiously, he must
essentially explain it in the very same way as the believing layman. But since the
learned interpreter of Scripture has learning and scientific cultivation in advance of
% --- p. 65 ---
the layman, this cultivation and learning will enable him to interpret Scripture in such
a way that he becomes, in a human sense, a true scribe, a teacher for the simple who have
faith but lack learning.

But in a divine sense the interpretation of Scripture is a superhuman gift, and precisely
thereby exalted above every human art. In virtue of this its double nature, interpretation
strives after a double perfection, a formal and a real one.

As regards interpretation's formal perfection, he will be the best interpreter who has the
most cultivation; but as regards the substance of faith-knowledge, he will be the best
interpreter whose existence has the greatest inwardness in faith. In the first respect the
scribe can instruct the layman; in the last respect the simple layman can just as well
instruct the scribe as conversely be instructed by him. The humane element in Christianity
is that the one can thus teach the other; but the Christian element, the absolutely divine,
is that the one at bottom cannot teach the other at all.

Where, now, does the theologian mean to go with his religious-scientific interpretation?
Religious-ethical interpretation of Scripture stands partly outside science, partly above
science, but never in science!

If it is now really true that there is not and cannot be any intermediate sphere in which
religious-ethical interpretation could be derived from, determined by, or governed by the
learned-scientific: then this is a new proof that \emph{Christianity is higher than science,
and the faith of the Gospel heterogeneous with theology.}

%% ============================================================
%%  SIXTH LECTURE (pp. 66--80)
%%  NB: the Indhold prints "66--86"; that is wrong. VII opens at p. 81.
%% ============================================================
\chapter*{Sixth Lecture}
\addcontentsline{toc}{chapter}{Sixth Lecture}
\label{ch:forel6}
\markboth{Sixth Lecture}{}

\begin{center}\itshape
The theological interpreter of Scripture appeals to the Bible's inspiration
in the whole, but may on no account acknowledge its inspiration in the
particular. Since there is thus doubt about every particular, there must of
course also be doubt about the whole: for the theologian the Bible's
inspiration is an empty phrase. The evangelical interpreter of Scripture,
who in accordance with the law of existence everywhere sees the whole
concentrated in the particular, conceives every single passage of Scripture
as inspired, under the ecclesial presupposition that the whole Bible is
inspired: the faith of the Gospel is heterogeneous with theology.
\end{center}

% --- p. 66 ---
To interpret the Bible, the Church's Holy Scripture, is in other words to find out,
fix, and present the essential relation between meaning and expression, spirit and
letter (πνεῦμα and γράμμα). Inspiration here shows us its two-sided significance. In
the inspired Scripture content and form are inseparably bound together; everything
must be understood according to the letter: he who slights the letter offends the
spirit. But the servant of the letter, who sees only the letter, offends the spirit
too; for the letter killeth, but the spirit giveth life.

The hermeneutic task of determining the relation between spirit and letter has from
theology's first beginning occupied the theologians through the Church's many
centuries, without one's yet having been able to mediate the opposition
% --- p. 67 ---
that already early gave itself a scientific expression in the dispute between the
Alexandrian and the Antiochene schools.

The allegorizing tendency with all the fantasticalities belonging to it has doubtless
been more and more forced back by ever-advancing scientific rigor; but
grammatical-historical interpretation has nevertheless been just as little able to
fathom the spirit as to maintain the letter.

In a certain sense the grammatical interpreter can doubtless maintain the letter,
insofar, namely, as he can by the method acknowledged above find out the text's
literal meaning; but this literal meaning is then at the same time transformed into a
purely human meaning. ``God himself surely did not, after all, think this finite
thought, speak these finite words, construct these grammatical sentences!'' But what
then becomes of the spirit, together with the deeper divine meaning?

This religious question cannot be answered scientifically. For the scientific view
the spirit in the individual words and sentences becomes a purely human spirit, the
meaning a purely human meaning. The deeper, the superhuman, the divine meaning
(insofar as such a thing is assumed at all) is not to be found in any single place.
The deeper meaning first comes to view when one puts together a multiplicity of
passages of Scripture and subsumes them under the idea's generality --- that is, in
other words, when one looks to the whole. The particular is human, the whole divine!
So it is quite characteristic that theology in recent times has so exceeding a
propensity to deify the whole.

If one lays before a theological interpreter of the Bible a definite passage of Holy
Scripture, and asks him whether he now also dares to acknowledge this passage as an
actually inspired passage, so that these words, these expressions, these sentences
are all penetrated by God's Spirit: then he surely comes into great
% --- p. 68 ---
embarrassment. If, on the other hand, one will dispute with him about the passage's
relation to the whole writing and the whole writing's relation to Scripture's
totality, he shall at once be willing to acknowledge inspiration. Choose the Gospel
of Christ's birth, this single definite Gospel as it is found in Luke; then ask the
theologian whether in the name of science he dares to acknowledge this evangelical
narrative as one inspired by God, as an in every way infallible and trustworthy
narrative: he comes into great embarrassment, and attempts to evade the question; but
if on the other hand anyone will reason with him about Luke in the whole, about the
four Gospels in the whole, about the New Testament, indeed about the whole Bible in
the whole: then the theologian shall be willing to defend both inspiration and the
deeper meaning and the divine testimony, and his defense shall become the more
vigorous the more it is allowed to spread out into the whole.

Only the particular wounds; the whole comes too near to no one. The whole is at a
distance, only the particular is present; the whole is in hovering indefiniteness,
only the particular is definite and decisive; the whole is for the scientific
observer's objective imagination, the particular is for the believer in
appropriation's inwardness.

The opposition between the theological and the evangelical interpretation can
therefore, with respect to the relation between spirit and letter, be expressed thus:

\emph{According to the theological interpretation the particular is posited as a
moment in the whole; according to the evangelical interpretation the whole is posited
as concentrated in the particular.}

To illuminate this opposition more closely I choose a definite example:

\emph{When anyone smiteth thee on thy right cheek, turn to him the other also}
(Matt.\ 5:39).

% --- p. 69 ---
In these words of Christ there is not only unity of spirit and letter, but here there
is at the same time difference. One may readily grant the scientific interpreter that
this commandment cannot be taken straightforwardly according to the letter without
the letter's killing the spirit; but it is then also certain that only
religious-ethical interpretation is able to raise itself to the point where the
spirit again unites itself with the letter, so that the ideal itself is seen in the
literal.

If anyone takes the words straightforwardly according to the letter, and thus
understands them spiritlessly, the misunderstanding will easily show itself; one
surely needs no scientific guidance to see this. A human being can, when he receives
a blow on the right cheek, readily turn the left to it without therefore being a
Christian; just as conversely a human being can be a true Christian without therefore
turning the left cheek every time he receives a blow on the right. That the believer
should now make himself guilty of so spiritless, arbitrary a misunderstanding of the
literal expression is by no means to be feared; on the other hand it is in a high
degree to be feared that the scientific interpreter, who here comes so easily to a
rational separation of spirit and letter, will explain the commandment away for us so
spiritually that he entirely effaces Christianity's mark.

With such an explaining-away it goes as of itself. Once one has found out that
Christ's commandment cannot be understood literally, what lies nearer than to draw a
general morality out of the figurative expression, and to elucidate this general
morality by a juxtaposition of several passages of Scripture belonging here. Let the
general morality now be that a disciple of Christ shall know how to govern himself,
show himself conciliatory and yielding, be willing to be reconciled with his
adversary, overcome evil with good, and so on. As the scientific interpreter now
expounds this general morality in all generality according to ``the analogy of
Scripture,'' it occurs to him that other
% --- p. 70 ---
moralists and wise men of antiquity have put forward similar moral precepts in
similar figurative expressions; such classical passages, be they from Confucius or
Pythagoras or Xenophon or Plato or Cicero or Horace or Seneca, or whoever it may be,
ought then also to be cited. By citing such passages the interpreter not only shows
his great reading and his learned survey, but the interpretation at the same time
gains in scientific rigor: Christ takes his place in the ranks of the noblest spirits
who have lived on earth, and the noblest spirits take their place in turn in the
ranks with Christ; so the analogy of faith is widened, and Scripture's divinity is
recognized in the whole.

Otherwise, indeed quite otherwise, does it stand with religious-ethical
interpretation. The religious man, who in faith's inwardness strives to understand
Christ's commandment, does not begin by scattering his attention over a multiplicity
of passages of Scripture, does not give himself over to leafing through the Bible in
order to entertain himself and others with an interesting comparison of
``unfigurative'' and ``figurative,'' ``prosaic'' and ``poetic'' statements; still less
does he lose himself in learned investigations in order to discover what other sages
may possibly have said about the same thing; no, the religious man collects his mind
and fixes his gaze on the one commandment that here aims precisely at him as a
commandment of Christ; the religious man loses himself in the appropriation of the
commandment's particularity. When he then reads the commandment and understands that
this commandment concerns him himself in particular (as it concerns every disciple of
Christ in particular); when he grasps and understands that as Christ's commandment it
must be for him an unconditional commandment, a commandment according to which he is
to live: then the figurative expression will at once teach him that the commandment
has all the Christian marks.

\emph{``When anyone smiteth thee on thy right cheek'':} --- yes, what is more natural
than that thou, the ill-treated one, shouldst seek to avenge it on thine adversary;
what is more natural than that thou, the wronged one, shouldst seek to obtain thy
right, indeed shouldst involuntarily be tempted
% --- p. 71 ---
to take the law into thine own hands? --- Never is self-love more dangerous than in
the moment when it makes itself one with the sense of justice and incites to
requital. As self-love now flares up, Christ's commandment comes and would strike it
down. The first natural impression that the natural man receives when the commandment
so instantly and unconditionally will annihilate the ennobled self-love, justified in
the moment of anger, is --- the impression of offense.

That one is in general and to a certain degree to deny oneself, that one is in
general and to a certain degree to be conciliatory, yielding, and willing to be
reconciled with one's adversary, that everyone can understand who has any moral
sense: the Israelite can understand it, the noble pagan can understand it! but that
such self-denial as Christ's commandment here enjoins should really be meant
literally and in earnest --- no, that only a believing Christian can understand, when
in unconditional obedience he attempts to act accordingly.

The pagan ennobles self-love, but does not deny it (even in Moses it is said: eye for
eye, tooth for tooth); only in Christianity is self-love judged, and self-denial
posited as the inviolable condition for entering into the kingdom of heaven. This
unconditional self-denial cannot be acquired by a merely human resignation. A pagan
sage, a Stoic, will by continued exercise be able to bring it to an astonishing
degree of self-mastery; will perhaps be able to bring it to the point where he never
acts rashly, is never taken unawares by an outburst of passion; will perhaps win such
% The Danish here carries a raised speck (a broken comma?) after „Magt“; no
% punctuation is called for and none is transcribed there or here.
power over himself that he calmly receives the blow on the right cheek without
changing a feature, without even finding it worth the trouble to strike back; but
here too is the limit of resignation. If one goes beyond this limit, there shows
itself the defiant arrogance, the ice-cold contempt that even condescends to play
with injustice; as if the Stoic would not merely endure the blow on the right
% --- p. 72 ---
cheek, but, as proof that he did not regard it, would even as it were mock the
mockery, and turn the left cheek to.

Such defiance and contempt, in which self-love secretly triumphs --- such defiance and
contempt is nevertheless so vastly different from Christian self-denial that a human
being had better, in all frailty, requite the unjust blow and afterward repent his
rashness than thus deify his own pride and harden himself in an inhuman arrogance.

When, therefore, Christ's commandment presupposes that someone could inflict an
injurious wrong on thee, and smite thee on thy right cheek, then the commandment also
presupposes that such a blow pains thee, that thou as a human being in flesh and blood
bitterly feelest the injury and the mockery, indeed it presupposes at the same time
that it even stood in thy power at that same moment to requite the injustice by
requiting the blow. When, then, the commandment nevertheless not only forbids thee to
strike back when anyone smites thee on thy right cheek, but even expressly commands
thee to turn the left cheek to thine adversary: then this commandment is to natural
reason an offensive commandment, a commandment of unnatural and inhuman severity.

For the believer, on the other hand, this commandment sounds as a divine commandment
of love, a commandment that cannot indeed be fulfilled by the natural man, but is
nevertheless to be fulfilled, because its fulfillment not only requires of the human
being an unconditional obedience, but at the same time points to a superhuman
assistance. Only Christ, who promises such assistance, can in earnest put forward such
a commandment; but when the commandment is in spirit and truth conceived as Christ's
commandment, the ideal too will be recognized in the literal expression. If Christ
would inculcate self-denial in a definite example, where could an example be found
more telling than this? How decisive, how unconditional must self-denial not be, when
he who receives a blow on the right cheek is at once and as if without deliberation to
turn the left cheek to! And what
% --- p. 73 ---
force of inwardness must that human being not possess who is to fulfill this
commandment! for everyone who practices self-denial in such a way that the world
understands that it is self-denial, and that the self-denial is reasonable,
nevertheless gets the reward that the world acknowledges him for a virtuous man (for so
fair-minded is the self-loving world that it is glad to praise self-denial, if only
someone else will practice it); but he who without resistance or complaint turns the
left cheek to when he has received the blow on the right --- what does the world say of
him? Of him the world does not say that he is a reasonable and virtuous man; but it
says of him that he must either be an intolerable, defiant
% The Danish here sets „Mennneske“ with three n's --- printer's error, kept as
% printed there. The English has the correct word.
human being who despises everything, or else a simple fool.

Thus this commandment has all the Christian marks by which it can literally be known
that in spirit and truth it must be Christ's commandment. The whole is here concentrated
in the particular; the whole of Christian ethics with all the works of love is compressed
into the one commandment: \emph{when anyone smiteth thee on thy right cheek, turn to him
the other also.} If anyone will attempt to order his life according to this commandment,
then the Gospel promises that God's Spirit will stand by him, and existence help him in
principle to discover the whole of Christianity; but he who will not and cannot live
accordingly can neither grasp what lies in the commandment; he understands it only in
imagination, in possibility (as we at this moment doubtless understand it in imagination
and possibility); but in actuality he does not understand it.

That the religious-ethical interpretation of the passage of Scripture here treated is
entirely heterogeneous with the scientific can now easily be proved.

Insofar as the author relates himself scientifically to his text, and refers the
significance of each individual passage of Scripture to the general categories of the
totality, he cannot even in representation get so far as to touch existence; for
% --- p. 74 ---
existential thinking distinguishes itself from the scientific precisely in this, that at
every point it touches particularity, the individual human being, the individual
commandment, the individual relations and conditions under which the commandment finds
its application, and so on. But even if one would reckon among the scientific that
interpretation which converts the essential content into nothing but existential
determinations, and assert that the sample we have here given of a religious-ethical
exegesis could at bottom just as well be called scientific as the
philological-historical: then this is a pure misunderstanding.

It is indeed true that religious-ethical interpretation, insofar as it is communicated
in a lecturing discourse, has a certain likeness to the scientific --- namely this, that
it is lectured. But insofar as religious-ethical interpretation is lectured, it points at
every point beyond itself, declares itself insufficient, and abolishes itself. The
existence with which the lecturing discourse deals is namely no actual, but only an
imagined existence. Therefore one cannot conclude either that someone is really religious
because he has a certain facility in putting forward religious-ethical explanations. In
the form of a discourse the religious-ethical explanation is therefore still only an
imagined explanation, an explanation that differs from the scientific merely in this,
that it turns toward existence, requires existence, and makes the existential tasks
recognizable.

Religious-ethical interpretation has its truth solely in actual existence; only in actual
existence is it in earnest heterogeneous with the scientific.

In actual existence the believer who interprets his Bible is no imagined figure, but an
actual human being, and the understanding of Scripture's word no imagined but an actual
understanding. Actual understanding is only in actual application; the circumstances of
life and the cases in which Christ's commandments find their
% --- p. 75 ---
application are then not imagined either, but actual; but the actual case is each time an
individual, definite case.

If we now hold fast to this point of view, it follows that religious-ethical
interpretation is not only different from the scientific, but at the same time so
heterogeneous with it that it is even, in the strictest sense, impossible for the one
human being to interpret the Bible for the other, since it is impossible for the one
human being to put himself entirely into the other human being's existence. How far the
believer in actuality understands and appropriates Christ's commandment: ``when anyone
smiteth thee on thy right cheek, turn to him the \emph{other also},'' depends solely and
alone on how he lives, and this everyone must know within himself: the relation to God is
a relation of conscience.

The self-denial that the commandment enjoins is unconditional; but every human being's
life is subject to finitude's conditions. It is not granted to the human being to carry
the unconditional through in every moment, and yet it demands itself of him in every
moment. The unconditional can besides work in manifold ways in the different conditions.
A human being can in countless ways exercise himself in religious self-denial, can in
countless ways display this self-denial, sometimes even in such a way that the most
inward self-denial shows itself outwardly as presumptuous self-will. Under such
conditions it is entirely impossible for the one human being to advise the other:
everyone must in faith advise himself. When this turning point sets in (and it sets in
every time existence presses in with real earnestness), all mutual instruction falls
silent; everyone grasps in his own bosom --- after the understanding, everyone encloses
the commandment within his own inner being, and shuts himself in, so to speak, together
with the commandment, in order on his own responsibility to understand how this
commandment in this given case can and ought to be applied.

As the earnestness of existence thus halts mutual
% --- p. 76 ---
instruction, admonition, and guidance, and separates human beings, so that everyone in
the hour of decision must advise himself, the meaning can nevertheless not possibly be
that obstinate, arbitrary self-will should draw new advantage from this. On the contrary,
so far is it from being the case that the individual is here referred to arbitrary
self-will, that he now first fully understands how greatly he needs counsel, how
impossible it is for him to advise himself. To natural self-feeling the human being
cannot entrust himself, for self-feeling is in collusion with self-love; but here it is a
question of self-denial. Reflecting understanding cannot solve the task for him either;
for partly the understanding is bribed by self-love, partly it becomes entangled and
caught in its own reflection. Instead of solving the task, reflection will only make it
far harder still, by making it dialectical. On the one side an absolute self-denial is
required, an unconditional self-sacrifice; on the other side it is required that the human
being shall recognize his limit. A human being must after all always yield something to
the outward circumstances at hand; one cannot, as it were, burst all existence's
conditions at a single wrench in order to grasp the unconditional. For reflection the
commandment comes apart, and divides into two mutually contradictory commandments: grasp
the unconditional! and: recognize thy limit! This contradiction cannot be mediated by
continued reflection: the longer one reflects, the greater the difficulty becomes;
besides, in the moment of decision there is no time to reflect, for the moment is precious:
what matters is to resolve and to act.

If the human being is now thus referred to himself without nevertheless being able to
advise himself, then there would be no counsel for him at all if God's Spirit did not bear
witness with his spirit in conscience, and prompt him as to what he should do in the
decisive moment. So says the Gospel: the Spirit that teaches the human
% --- p. 77 ---
being the right choice in the hour of decision; the Spirit that determines self-denial in
relation both to the outward conditions and to the abilities and cultivation of the human
being concerned; the counseling and governing Spirit that lifts the human being above
himself in the moment of self-denial; the Spirit of power that makes the wronged one
become stronger than his adversary in that he --- wondrously enough! --- becomes stronger
than himself; the Spirit that makes it possible for him who receives the blow on the right
cheek to turn the left cheek to in the superiority of meekness and humility: could this
Spirit well be any other than God's own Holy Spirit?

What in science is without significance and in theology an empty-solemn turn of phrase
--- \emph{namely, that the Holy Spirit is the only right interpreter of Holy Scripture}
--- becomes an actual truth in the believer's existence; for the Holy Spirit is a Spirit
of existence.

In actual existence the Spirit will always answer to faith, as faith to the Spirit. Dare
to believe, and thou shalt receive the Spirit's power to grasp the unconditional; dare to
deny thyself, and thou shalt not only receive power to break with the considerations of
worldly prudence, but also patience and wisdom to understand the conditions of finitude
and frailty, for the Spirit of power is at the same time the Spirit of counsel and wisdom.
This is the law, the revealed law of God, by which every human existence has to order
itself: that he who has a strong faith receives a strong spirit, and understands
Scripture's word in spirit and truth; but he who has only a weak faith receives a weak
spirit, and understands the Word as through a shadow. This is the law, the inviolable law
of the interpretation of Scripture: everyone who hears the Word and acts accordingly
experiences the Word's power and is to be likened to a wise man who builds his house upon
the rock; but if anyone hears the Word and does not act accordingly, then he can never
recognize the Word's
% --- p. 78 ---
power either, but deceives himself by an imagined understanding and is to be likened to
the foolish man who builds his house upon the sand.

The theologians' opinion that the Holy Spirit will surely engage with their learned
researches and at length come to their aid with a true scientific interpretation of
Scripture is grounded neither in science's nor in revelation's essence, but is, plainly
said, neither more nor less than a theoretical confusion of heterogeneous tasks. This
confusion expresses itself every time one looks at the particular; for the theologian, as
theologian, never has the Holy Spirit in the particular, but always in the whole, never in
actuality, but always in possibility, never in existence, but always in imagination.

But the Spirit that is thus in the whole without coming to view in the particular is a
nebulous and hovering Spirit, a Spirit that assuredly offends no human being, but
assuredly saves no human being either, a Spirit that founds turns of phrase as a poet
founds imaginings.

Theological science's scientific defenders will without doubt find this assertion somewhat
personal; to this I answer that the assertion is personal only because it is true, and I
shall prove to superfluity what I say.

It is well known that not only the orthodox theologian, but also the rationalist, not only
the speculative but also the unspeculative rationalist appeals, in interpreting Holy
Scripture, to
% In the Danish the „e“ of „hellige“ is blotted with ink in this copy; the
% reading is nevertheless certain from the sense and the surrounding letters.
the Spirit's testimony in the whole (\textit{testimonium spiritus sancti theoreticum}).
From this one can now draw two conclusions: either the Holy Spirit, which thus gives its
testimony in all these interpretations, is itself both orthodox, heterodox, speculative,
unspeculative, and so on; or else the Holy Spirit's testimony must be exalted above these
finite schools' finite oppositions. But whichever of these
% --- p. 79 ---
assumptions one prefers, the prevailing theology must perish. In the first case the
theologians' Holy Spirit is a self-contradictory and self-conflicting Spirit; but a
science whose spirit is in essential conflict with itself cannot possibly subsist. In the
last case the Spirit to whose testimony the theologians appeal is a supra-scientific
Spirit, a Spirit in whose essence theology itself therefore has no share, since theology
is and will be science.

As regards religious-ethical interpretation of the Bible, it is by no means ``the
theologians'' that we need; but insofar as a human guidance in this direction is possible
at all, an existential author is doubtless the guide we could greatly need, if only our
theoretical conceits did not prevent us from understanding the guidance.

By an existential author I mean such a one as not only has a quite peculiar gift for
exhibiting existence's categories and the corresponding tasks in all possible relations,
but an author who at the same time has the self-denial to be willing to sacrifice a life's
powers on working out these tasks; an author who in the earnestness of jest is not even
afraid to work against himself, in order the more honestly to serve his idea; an author
who, when it must be so, would rather bear the misinterpretations of the day than give up
the understanding with an eternal truth; a poet in the immediate enamorment with which he
looks upon the ideal; an ethicist in the strict understanding that it is precisely his own
existence that the ideal will require; a religious individuality in the humbling
recognition that he nevertheless cannot --- what every human being ought to be able to do!

Thus, roughly thus, do I picture to myself the authorial existence that would surely be
most nearly equipped to wage war against the present age's ``delusions of sense,'' and to
make Christianity recognizable to modern reflection.

% --- p. 80 ---
That the theologians --- I mean the genuinely scientific ones, especially the genuinely
speculative-scientific ones --- would of course prefer to overlook such an authorial
existence, or, if the phenomenon in literature were now so striking that it could no
longer, at least for conversation's sake, be overlooked, would then interpret it as --- an
``unscientific'' one-sidedness! --- a ``remarkable'' one-sidedness! --- a ``brilliant''
one-sidedness! --- an ``ecstatic'' one-sidedness! is after all easy enough to understand
--- in conversation. This many-sided interpretation proves precisely that the theologians
must have their own many-sided hermeneutics.

Let everyone then interpret it according to his hermeneutics as he will: I interpret it
thus, \emph{that Christianity is higher than science, and the faith of the Gospel
heterogeneous with theology.}

%% ============================================================
%%  SEVENTH LECTURE (pp. 81--96) --- contains the p. 82/83 scan seam
%% ============================================================
\chapter*{Seventh Lecture}
\addcontentsline{toc}{chapter}{Seventh Lecture}
\label{ch:forel7}
\markboth{Seventh Lecture}{}

\begin{center}\itshape
In the treatment of sacred history --- ``the Life of Jesus'' --- the
theologian will work against the freethinker; but the scientific criticism
on which he proceeds shows plainly enough that in principle he is at one
with the freethinker. The believer, who in his conception of the Gospels'
Christ holds exclusively to religious-ethical criticism, is in principle at
one with the evangelists: the faith of the Gospel is heterogeneous with
theology.
\end{center}

% --- p. 81 ---
As is well known, Germany's theologians are at present eagerly occupied with their
critical-scientific treatments of the life of Jesus (Das Leben Jesu). ``The critical
introduction'' has to that end called up theology's various powers --- hermeneutics and
dogmatics and ethics and apologetics, together with polemics and irenics --- to work with
united forces, so that one might thus with united forces work toward the common goal, the
solution of the great task. The plan has accordingly been set in motion; but, it cannot
be denied, the undertaking has hitherto had quite a peculiar fate.

The modern idea of a ``life of Jesus,'' which was conceived and born of a famous
theologian, had namely the remarkable fate that immediately after its birth it was adopted
by an acute philosopher, so that it might be all the better reared and formed in the
spirit of modern science itself. It was the critic David Strauß who fetched the idea from
Schleiermacher's lectern, sent it
% --- p. 82 ---
to school with Hegel, and thereafter let it appear in his celebrated work ``Das Leben
Jesu'' --- to the terror and offense of theologians of every color. Never has theology
suffered a more complete defeat, and never has negative criticism won a more complete
victory, than precisely in this scientific collision. Many theological writings have gone
forth, many theological plans of attack been attempted, much theological ink spilled for
the scientific refutation of the great freethinker: --- all in vain!

Doubtless Strauß has in several respects been corrected by his many opponents; among
others the theologian J.\ H.\ A.\ Ebrard has in his ``Wissenschaftliche Kritik der
evangelischen Geschichte'' (part 2) successfully pointed out a number of arbitrarinesses
in ``the destructive method''; but all this is of no help so long as the foundation on
which the Straussian criticism rests has not itself been destroyed. The point is that
Strauß does after all have a principle, whereas theology is and remains without principle.

The Straussian principle is namely nothing other than the modern consciousness's own
generally acknowledged fundamental principle: the principle of positing the idea of reason
as the absolute, deifying the laws of nature, and raising science to religion. However
many small alterations religious feeling and subjective humanity may otherwise make in
this fundamental principle, it is nevertheless seen through all modifications and
interpretations to be the very same.

The Straussian endeavor to raise science to religion cannot possibly be heterogeneous with
the theological endeavor to raise religion to science; for both endeavors proceed from the
presupposition that religion and science are homogeneous. But when the theologians in so
essential a sense stand at the very same standpoint as Strauß, how then should the
theologians be able to refute Strauß?

We see it plainly in Ebrard's relation to Strauß. Strauß proceeds from the presupposition
that every belief in miracles is super%
% --- p. 83 ---
stition, and then proves from this that all the miracle stories are myths, in the New
Testament just as much as in the Old; Ebrard proceeds from the presupposition that God can
work miracles, and then proves from this that the miracle narratives in the four Gospels
are trustworthy; Strauß is a naturalist, Ebrard a supranaturalist; Strauß is unbelieving,
Ebrard believing; Strauß tears down, Ebrard builds up: to that extent Ebrard doubtless
stands at quite another standpoint than Strauß. But if we look more closely, it will soon
show itself that the opposition is not after all so great as it seems at first glance, but
that it can, as mediation says, quite well be abolished in a higher unity.

Despite all the disagreement that otherwise prevails between the two opponents, they are
nevertheless both agreed on a very important and essential point; they are namely both
agreed that in treating ``das Leben Jesu'' one must of course stand at science's
standpoint. This presupposition the theologian shares with the freethinker, and with this
presupposition the freethinker crushes the theologian.

A supranaturalist at science's standpoint is a contradiction; a scientific critic at the
standpoint of belief in miracles is a contradiction; if Christianity is to be explained in
science, then it is already explained, then it is already settled what it is: it is a
mythology; if sacred history is to be criticized by science, then it is already
criticized, it is already settled what it is: it is a chain of myths, one great myth.

``The divine'' --- says Strauß --- ``cannot have happened in this way; what happened in
this way cannot have been divine --- thus will the difference express itself, and when
interpretation strives to settle it, it will strive either to present the divine as not
having happened in this way, thus denying the old sources their historical validity, or to
show that what happened
% --- p. 84 ---
is not divine in this way, thus explaining away the absolute content of these books. In
both cases the explanation can go to work partially or impartially; partially, when it is
blind to the consciousness of the difference between the new cultivation and the old
sources, and only imagines itself to be finding out the latter's original meaning;
impartially, when it clearly recognizes and openly confesses that it sees what those old
writers relate otherwise than they themselves have regarded it.''

``The ancient, especially Oriental, world'' --- it is said further --- ``was, in
consequence of its preponderantly religious tendency and its slight knowledge of the laws
of nature, of the connection in earthly, finite being, something so loose that it was able
to leap from any point in it over into the infinite, and at every single alteration in
nature and the human world to regard God as the immediate cause. From this standpoint of
consciousness biblical history too is written . . . . But the more recent age has a series
of the most laborious researches, continued through centuries, to thank for the insight
that everything in the world hangs together by a chain of causes and effects that tolerates
no interruption . . . . This conviction has so become the consciousness of the more recent
world that the opinion or assertion that a supernatural cause, a divine influence, has
somewhere intervened immediately is in actual life regarded straightforwardly either as
ignorance or as deception; and it has advanced to an extreme in the more recent
Enlightenment's view, which, directly opposed to the biblical one, either removes divine
causality entirely, or at any rate pushes it so far back that only in the act of creation
would it be an immediate one --- that is to say, that God would work upon the world only
insofar as at the creation he gave it a definite arrangement. From this standpoint, which
sees in nature and history a fixed web of finite causes
% --- p. 85 ---
and effects, the biblical narratives, in which this web is at countless places torn apart
by an intervention of divine causality, cannot possibly show themselves as history.''

One sees from this how insignificant it is to want to refute Strauß in the particular and
to try one's hand at scientific counter-proofs against his destructive criticism. If in
profane history it is already far easier for criticism to tear down than to build up, how
meaningless then does the scientific attempt become of trying by critical proofs of
probability to maintain the miracles in sacred history. Ebrard's theological venture has
accordingly not been able to win the theologians' approval either. Let one read, for
example, Bleek's review of Ebrard's writings, and one will get a very disheartening
impression. The as acute as he is profoundly learned Bleek goes through Ebrard's
assertions, and winces at every step, like one walking on needles. ``Hardly credible! ---
Improbable! --- Unprovable! --- Highly improbable!'' are the exclamations that escape the
otherwise so well-disposed reviewer at every moment. Meanwhile the reviewer Bleek is
caught in exactly the same misunderstanding as the author Ebrard. The misunderstanding is
that one will judge the miracles scientifically, and apply the measure of probability to
that which by its nature is and must be the improbable.

The dispute with Strauß transforms itself on the instant into a dispute between religion
and science. It is the critic's merit to have given the matter so strong a push that it has
precisely thereby been driven to this decisive turning point. The dispute of the faith of
the Gospel with the Straussian criticism is a dispute of competence, a dispute of
sovereignty, a dispute about the supreme power, the highest authority in the world of
spirits.

If science is the highest instance, then Strauß is in principle right, then it is,
frankly spoken and plainly said, all over with
% --- p. 86 ---
Christianity, then religion lies at its last extremity and awaits only a gentle dissolution
--- in humane culture.

If, on the other hand, Christianity is higher than science, if it has its own supernatural
principle and moves in a sphere into whose depths no human knowledge penetrates, then it
will also still, as before, rule with a sovereign's authority in its own kingdom, and make
every human being answerable for how he understands its authorization and its relation to
science.

Science, as science, cannot possibly arrogate to itself the supreme authority. Insofar as
science knows what it knows, it knows at the same time that there is something it does not
yet know, and something it will never come to know in this existence. Science knows within
itself that it exercises a magnificent and beneficent influence on humanity in leading the
generations forward from enlightenment to enlightenment; but it also knows that it has, and
can have, only a very mediate influence on actual human beings' actual existence. Science
knows within itself that the great events in world history --- the oppression of peoples,
the victories of revolutions, the wars of nations --- are not settled by pure grounds of
reason. Science knows within itself that there are powers, dark, terrifying powers in
existence that never bow their necks under its dominion; science knows that the new
thoughts, the new ideas, the new means that it offers the human race at every step must
serve the evil as well as the good, the unjust as well as the just. Science knows within
itself that the knowledge of truth that it bears and upholds in continually advancing
development never reaches and never can reach so far as to answer for us the highest
questions, the most important of all for human existence: is there a God? is there an
eternal life?

By no means do I assert that science, as science, straightforwardly leads to a denial of
the ground of things or to a denial of
% --- p. 87 ---
eternal life. Far from it! Philosophy cannot live without the idea: when it misconstrues
the idea, it misconstrues itself; when it dissolves and denies the idea, it dissolves and
denies itself. For philosophy the idea is as God, God as idea: \emph{the idea is divine};
for philosophy the eternal is as idea, the idea as the eternal: \emph{the idea's life is an
eternal life}. Higher than to this explanation philosophy has hitherto not come, and can
never come either.

Philosophy's idea can now in manifold ways be developed, articulated, and formed in the
dialectical formations of conceptual systems: otherwise by the one philosopher than by the
other, otherwise by an Aristotle than by a Plato, otherwise by a Schelling than by a Hegel.
Under these systematic formations idea-knowledge doubtless makes essential progress,
insofar as the idea itself clarifies itself along with the clarified concept. But the
idea's fundamental essence is and nevertheless remains through all the systems the one and
the very same fundamental essence, just as philosophy too in all philosophies is and
remains the one and the same philosophy. Plato's and Aristotle's God is also Schelling's
and Hegel's God, that is: the absolute idea.

The great advances that philosophy in recent times is supposed to have made with
Christianity's assistance are illusory. What philosophy has borrowed from Christianity it
must give back again; this loan can never become its property: Christianity is
heterogeneous with philosophy. Study the Hegelian philosophy, go through this great,
rigorous, imposing conceptual system point by point from the beginning to the end; then
put to yourself the question what philosophy has here at last fathomed about God's
existence and the soul's immortality; hold fast to the concept, let yourself not be
beguiled by the objective imaginative luster of the idea: you shall assuredly see that the
Hegelian God with all his logical trinitarian movements is still the old
% --- p. 88 ---
% The Greek is set in antiqua italic against the surrounding Fraktur.
Aristotelian world-God (\textit{τὸ πρῶτον κινοῦν, οὐ κινούμενον}), that first mover who is
not moved, i.e. the absolute idea; indeed you shall see that Hegel, still in the nineteenth
century after Christ, has not come a hair's breadth further with the immortality of souls
than Plato already was in the fourth century before Christ.

The foundation on which philosophy erects its systems is the idea's foundation; but this
foundation is ambiguous, like the idea itself. What, pray, is the philosophical idea? An
amphibolous being! a being that does not itself know whence it comes and whither it will;
a being that now points out beyond itself upward toward the personal God and the kingdom
beyond, but now again denies the beyond, mocks the personal God, and declares itself the
divinity of the actual world.

The ambiguity rests on the fact that the idea looks otherwise to so-called pure thinking in
relation to essence and possibility than it does to immediate experience in relation to
existence and actuality.

Metaphysical thinking does not come to rest until it comes to rest in God as the wisdom
that thinks itself and all things, wills itself and all things, explains itself and all
things, of whom, through whom, and unto whom are all things: philosophy to that extent
finds it both thinkable and possible that the absolute Spirit is a Spirit free and
self-conscious in itself. Likewise the philosophical development of the idea can be laid
out in so subjective a direction that the human I recognizes therein its inner infinity,
its original power to conquer the barriers of existence, indeed of death, and to raise
itself by the moral world-order to eternal freedom in ever-increasing perfection.

But what philosophy thus promises us in relation to essence and possibility, it recalls
again when the idea is considered and tested in relation to existence and actuality. The
philosophical idea cannot free itself from matter's barriers,
% --- p. 89 ---
and therefore cannot raise itself above this finite world either. The idea is in the
existing things, as the things are in the idea. The idea has no independent existence: its
reason is nature's reason, its powers the world's powers; it ferments in the ground of the
elements, it awakens to life in the animal, to light and freedom in the human being.
Individuals must die, the human being as well as the animal; only the idea lives in all
generations: the idea splits apart in actual existence; for actual existence there is
therefore no personal God and no personal immortality.

From this it must surely be possible to discern how the religious man who will support his
faith on philosophy and science is deceived. A personal God, such as the one religion
proclaims, is to science an offense; a personal immortality, such as the one Christianity
promises, is to philosophy an offense! Science holds, as regards actuality, to the present,
the natural; but religion points to the beyond, the supernatural.

The dispute between the believer and the freethinker can then be more closely determined
thus: the freethinker is right in the assertion that religion and science do not suit one
another; but on the other hand he is greatly wrong when he draws from this the conclusion
that science can abolish religion. The possibility of a revelation science cannot deny;
the actuality of a revelation science cannot confirm; the truth of a revelation science
cannot test.

It has been supposed that such a separation of religion and science as the one I here set
forth would occasion great confusion, since one would then have to believe blindly, and
might just as well take it into one's head to believe the false revelation as the true.
Nothing but evasions! Let one only be so good as to reflect a little seriously on the
difference between subjective knowledge and objective knowing.

% --- p. 90 ---
It has been supposed that he who accepted a revelation without first testing its objective
truth in objective knowing neglected a duty toward truth, for which ethics ought to call him
to account. A poor argument, seeing that ethical truth cannot even itself be tested in
objective knowing.

To want to convince oneself beforehand of a scientific truth --- for example a mathematical
or a logical one --- before one accepts and applies it, is in order; but to want to be
convinced beforehand of an ethical truth before one accepts and applies it is nonsense,
since ethical truth refers to resolution and action. Pure scientific truth can be
established by pure scientific thinking: I need do nothing whatever but think in order to
get the truth proved. In ethical truth, on the other hand, the deed is itself a part of the
proof, indeed the proof itself; but when the deed is itself a part of the proof, indeed the
proof itself, then it is surely a pure folly to want to have the complete proof in advance
of the deed.

Even if in deep scientific earnestness one has long since forgotten the man who would not go
into the water before he had learned to swim, one ought nevertheless not to sink so deep
into the depths of theological knowing that one entirely forgets what the warrior sings in
``Wallensteins Lager'':

\begin{quote}
Und setzet Ihr nicht das Leben ein,\\
Nie wird euch das Leben gewonnen seyn.
\end{quote}

If Christianity is higher than science, then it surely follows of itself that evangelical
history must also be higher than scientific criticism. If Christianity is only for the
believer, then sacred history too is only for the believer; from which it follows in turn
that the freethinker can as little refute the believer as the believer can refute the
freethinker.

In imagination and possibility they may well understand one another, enter into one
another's train of thought, and to that extent negotiate and
% --- p. 91 ---
dispute with one another; but in actuality the one cannot by any objective proof overcome
the other or draw the other over to himself; each of them has his own existential footing:
it is hard existence that sets them apart. The freethinker exists in unbelief; his
conception of the life of Jesus is entirely consistent under the presupposition that
science is higher than religion; the religious man exists in faith; his conception of the
life of Jesus is entirely consistent under the presupposition that religion is higher than
science. Between the two now stands the theologian --- or rather, he does not stand at all;
he at least does not stand firm (for he has no existential footing), but sways and limps to
both sides, now toward faith and now toward science. The theologian's conception of the
life of Jesus is a conception without principle, in nothing but halfness and inconsistency;
this halfness and inconsistency I call mediocrity.

A presentation of the life of Jesus can now be laid out in such a way that the freethinker
with his ``riper consciousness,'' the religious man with his ``faith of the Gospel,'' and
the theologian with his ``mediocrity'' join in going through the evangelical history and
each expressing himself in his own spirit on each of the Gospels. Such a presentation is
then to be characterized as an existential presentation, insofar, namely, as it aims not at
a scientific juxtaposition of different theories, but at a living, passionate strife between
different characters.

The existential presentation of the life of Jesus must in the strictest sense of the word be
critical; but the criticism that here comes into application is, be it noted, not a
scientific but an existential criticism. Criticism comes, after all, from crisis, and surely
no one will deny that existence has its crises. The spirit that sets the strongest
oppositions in motion is then also the highest spirit, a spirit that effects the deepest
crises. When, then, a human
% --- p. 92 ---
being comes forward and proclaims himself God's only-begotten Son: what crises must such a
one not call forth!

Thus Christ comes forward: a human being in outward lowliness, and yet a human being with
divine authority; a human being who makes the seeing blind, and yet a human being who makes
the blind see; a human being who in this world has not that whereon he may lay his head, and
yet a human being who says to the dead: arise! and to the palsied: thy sins are forgiven
thee! Who, then, is this human being? --- The revealed God? The manifest deceiver? Here the
crisis begins.

Evangelical history bears existential criticism in itself: it shows us Christ in relation
both to the non-believing and to the believing.\footnote{Cf.\ Evangeltr.\ og den mod.\
Bevidsthed, pp.\ 210---227.} The non-believing see him at once as a deceiver; the believing
see him as God's Son; the great crowd changes its mind. These crises and this criticism
accompany Christ uninterruptedly through his whole public life, from Galilee to Judea, from
Jerusalem to Nazareth, from Capernaum to Golgotha. It was men of Galilee who became his
first disciples and believed him on his word, that from now on they should see heaven open
and God's angels ascending and descending upon the Son of Man; it was men of Galilee who in
Nazareth's synagogue were offended at him, and thrust him out of the city to cast him down
from the brow of the hill. It was Pharisees and scribes who hated him as a fanatic and a
seducer of the people; it was a Pharisee and scribe who came to him in the stillness of
night as to a teacher of Israel. It was a robber on the cross who reviled him, and it was a
robber on the cross who confessed him. It was a Roman who condemned him in Jerusalem when
the Jews raged: Crucify, crucify! and it was a Roman who acquitted him on Golgotha, when the
graves burst and the earth quaked.

% --- p. 93 ---
The relation in which Christ by his first appearance places himself to his contemporaries is
at the same time determined as the fundamental relation of principle in which he places
himself to the whole human race through all times. On the one side there follow him the
offended, who now would stone him, now thrust him down, and always mock him, because his
word is to them a hard saying which no one can hear. On the other side there follow him the
believing, who confess his name and cannot leave him, because in the confusion of the times
they know no other way out than Peter, when he said: \emph{Lord, to whom shall we go? Thou
hast the words of eternal life.} About these characters moves the crowd, the surging crowd
that comes and goes; the dependent crowd that is now stirred to awakening, now incited to
offense; the restless crowd that cries and cries, today: Hosanna! tomorrow: Crucify!

Here we have, as in a silhouette, the original presuppositions and conditions on which the
existential criticism of the faith of the Gospel must further be laid out. And now the world
may go as far as it will, and become as old and as enlightened and as clever and as
scientific as it will: it will nevertheless never get beyond these crises, it will never get
beyond this criticism.

The existential presentation of the life of Jesus ought, according to the evangelists'
direction, surely also to be pragmatic. But the pragmatism here in question is not the
scientific one. In scientific pragmatism the particular is constantly posited as a fleeting
moment in the whole. Hereby finitude's schema acquires so decisive a preponderance that the
spirit disappears. So the whole is parceled out according to the various journeys and the
various feasts; so the theologians cannot agree about the chronology; so one contends about
the Synoptics, so one wrangles about the Sermon on the Mount, so one reasons about ``the
Easter reckon%
% --- p. 94 ---
ing''; so learned science has forgotten faith, and moves in nothing but finitude.

In evangelical pragmatism, on the other hand, the whole is posited as concentrated in the
particular. In each single redeeming deed the whole Redeemer himself is present, in each
single Gospel the whole evangelical history. Finite temporal sequence is hereby reduced to a
vanishing moment in the grouping of the central events. The supernatural phenomena and
extraordinary deeds in the life of Jesus follow one another so suddenly, so surprisingly,
and burst the ordinary connection of the temporal so entirely, that a finite pragmatism
becomes impossible. No disciple could relate himself objectively to these
events,\footnote{``To become a historical eyewitness is thus easy for the contemporary
learner; the misfortune, however, is that to know a historical circumstance, indeed to know
them all with the eyewitness's reliability, by no means makes the eyewitness a disciple ---
as one can see from the fact that this knowledge means for him nothing but the historical. It
shows itself at once here that the historical in the more concrete sense is a matter of
indifference; we can let ignorance enter in relation to it, and let ignorance as it were
annihilate the historical; if only the moment still remains, as the point of departure for
the eternal, the paradox is present. If there were a contemporary who had himself restricted
his sleep to the shortest time in order to follow that teacher, whom he followed more
inseparably than the little fish that follows the shark; if he kept a hundred spies in his
service who everywhere spied out that teacher, and with whom he himself conferred every
evening, so that he knew that teacher's description down to the least detail, knew what he
had said, where he had been at every hour of the day, because his zeal let him regard even
the most insignificant thing as important --- was such a contemporary the disciple? By no
means. He could wash his hands if anyone would charge him with historical unreliability, but
no more than that either. If there were a contemporary who had stayed in foreign lands and
first returned home when that teacher had only a day or two of his life remaining; if the
contemporary were again kept by business from getting to see that teacher, and came only at
the last, when he was about to give up the ghost --- was this historical ignorance a
hindrance to his being the disciple, if the moment was for him eternity's decision?'' Johannes
Climacus. Philosophiske Smuler, pp.\ 85---86.} no spectator could pettily record every single
trait,
% --- p. 95 ---
no apostle could, so to speak, keep a diary of Jesus' deeds (cf.\ John 20:30---31; 21:25).

The existential presentation of the life of Jesus can doubtless not be philosophical, but must
nevertheless be laid out in such a way that the evangelical history is seen in the ideal's
light, as an eternal and sacred history exalted above past and present.

% The Danish here reads „Fremstilligen“ (missing n) for „Fremstillingen“ ---
% printer's error, kept as printed there. The English has the correct word.
The presentation of sacred history has then the peculiar determination of casting light over
the world's history and making the ideals recognizable, so that everyone must judge himself
when he judges of the revealed ideal that reveals the thoughts of hearts.

If there is no blessed life, then it is doubtless greater to take the city than merely to rule
over one's own mind, and then Christ too is an offense; but if there is a blessed life: yes,
then he who rules over his own mind is greater than he who takes the city, and then Christ is
the Savior of the world. If there is in faith no eternal life, then the individual must seek
his ideal in worldly society, and worldly society its ideal in the history of the race. The
spirit of history then reveals the very highest ideal of kingdom and power and glory when the
city is taken, when the heroic genius says \textit{jacta alea esto} and crosses the Rubicon
with his legions. But if there is in faith an eternal life; if the individual must buy it with
self-denial, preserve it in patience, and consummate it in self-sacrificing love: then he who
is crowned with thorns is surely also higher than he who is crowned with laurel; then the
ideal's way does not go over the Rubicon to the Capitol either, but the way goes over the brook
Kedron to the garden of Gethsemane; and then the highest watchword of resolution is
% --- p. 96 ---
not a \textit{jacta alea esto} either, but the highest watchword is in the words of the prayer:
Father, not my will, but thine!

% The Danish here reads „er der er saligt Liv“ for „et saligt Liv“ (cf. p. 95,
% „er der et saligt Liv“) --- printer's error, kept as printed there. The
% English has the correct word.
Let everyone judge of this as he will: he judges himself. But this stands fast: if there is a
blessed life, then Christ is the Savior of the world; if Christ is the Savior of the world,
then Christianity is higher than science; but if Christianity is higher than science, then
\emph{the faith of the Gospel is indeed also --- heterogeneous with theology.}

%% ============================================================
%%  EIGHTH LECTURE (pp. 97--111)
%%  NB: head reads "halvphilosophiske"; the Indhold reads "halvspeculative".
%%  Real variant in the 1850 printing --- both witnesses kept. Do not
%%  normalise. (Head -> "half-philosophical"; Indhold -> "half-speculative".)
%% ============================================================
\chapter*{Eighth Lecture}
\addcontentsline{toc}{chapter}{Eighth Lecture}
\label{ch:forel8}
\markboth{Eighth Lecture}{}

\begin{center}\itshape
The half-philosophical image of Christ that the speculating dogmatician
conjures up consists in a mystical fusion of the scientific and the
religious, and is on that account so infinitely different from the
evangelists' Christ that it falls even outside the religious-ethical
horizon: the faith of the Gospel is heterogeneous with theology.
\end{center}

% --- p. 97 ---
Although the theologians of course cannot accept what, in relation to the present age's
enlightenment, they must call ``the crude principle of inspiration,'' and are on that
account in much doubt about the particulars of evangelical history and the genuineness of
the four Gospels: they nevertheless console themselves that pious contemplation still
rediscovers in each of the evangelists the same historical main features, the same
fundamental religious view, the same image of Christ. The theological voice is never more
moved and full-toned than precisely when the talk is of the leading thought, of the
fundamental view, of the image of Christ. Even when one comes into embarrassment with all
the particulars, one can still cheerfully hold to the fundamental view; even when science
has transformed the evangelists' Christ into a half-philosophical fancy, one still always
sees an image of Christ. --- While the criticizing interpreter is chiefly interested in the
historical particulars in the life of Jesus, the speculating dogmatician is particularly
interested in the idea, in the fundamental view, in the image of Christ.

% --- p. 98 ---
In the idea and the fundamental view the theologian now supposes he reconciles Christianity
with science, in that out of Christianity's own content he develops a speculative
knowledge, a knowledge that in all essential points is supposed to satisfy faith.

In order to discover the error on which this attempt at mediation, of great antiquity in
the Church, rests, we will here dwell a little more closely on the often-touched opposition
between scientific and existential knowledge.

In scientific knowledge the knower's particular mood and subjective state signify nothing;
in existential knowledge everything depends on the knowing subject's peculiar state and
mood. In scientific knowledge the knower must abstract from himself, forget himself, give
himself over to contemplation's calm, and relate himself generally to the general object. In
existential knowledge the knowing subject must go into himself, place himself in a decisive
relation to the object --- for which reason everyone in this personal relation of
subjectivity, inwardness, and interestedness rightly calls the object his object. Scientific
knowledge, which moves in the general, is passionless, and ought to be so, since it is
impersonal; existential knowledge, which moves about the punctual, the particular, the
proper, is passionate, and must be so in the same degree in which it is personal.

From this it follows that the knowledge won in science has a direct significance only for
science, and must change its character when it is applied to existence; and conversely, that
the knowledge attained by existence's way has a direct significance only for existence
itself, and must change its character when it is translated into science. Since now these
two kinds of knowledge are in this way so heterogeneous that they cannot be exchanged for
one another, it follows in turn
% --- p. 99 ---
that they can neither relate themselves directly to one another nor be directly corrected by
one another. Scientific knowledge cannot be directly corrected by existence, any more than
the existential by science.

What has here been said has (as Aristotle already points out) a great significance for
rational ethics, but acquires a still higher and quite differently decisive significance for
that peculiar sphere which in a stricter sense we must call the sphere of the faith of the
Gospel. As regards the knowledge of Christianity, we must in the most definite terms hold
fast to the rule that science cannot be corrected by faith, any more than faith can be
corrected by science.

If we now apply this rule to our present theme, then the error in the attempt at mediation
spoken of will without doubt also come to light.

In order to work out the christological idea and construe the so-called image of Christ, the
speculating dogmatician goes to work in the following way. The christological main
propositions on which the dogmatic development of the idea is to be grounded he does not
borrow from any philosophy standing outside the Church; he fetches them all from Scripture
and the Church's doctrine. Since now every main proposition must to that extent be said to
be a proposition of faith, the dogmatician supposes that he thereby satisfies faith, builds
on faith's ground, knows in faith and out of faith. But what knowledge is it, then, that he
here strives after? Well, it is to be a scientific and, so far as possible, even a
speculative, at least a half-speculative knowledge. Now the error begins. The ecclesial
propositions of faith are converted into the form of objective knowing, and treated as
general scientific propositions; as these scientific propositions are worked together into a
scientific totality, the christological idea is then worked out, and the image of Christ
shows itself.

That such a christology by no means reconciles faith with
% --- p. 100 ---
knowledge, but rather shows itself as an ambiguous fusion of heterogeneous constituents, a
fusion that conflicts in equal degree with the nature of faith and of knowing, can be proved
without difficulty.

As regards the dogmatician's half-speculative christology, it must at once awaken misgiving
that it lays claim to a double acknowledgment, in that it will both be acknowledged as a
product of faith and yet at the same time be acknowledged as a product of science; it is
like an offspring that has faith for mother and science for father, and now with this its
double nature makes a double claim, one on the mother's and one on the father's behalf. With
this double claim one would now also expect that the christological idea was willing to
submit to a double testing, namely a testing in faith and a testing in science; but nothing
comes of this testing either. For if we apply faith's existential measure to the
christological idea, it evades the test by referring us to science; if we then will apply
science's strict measure, it again evades the test, and refers us to faith. In this way the
dogmatician's christological idea can at bottom never be tested, unless it were before a
tribunal that is higher than both faith and science; but where, then, is this tribunal?

The christological idea does not satisfy science; for what science is to acknowledge, it must
also be able to test; but the christological idea science cannot test. For since this idea
has its origin not from the natural system of reason in the present, but from revelation's
kingdom beyond, it is of miraculous origin and has a constituent in it about which science
can say nothing, since it cannot test it --- namely the supernatural constituent. --- Christ
became man when he was conceived by the Holy Spirit and born
% --- p. 101 ---
of the Virgin Mary. Half-speculative christology can now doubtless put forward various things
in general about what it calls the incarnation's ``cosmic miracle''; but about \emph{this}
incarnation, about \emph{this} conception and \emph{this} birth science can nevertheless
settle nothing, since science in no way sees itself able to test the miracle. --- Christ is
true God and true man. Half-speculative christology can now put forward various things about
the God-man in general, about Christ's double nature and double will in general; but about the
actual unity of the two natures, how it came about that Christ could really suffer and die on
the cross although he himself in the moment of suffering and death continued to be the true
God, the God who bears heaven and earth --- of this science cannot judge, of this it can
express itself neither affirmatively nor negatively. But if science cannot judge christology's
content, how then should it be able to acknowledge the christological presentation. ---
Christ rose from the dead, ascended into heaven, and sat down at the right hand of the Father.
There was a time when science in earnest supposed that it could deliver a judgment on this
exaltation of Christ: alas, what vehement controversies have the theologians not conducted
about the ubiquity of Christ's body, about \textit{ubiquitas absoluta} and \textit{ubiquitas
respectiva}, and what offenses have these controversies not awakened! When such theological
controversies can no longer set science in passion, should this not be a proof that one has now
at last come to consciousness that such things lie outside the bounds of human knowing; but
what lies outside the bound of our objective knowing cannot possibly be treated scientifically.

Neither does the christological idea satisfy faith, and that for the simple reason that every
proposition of faith loses its character when it is converted into the form of knowing. The
so-called
% --- p. 102 ---
reconciliation of faith and knowledge begins involuntarily with the fundamental fault of
placing faith on the same level as knowing, and thus moves in a circle.

To speak of knowing as of a free-standing concept independent of the knower is in order, since
knowing is objective; but to treat faith as a general being, grounded in itself, independent of
the believer, is unreasonable, since faith is subjective. Faith is only for the believer; the
believer is an actual human being existing in subjectivity's inwardness; faith's category is
existence's highest category, its knowledge existential. If one nevertheless commits the
blunder of abstracting from the believer when one speaks of faith, just as one abstracts from
the knower when one speaks of knowing, then one does indeed get the result that faith is
objective just as much as knowing; but therewith one has at the same time killed faith, for
faith dies in objectivity.

That knowing is objective is recognized from the fact that it is grounded on universally valid
presuppositions, on objective proofs, and will have nothing to do with subjective grounds; that
faith is subjective is seen from the fact that it is unconditionally grounded in subjectivity's
inwardness, and therefore will know nothing at all of the objective grounds. As now he who
supposes he is helping science by mixing subjective grounds into the objective proofs offends
and bungles science, so one again offends and bungles faith when one attempts to support
subjectivity's decisive ground by objective proofs.

In order to justify his half-speculative christology, the dogmatician is --- characteristically
enough --- on the whole more attentive to the Epistles than to the Gospels, and, as regards the
Gospels, more attentive to John than to the Synoptics. The reason for this is evident enough.
The Gospels, especially the ``synoptic'' ones, keep at every point within existence, whereas
% --- p. 103 ---
the Epistles, on the other hand, which put forward the doctrine in greater generality, seem to
come far nearer to scientific discourse. In the Epistles, therefore, especially in the Pauline
Epistles, one has supposed one found the type for a speculative christology, and the Apostle
Paul has in this way been set up as a model for speculative dogmaticians.

That there really do occur in the Apostle Paul, as also in the evangelist John, many
propositions and expressions that dogmatizing speculation can in its own way take to its
credit shall never be denied; but from this it does not yet follow that these Johannine or
Pauline propositions themselves rest on a speculative ground.

The Pauline doctrine of Christ is just as little speculative as the Johannine doctrine of the
Logos is speculative. What John and Paul express in this direction they express not in the
name of science, but in the name of faith, with faith's subjective assurance, with
inspiration's divine authority. When it is said in John's Gospel: ``In the beginning was the
Word, and the Word was with God, and the Word was God. The same was in the beginning with God.
All things were made by it, and without it was not one thing made that was made'' . . . ; when
it is said further (v.\ 14): ``The Word became flesh and dwelt among us,'' and so on --- then
these propositions doubtless sound speculative to a speculative ear; but if we look a little
more closely, it will soon show itself that such statements neither are nor can be speculatively
meant.

If the Johannine doctrine of the Logos were to have speculative significance, then it would
surely have to be laid out in such a way that it could be justified metaphysically; but
% The Danish here reads „mataphysisk“ for „metaphysisk“ --- printer's error,
% verified on the page image and kept as printed there. The English has the
% correct word.
this is not at all the case. The proposition \emph{the Word became flesh} is no scientific
proposition and can never become one. ``The Word became flesh'': the supernatural took on our
nature, the supra-historical became historical; such a composition of such heterogeneous
factors
% --- p. 104 ---
science cannot master, and what science cannot master it will not judge either. If science were
to settle anything about the proposition \emph{the Word became flesh}, then the law for the
supernatural's interaction with the natural would first have to be discovered (as natural
science has discovered laws for the interaction of natural things); but this law is to science
a pure X, an entirely unknown magnitude:
% The X is set in the surrounding Fraktur fount in the Danish, not antiqua, so
% no \textit{} there or here.
here no scientific test can be instituted; but what science cannot test, science will not judge
either. If the Johannine proposition is to be grounded scientifically, then things look
sorry indeed, both for John's Gospel and for Christianity; for the grounding is really
entirely impossible. But if scientific grounding is in this case entirely impossible, then it
is surely also, as regards faith, entirely superfluous.

The Johannine proposition is an existential proposition of revelation. When John says: ``The
Word became flesh and dwelt among us,'' he says it in the name of faith; the evangelist's faith
cannot be separated from the evangelist himself: faith cannot be separated from the believer.
In speaking out of his faith, the evangelist does not speak about faith in general, but he
speaks as a believer to believers: the relation is existential, not scientific. The others who
would convince themselves of the truth in the evangelist's word cannot then be convinced by
objective grounds in general either. The only ground on which the conviction can be built is
that the evangelist is inspired; but that the evangelist is inspired cannot be proved by
objective grounds either: it must be believed. If I am to believe the evangelist's word, then I
must believe that the evangelist is inspired; if I am to believe that the evangelist is
inspired, I must also believe the evangelist's word: the relation between the evangelist himself
and all those who by the evangelist's word are to come to faith is again existential, not
scientific.

% --- p. 105 ---
From this one sees at the same time what one is to judge of the christological passages found
in the Apostles' Epistles. Already the manner in which these Epistles came into being must in
the highest degree warn us against scientific misunderstandings. How, pray, does such a letter
come about? Confusion arises in a congregation; self-willed people, immoral persons, false
teachers awaken offense. The report of this reaches the absent apostle and awakens his concern.
With the zeal of one concerned, with his gaze fixed on the offenses taking place, with all the
haste that the danger of the moment makes necessary, the apostle now sits down and writes his
letter. In this letter, whose purpose is to punish, warn, admonish, instruct, encourage, and
comfort, the christological propositions of revelation are now interwoven into the practical
instruction, so that the one as it were involuntarily flows from the other; as for example when
Paul (Phil.\ 2:1--13) will inculcate humility in the Philippians, and he then (v.\ 5--13) says:
% The Danish prints the second reference „(V 5--13)“, without the period after
% V that „(Phil. 2, 1--13)“ has --- printer's error, kept as printed there.
``For let this mind be in you which was also in Christ Jesus, who, being in the form of God,
thought it not robbery to be equal with God, but emptied himself and took upon him the form of
a servant, and was made in the likeness of men; and being found in fashion as a man, he humbled
himself and became obedient unto death, even the death of the cross. Wherefore God also hath
highly exalted him and given him a name which is above every name; that at the name of Jesus
every knee should bow, of things in heaven and things on earth and things under the earth, and
every tongue confess that Jesus Christ is Lord, to the glory of God the Father. Therefore, my
beloved,'' and so on.

But what the apostle thus writes in the danger of the moment, in the anxiety of concern, in the
zeal of love, so that his heart burns --- behold, this the theologians then in all scientific
peace and leisure make the object of objective scientific contemplation: does one
% --- p. 106 ---
not then notice the disproportion! What faith's apostle writes to believers for faith's
furtherance in faith's name --- behold, this the theologians then suppose they can treat as
half-speculative contributions to a half-speculative understanding of the half-speculative
relation between Christ's divine and human nature (κρύψις or κένωσις?), between the God-man's
humiliation and the God-man's
% „eller“ is set in antiqua between the two Greek words in the Danish;
% transcribed plain there, and left plain here.
exaltation (\textit{status exinanitionis} and \textit{status exaltationis}); does one not then
notice the disproportion!

If already the apostolic letter-form and the peculiar conditions under which these letters came
into being must warn us against the speculative-scientific conception, how much more must we not
be warned when we consider that these letters have streamed forth from inspiration's spring. The
spirit that inspires the scientific author is not the apostolic one. If an inspired apostle
arose in the present age, he would surely have something else to do than to be an author, and
would then teach us something quite other than what we can learn from scientific books. Where
religion's spirit has a power such as it must have in an apostle's life, the religious is far
too jealous for itself and far too zealous for its proclamation for the inspired one to get
peace for scientific contemplations and aesthetic enjoyment of ideas. An apostolic word is an
apostolic deed; an apostle truly has no time to build theories; even when he writes his longest
letter and communicates the doctrine in the fullest form, every word is seized out of his own
life, and the communication point for point existential.

In half-speculative dogmatics the incarnation and inspiration are now regarded as the two great
fundamental miracles that condition and explain one another mutually. He who does not accept
inspiration cannot grasp the incarnation either; as one judges of Christ's apostles, so one
judges of Christ himself: the true image of Christ must be seen in the light of the apostolic
spirit. What dogmatics thus says in general we must also grant
% --- p. 107 ---
in all generality; but does this general concession nevertheless lead us into speculative
contemplations?

As the principle of inspiration is a supra-scientific principle of existence that holds only for
the believer in faith, so also the principle of incarnation is a supra-scientific principle of
existence, and the relation between the two an existential, not a speculative relation.

The existential doctrine of Christ that can be derived from the Apostles' Epistles is at all
points in agreement with the doctrine of Christ that lies in the confession of faith and in the
Gospels; but the half-speculative christology that the theologians abstract from Scripture and
the Church's doctrine will never rightly agree with the Gospels' Christ.

Theological christology will now, for example, either not let Christ be tempted at all, because
his temptation in the wilderness seems to it so unreasonable (in this case the dogmatician then
straightforwardly defies the evangelists); or else scientific christology does indeed grant a
temptation, but without entering more closely into the temptation of which the evangelists
speak. The evangelical history of the temptation is not at all laid out for scientific
treatment. What a half-speculative christology has to tell us about Christ's temptation in
general is, on the other hand, when it is said most profoundly of all, precisely calculated to
make a scientific impression. But this scientific impression --- how different it is from the
wondrous, repellent, offensive impression we feel when we read with a ``critical eye'' the
Gospel of the temptation in the wilderness.

He who reads half-speculative christology gets an impression as if human understanding were now
in the best understanding with the revealed fact, as if now everything, even the most wondrous
of all, even the temptation in the wilderness, were quite near to being resolved into science's
christological idea; he, on the other hand, who gives himself over to the impression of the
evangelical history of the temptation feels involuntarily that he must either reject
% --- p. 108 ---
this temptation as an absurd fancy, or else take reason captive under faith's obedience.

If the Gospel now at all points makes the impression that one is to believe with understanding
\emph{against} understanding, while dogmatics at all points endeavors to produce the impression
that offense's sting had disappeared in science, so that at speculation's present stage one
should really only believe \emph{with} understanding: then it is surely evident enough that the
Gospel and dogmatics work --- not with one another --- but against one another.

Therefore the theologians have also, not only on historical-critical but at the same time on
metaphysical-dogmatic grounds, always so much to correct in the four evangelists; it is so very
seldom that an evangelist is so fortunate as to be able to suit them properly. In recent times
the evangelist John, for example, has had to hear ill of it because, contrary to christology's
scientific rule, he lets Christ pray aloud at Lazarus's grave. ``Das Unpassende dieser Wendung''
--- says De Wette --- ``muß man anerkennen . . . . Denn in der That dürfen und müssen wir uns
Jesu Bewußtseyn von Gott und seiner Stellung zu ihm als ein ununterbrochenes, nicht an einzelne
Gebetsmomente geknüpftes denken.''%
\footnote{Cf.\ Evangltr.\ og den mod.\ Bev., pp.\ 438--58.}

But suppose even that the dogmatician's half-speculative image of Christ were in all main
features as like the evangelists' Christ as a copy can resemble its original --- we should be
no nearer for that. The dogmatic image of Christ falls under a quite different horizon, and is
entirely heterogeneous with the evangelical Christ. The evangelical Christ stands in actual
existence, and at every moment makes existence's demand of everyone who approaches to
contemplate him; but the dogmatic image of Christ stands only in contemplation,
% --- p. 109 ---
and belongs to the aesthetic-scientific gallery of ideas, where everyone who has any ``deeper
sense'' can go to enjoy the spectacle. To paint an image of Christ with words and thoughts is
doubtless a greater art than to paint it with lines and colors; the task has accordingly not yet
been solved by any speculative-Protestant dogmatician, as it was long since solved by the
Catholic painters of the Middle Ages.

But even if the task were solved with astonishing success; even if there arose a
dogmatic-speculative genius who succeeded in painting us an image of Christ with such
philosophical profundity, such inwardness of piety, such purity in thought, such clarity in
expression, that Christ now really stood before us in this image as ``the immediate appearance
of divine holiness on earth,'' so that ``the highest love and the heart's humility and devotion
in God were here seen to melt together into one with the expression, inspiring both fear and
reverence, of a righteousness whose foundation nothing can shake, and with the mystical
expression of a power that could bid all forces and all powers serve it''; --- suppose that the
solution of the task succeeded for the dogmatician on so extraordinary a scale: then such an
image of Christ would nevertheless still be just as far from making a specifically Christian
impression as the painter's image.

For just as the painter's image makes an aesthetic-artistic impression, and therefore gladdens
and attracts all those who have a sense for sacred art without regard to whether they are in
actuality believing or unbelieving, so the aesthetic-speculative image of Christ would make an
aesthetic-speculative impression, and therefore in a high degree gladden and attract all those
who have a sense for speculative imaginative intuition; but this image of speculation would
nevertheless by no means stand before the beholder as ``\emph{the sign of contradiction},''
would nevertheless not have power to ``\emph{reveal the thoughts of hearts},''
% --- p. 110 ---
would nevertheless not be ``\emph{mighty in word and deed}'': would nevertheless not be able to
make a true evangelical impression as ``\emph{the Savior of the world!}''

If one cannot relate oneself aesthetically-speculatively to the Christ of the faith of the
Gospel%
\footnote{On this cardinal point let one read ``Philosophiske Smuler'' by Johannes Climacus.},
then neither can one relate oneself in evangelical faith to speculation's image of Christ.

It can indeed sometimes happen that one or another comes forward and testifies that it is now
really the dogmatic christology that has properly first grounded and clarified his faith for
him, indeed that only in the profoundly speculative image of Christ has he really properly seen
who Christ is; but with such testimonies one ought nevertheless to be cautious. It is
unhappily only too true that we human beings in the nineteenth century are very inclined to
confuse our aesthetic-speculative interest in the idea with the religious-ethical interest in
the matter; how easily then may it not happen that someone really supposes himself to be
strengthened in Christianity, although he has, properly speaking, only been strengthened in an
aesthetic-speculative love for an aesthetic-speculative image of Christ!

By many ways Christianity comes to human beings, by many ways human beings come to
Christianity. We readily grant that theological christology can also in its own way become the
occasion for a human being to go into himself and reflect on the Gospel; but we must therefore
none the less make a distinction between the matter itself and the occasion of the matter.

One can with reason say that it stands with the dogmatic images of Christ in the Protestant
Church as with the aesthetic images of the Madonna in the Catholic: various things are told of
their power of conversion, but this power of conversion is somewhat ambiguous. In order to
remove the ambiguity, let one resolve
% --- p. 111 ---
to give faith what is faith's, and science what is science's.

By this apportionment dogmatics' half-speculative image of Christ will, insofar as it consists
in a mystical fusion of faith and knowledge, doubtless for a great part fade and vanish; but the
Gospel's Christ will conquer and win, when it is recognized in truth and proved in deed
\emph{that Christianity is higher than science, and the faith of the Gospel heterogeneous with
theology.}

%% ============================================================
%%  NINTH LECTURE (pp. 112--128)
%% ============================================================
\chapter*{Ninth Lecture}
\addcontentsline{toc}{chapter}{Ninth Lecture}
\label{ch:forel9}
\markboth{Ninth Lecture}{}

\begin{center}\itshape
Theology's half-speculative dogmatics cannot justify its ecclesial
significance from the history of the Church, since the history of the Church
proves, on the contrary, how dogmatic prejudices have in this respect harmed
the Church, and what a thoroughgoing opposition there is here between
theological and ecclesial dogmatics. Theological dogmatics must fall before
the proposition that Christianity is a paradox: the faith of the Gospel is
heterogeneous with theology.
\end{center}

% --- p. 112 ---
The series of objections advanced in these discourses, both against theology in general and
in particular against the present age's theologians, will perhaps have a repellent effect on
many --- even to the degree that, without discussing the matter more closely, they at once
reject the view here set up as a one-sided and arbitrary conceit.

The misgivings that press forward on this occasion are then doubtless not few either. To keep
to but one question: ``how should Christianity be able to tear itself loose from theology,
since the Church cannot possibly do without its firm structure of doctrine, its dogmatics?
Christianity is indeed a life, but it is also a doctrine, and it is the one only insofar as it
is the other. It is the Church's requirement that the Word be taught purely and cleanly; but
how should God's Word be able to be taught purely and cleanly when one casts contempt on the
objective, systematic concept of doctrine? Doubtless the Apostles' Creed furnishes a
% --- p. 113 ---
type for the Church's doctrine in general, but the Church's history shows sufficiently why one
could not stop at this first general foundation. Under continued controversies with the various
heretical parties the Church soon saw itself compelled to articulate the Christian content in
ever more definite expressions, in order thereby to ward off the un-Christian. Thus the symbols
were formed at the ecumenical synods (at Nicaea 325, at Constantinople 381, at Chalcedon 451,
at Constantinople 680) under vehement theological struggles. These theological struggles are to
be regarded not merely as theoretical party controversies, but at the same time as the
magnificent religious crises the Church had to go through in order rightly to take possession of
its Christianity and secure objective truth against the arbitrariness of subjective opinions.''

``If from the beginning one has thus not been able to refrain from entering upon a learned
scientific treatment of the individual doctrinal propositions, then it follows of itself that one
has not been able to refrain from going through all the dogmas in their inner organic connection
either --- or, in other words, from treating the whole doctrine of the Church systematically.''

``As it went with the formation of the Catholic concept of doctrine, so it went later also with
the formation of the Protestant. Scarcely had the Reformation's religious breakthrough taken
place before a scientific reworking of the Church's doctrine also followed. Luther was more
pastor than theologian; Melanchthon was more theologian than pastor; Calvin was with equal
distinction both theologian and pastor. From this one sees that theology is not an arbitrary
conceit of individual learned men, but an organic product of ecclesial life.''

``Doubtless it is not necessary for the individual layman to procure systematic insight into the
individual parts of the concept of doctrine; but it is nevertheless none the less necessary for
the ecclesial teaching order. If the congregation cannot do without the confession and the
Gospel, then
% --- p. 114 ---
the teaching order cannot do without the concept of doctrine and dogmatics either. Doubtless the
chief Christian doctrines cannot be set forth for the congregation in their dogmatic connection
or scientific grounding; but it is nevertheless none the less a truth that Christian preaching
stands in an organic reciprocal relation to Christian dogmatics, so that one can therefore say:
to such a dogmatics corresponds such a sermon, and to such a sermon such a dogmatics.''

``If it must now be granted that preaching presupposes dogmatics, and dogmatics theology: should
there then not be contained in this a serious admonition as to how questionable it is to want to
undermine theological science --- assuming that one does not therewith at the same time intend
to undermine Christian doctrine?''

We have now seen that there are misgivings. But however weighty these misgivings may appear to
us at first glance, they are nevertheless, seen more closely, of far too fleeting a nature for
one to be able to build on them any reliable defense of the prevailing theology. --- It will
hardly have escaped the attentive observer that such misgivings, occasioned by surveys of church
history, always move in so vague a generality that they precisely thereby divert attention from
the real question of principle. The church-historical considerations by which one might in this
respect attempt to defend Christian science and re-establish the disputed authority of
theological dogmatics can at once be given such a turn that they let themselves be used with the
greatest emphasis for a direct attack on theology itself.

In order, then, to get away from the vague generality --- back to the question of principle ---
I shall here permit myself to direct attention to some prejudices, which one might perhaps best
designate by the expression \emph{dogmatic prejudices.}
% --- p. 115 ---
Dogmatics lays claim to being treated systematically; I will this time comply with that, and
treat its prejudices systematically. I place these systematic prejudices in a systematic
relation 1) to heresy, 2) to Christian zeal for faith, 3) to higher scientific rigor.

1) When the Church makes use of an age's scientific forces to discover the Christian element in
a disputed doctrine, it by no means lies in the Church's idea to let science, as science, settle
what is Christian and not Christian. The Church uses science only as means and instrument in
order to arrive at clarity about its own tasks. Science has the office of proposing various
modes of expression and formulas for the doctrine's content; but the Church reserves to itself
the choice among these forms of expression according to its own judgment: only the Church has
authority to stamp the definite doctrinal proposition with faith's seal, to sanction it as
dogma. When the doctrinal proposition is sanctioned, it is an object for believing appropriation,
and has nothing further to do with science. This is the relation of principle between the Church
and science.

But since science nevertheless had from the beginning to be set in motion every time a new dogma
was coming into being, the blunder lay only too near at hand that science, while it worked for
the Church, now also at the same time worked a little on its own account and speculated a little
about the dogma, in order after all to make it a little comprehensible. So it came about that
science became theology, and theology nourished heresy. --- How does heresy arise? It arises
from the fact that one will reason one's way to the dogma and reason about the dogma.

Had the Church been able to keep faith's principle pure, then the conflict between orthodoxy and
heresy would from the beginning have transformed itself into a direct and open conflict between
faith and unbelief; the heretic's unbelief would then have shown itself precisely
% --- p. 116 ---
in this, that he corrupted the doctrine because he reasoned about the dogma. Instead of reasoning
with the heretics, the orthodox would then have had to confine themselves to preaching conversion
to the erring. This principle the Church could not carry through; its teachers became
theologians, and the theologians then shared with the heretics the error that the Church's
doctrine ought in general to be an object of reasoning and speculating. The vehement heresy
conflict thus did not become any pure conflict of principle between faith and unbelief, but was
corrupted into a half-scientific conflict between a supposedly true and a supposedly false
speculation, a supposedly true and a supposedly false theology.

This conflict then bore in itself the contradiction that is quite fitted to inflame the passions
to the very utmost. Neither of the contending parties could here come to see its own wrong. The
orthodox theologians felt within themselves that they had the Church and Christianity on their
side, and that their theology consequently must be right in faith; but now the same theology was
at the same time to be right in speculation: this prejudice inflamed the heretics. If there was
to be reasoning and speculating at all, then the heretic held himself fully and firmly convinced
that his own speculation was far more reasonable and far more sensible than the Church's. Far
from seeing that it was he himself who was astray, the heretic now rather believed himself called
to save the Church from a disastrous aberration, from a doctrine so unreasonable that it
conflicted both with true speculation and with sound understanding.

% Both maxims are letterspaced in the Danish print; verified at 400 dpi.
``\emph{What one is to believe, one ought so far as possible also to see one's way to
comprehending}'' was the orthodox theologian's principle; ``\emph{what one cannot comprehend,
one ought not to believe either}'' was heretical theology's principle. The one of these
principles supports the other. Arius became a heretic, and why? because he could not comprehend
that the Son had the same
% The Danish here reads „ramme“ for „samme“ --- printer's error (an
% unmistakable Fraktur r), kept as printed there. The English has the correct
% word.
% --- p. 117 ---
essence as the Father; although Athanasius did endeavor to make it comprehensible to him.
Nestorius became a heretic, and why? because he could not comprehend that the two natures were
one in Christ's person; although Cyril with all his adherents labored to make it comprehensible
to him. Pelagius became a heretic, and why? because he could not comprehend how human freedom
was resolved into divine grace; although Augustine summoned up all his profundity to make it
comprehensible to him.

To touch on but one main point in the Protestant Church: how have the Lutherans and the Reformed
not denounced one another as heretics in the dispute over the Lord's Supper, so sorry for both
parties! Luther was right, he felt it deeply; it was this feeling that made him so strong in
confession, so bitter in conflict. Yes, assuredly Luther was right; modern Christendom is hardly
conscious of what the Church owes to the man who on this point stood so firmly on the text's word
and defied a world with faith's Gospel. Yes, Luther was right --- yet not as theologian, but as
pastor. His conception of the Lord's Supper explains itself precisely from the fact that he
originally saw the sacrament not with the theologian's but with the pastor's eye. What Luther as
theologian has put forward about the doctrine of the Lord's Supper is significant only for his
own time; his theologoumena are vanishing, but his pastoral testimonies will remain as long as
faith is --- faith.

Luther himself was not clear about the relation of principle between religion and theology. The
new doctrine needed science's assistance; Melanchthon was a theologian, and the theologian then
supposed that Calvin was right; for Calvin was also a theologian. But the Lutheran Church, which
was unconditionally right in faith, would now at the same time also be unconditionally right in
theology; so it went the longer the worse. Lutheran theology worked itself weary in unspeculative
speculations, and has long had to see its favorite theories about \textit{communicatio
idio%
% --- p. 118 ---
matum} and \textit{ubiquitas corporis Christi} abandoned to dissolving criticism.

The old dispute between the Lutherans and the Reformed has fallen asleep in the age's
indifferentism, but wakes up again every time worldly prudence sets about burdening consciences
with a union calculated on mere humanity. When, then, will the dispute end? It ends only in
spirit and truth when both churches, the Lutheran as well as the Reformed, resolve to give up
their theological-dogmatic prejudices and hold in faith's inwardness to the word of institution,
to the Gospel's text.

2) Ecclesial zeal for objective orthodoxy has often been confused with true Christian zeal for
faith; this confusion has chiefly its ground in a dogmatic prejudice.

The dogmatician relates himself objectively to Christian truth, and then supposes he is spreading
Christ's kingdom by procuring for this truth a certain objective victory in the objective world.
Here is the point where theological dogmatics coincides with the state church, whether this be
otherwise Catholic or Protestant.

Orthodox theology's objective dogmas are put forward as pure objective truths, and in such a way
that no one who accepts Christianity at all shall be able to avoid becoming convinced of these
truths, unless he is either stupid or obstinate: now comes the self-contradiction. The Christian
truths are assuredly objective; but they are not objective in the sense in which mathematical,
philosophical --- in short, scientific --- truths are objective.
% NOTE (from the transcription): spacing.py flagged a run at „ere ikke“ here.
% The 400 dpi image shows ordinary spacing; the apparent gap is a later
% reader's pencil underline running under the whole clause. No emphasis in the
% print, and none here.

Scientific truths gain entry by conviction's objective way. It would here be as great a folly to
apply compulsion as to show obstinacy. With the Christian truths it is otherwise. These truths
can only make their way in by conviction's subjective
% --- p. 119 ---
way. If they are nevertheless to be inculcated as pure objective, general truths, they must of
course meet with contradiction; this contradiction then becomes a sign of obstinacy, and against
the obstinate one applies compulsion.

So it goes: orthodoxy's strict theologians discover a heresy; they summon the heretic to appear
at a hearing, a colloquium, a disputation, so that he may be refuted by sound grounds and
scientifically convinced of his dogmatic error. What happens? In the course of the disputation it
shows itself that the heretic cannot be refuted, since he of course has just as many grounds for
his opinion as the orthodox have for theirs. Then passion comes to a boil, and the matter takes a
subjective turn. The accused is now no longer regarded as a poor erring human being who lacks
true insight because he lacks dogmatic grounds, but he is treated as a stiff-necked, hardened
sinner who out of sheer obstinacy will not let himself be convinced. But convinced he shall be
--- by objective grounds. If science's grounds will not suffice, the state comes to the aid with
its grounds; and behold, then the hardened heretic is also quite properly convinced subjectively
--- by objective grounds.

But what do we now call it, to treat a purely subjective matter as if it were purely objective? In
relation to the scientific misunderstanding we call it stupidity; but in relation to the
passionate zeal that prefers to work by outward compulsion, we call it tyranny.

That the Catholic Church so early became a church of compulsion, and as a church of compulsion
gave itself over to tyranny in doctrine as in constitution, it has chiefly its theologians to
thank for, for in its original idea this did not lie at all. That the Protestant Church, which
strove for the right of consciences, so quickly became a church of compulsion that tyrannized
over consciences, it has
% --- p. 120 ---
chiefly its theologians to thank for; for it was especially they who by their systematic concepts
of doctrine made the subjective matter systematically objective.

``But'' --- one says --- ``these blunders were long since abolished. What crying injustice to
blame orthodox theologians for what plainly lay in the spirit of the age!'' --- Let me not be
misunderstood. I do not judge consciences; I speak only of the principle. As it ought to be
remembered with gratitude that there were many orthodox theologians who in sincerity's strength
pleaded the cause of Christian freedom in the face of the threatening state church, many orthodox
theologians who in that sorry period of compulsion dared to repeat with Melanchthon:
\textit{Tyrannis est inimica ecclesiæ}; so it ought also to be dutifully acknowledged that there
are to this very day orthodox theologians, theologians who in sincerity's strength defend the old
strict concept of doctrine in its whole severity in spite of modern enlightenment. At a time when
each and all plead freedom's cause, it is doubtless a grateful task to speak up for the Church's
freedom too; but it is no grateful task at all to maintain the strict dogma and teach the
present age's sensible people that one is to believe --- against understanding.

Here, then, there can be no talk at all of the theologians' subjective conviction (for this would
precisely be, once again, to transform a subjective matter into an objective one); but here there
shall and ought to be talk of the theological principle. In principle, the objective dogmatics
that will help faith along by scientific grounds commits exactly the same blunder as the
objective state church that will help faith along by civil compulsion: they fit together so
exactly that he who will consistently hold to the first ought also to be consistent enough to
hold to the last. If the Christian dogma is in actuality so objective that in relation to science
it can be defended with objective grounds, then
% --- p. 121 ---
the state church is truly also so objective that in relation to religious upbringing it can and
ought to apply civil compulsion: then let us have the old church discipline back!

3) As we now in conclusion consider the dogmatic prejudices in their relation to higher
scientific rigor, we must first and foremost make clear to ourselves what it is that is really
here in dispute, and then also what is not here in dispute.

Here it is not in dispute whether it is a need for the Church to have its Christian dogmatics; but
here it is in dispute how the dogmas are treated in Christian dogmatics. Here it is not in
dispute whether the Christian propositions of faith are objective or not; but here it is in
dispute what the essential relation of the objective propositions of faith is to subjective
faith. Here it is not in dispute whether a thorough scientific knowledge is a necessary
presupposition for the systematic treatment of the dogmas; but here it is in dispute what the
relation of principle is in which science, as science, ought to stand to the system of dogmas.

That the Church must have its dogmatics --- insofar as by dogmatics one understands a thetically
coherent presentation of the Christian dogmas --- follows directly from its historical
presuppositions and its reciprocal relation to the age's cultivation. In the present time, when
the Christian propositions of faith have long since been explained away, construed, and
misconstrued in so many ways that Christendom has almost forgotten what the Christian is, the
congregation's teachers can undeniably profit from getting to see the Church's original dogmas
once again in their right shape, and from receiving an impression of the imposing distinctiveness
of the ecclesial structure of doctrine.

But the dogmatics that the Church's teachers thus need is assuredly not a half-speculative
dogmatics, a dogmatics that must first be mediated with the age's ideas, and on that account must
first be produced by some dogmatic genius or other. No, the
% --- p. 122 ---
dogmatics the Church needs is none other than the dogmatics the Church itself already has,
insofar as it already has its dogmas. The whole labor then simply confines itself to the task of
revising these dogmas and freeing them from speculations of every irrelevant kind.

But, it goes without saying, with this meager task theology of course cannot be content; it
demands a more, it demands that something be achieved here for science as well. So there prevails
the dogmatic prejudice that every new dogmatician who writes a new dogmatics must treat the
dogmas in a new and original way. This prejudice, which has by no means diminished but rather
increased with theology's increasing scientific rigor, then in turn produces a remarkable
opposition between ecclesial and theological dogmatics. Ecclesial dogmatics is in constant being,
for Christianity is that which is, and the Christian has been found; theological dogmatics, on
the other hand, finds itself in an unceasing becoming: it constantly looks as though one could
not yet rightly find wherein the Christian actually consists, as though the true and complete
concept of doctrine could never rightly be finished. Ecclesial dogmatics is essentially one and
the very same, for the Church maintains itself in maintaining its type of doctrine; but
theological dogmatics, which mediates faith's content with science's ideas, is in unceasing
alteration, for here in the Protestant Church there come to be just as many dogmatics as there
are original dogmaticians. --- Either, then, Christian doctrine itself must be alterable, like
theology's alterable dogmatics, or else one must confess that theological dogmatics is
heterogeneous with Christian doctrine, as the alterable is heterogeneous with the unalterable.

``But'' --- one says --- ``there is nevertheless a third and higher thing, a medium in which the
alterable and the unalterable are bound into unity: this medium is the idea. If the idea is taken
away,
% --- p. 123 ---
dogmatics becomes spiritless; but the Church can surely not be served by having an
idea-forsaken --- that is, a spiritless --- dogmatics.''

It is the modern notion of the dogma's scientific perfectibility that haunts the prejudice I
shall now discuss. The theological dogmatician will of course not deny that the dogma is a
proposition of faith, and as a proposition of faith most immediately an object for practical
appropriation in faith; but then he asserts at the same time that the dogma is on the other hand
an objective proposition of truth, a proposition of truth that as such must become an object for
an ever deeper appropriation in objective knowing. As a proposition of faith the dogma is
unalterable: the Church itself has only one dogmatics; but as a scientific proposition the dogma
is at the same time alterable, insofar as it ferments in an ever-advancing process of knowledge,
and develops in a multiplicity of different forms of the idea: the Church's theology has many
dogmatics.

``It is therefore the same dogma that in the idea shows itself both as unalterable and as
alterable; it is the same dogma that strikes root in the heart and unfolds its blossom in
speculation; it is the same dogma that becomes an object for a double knowledge, a subjective one
in faith, an objective one in knowing; it is the same dogma and the same faith that goes through a
double development, a practical one in life, a theoretical one in the school.'' --- Herewith I
believe I have set forth the speculating dogmatician's ``fundamental view'' as sharply and
definitely as possible.

But if now the dogma that is an object for subjective knowledge in faith cannot in any way become
an object for objective knowledge in knowing without at the same time altering its nature and thus
ceasing to be the same dogma; if now the faith that is to develop practically in life ceases to be
the same faith when it is set forth theoretically in the school; if now the presupposition that
the Christian dogma
% --- p. 124 ---
is to be known objectively because it is itself objective shows itself to be a false
presupposition: what then?

``\emph{Christianity is the paradox!}'' --- sounds another voice --- from another side. --- The
paradox! This word strikes school theology with a stunning blow. The paradox! But what, then, is
the paradox? Well, it is precisely that which science cannot explain. Yet from science's not being
able to explain the paradox it does not yet follow that the paradox, as some suppose, is merely a
catchword for paradox-mongers, an arbitrary cue of an ``ecstatic'' poet-thinker. From science's not
being able to explain the paradox it does not yet follow that it should be impossible to explain
what the paradox is and wherein it consists. Paradox is a Greek word which, if that could serve
to recommend it, is even found in the New Testament.

That I may here in my own way give a short explanation, I say: the paradox expresses the Christian
object's repulsion of the knowing subject, and consists in the tension under which the believer
with inwardness's infinite passion holds fast the object of faith.

Of all other objects it holds that the objective will be known objectively, and that there must
therefore be an objective relation of attraction between the knowing subject and the object that
is to be known; only the object of revelation makes an exception, for here the relation is as
repellent as possible. All other objects --- those of history, of natural science, of mathematics,
of philosophy --- let themselves be appropriated in objective knowing; the knower can come into
direct understanding with these objects, can penetrate deeper and deeper into their essence; only
with the object of revelation can the knower not come into direct understanding: understanding is
in conversion, in continued conversion; understanding is in being turned about, in continued being
turned about; the understanding is paradoxical. All other objects will be known objectively,
\emph{because they are ob%
% --- p. 125 ---
% NOTE (from the transcription): the letterspacing stops at the page turn ---
% „fordi de ere ob-“ is spaced on p. 124, but the turnover syllable
% „jective;“ opens p. 125 in ordinary spacing. The emphasis is carried over the
% whole word in both files.
jective;} only revelation's object will, both before faith and in faith and after faith, be known
subjectively, \emph{although} it is objective.
% Only „uagtet“ / "although" is letterspaced; „det er objectivt“ that follows is
% set normally (checked at 400 dpi), and is left unemphasized here.
The miracle, for example, is an object, and yet I cannot relate myself objectively to the miracle
in objective knowing; but the paradox is that I must relate myself subjectively to the objectively
repellent miracle, and appropriate it in faith's inwardness.

In the dialectical writings of our most recent literature the paradoxical relation between faith
itself and faith's object is, so far as I can judge, presented with such assurance, such acuteness
and precision, that one can see exactly that these factors must in an absolute sense be determined
for one another: to faith as faith corresponds such an object, to such an object corresponds faith
in its highest potency as faith\footnote{Since faith and faith's object are both equally
paradoxical and consequently homogeneous, the relation between these homogeneous factors is again a
direct and not a paradoxical relation. From this it is explained that the believer who lives in all
simplicity does not notice the paradox.}.

It is not unknown to me that the theological consciousness takes umbrage at the paradox, and
regards it as a sign of a monstrous one-sidedness in the subjective. That I find in this quite a
singular agreement between the theological consciousness and my own consciousness I must on the
present occasion permit myself to remark, since it might perhaps cast light on one or another point
in this involved matter.

The theological consciousness now tells us, in the year 1850, that the dialectical writings to
which we here allude are doubtless exceedingly brilliant, piquant, and apt, but that the
``fundamental view'' nevertheless, when one looks a little more closely, betrays itself as a
striking one-sidedness: the same thing my consciousness told me in the year 1843. The theological
consciousness is offended at the paradox because it looks like holding out a while yet
% --- p. 126 ---
--- perhaps even until the year 1853\footnote{This date is admittedly given at random: yet it
might, if only one will let oneself be instructed by ``the signs of the times,'' perhaps in time
acquire a ``deeper'' significance for those who have ``deeper sense.''}; my consciousness became
offended at the paradox because it came so suddenly, so strangely, with ``\emph{Fear and
Trembling}'' in the year 1813.
% The Danish here reads 1813, an error for 1843 (cf. p. 125, „i Aaret 1843“, and
% the date of Frygt og Bæven). Verified at 400 dpi and kept as printed there.
% The English gives the intended 1843.
The theological consciousness still sees in the present year, 1850, as clearly as before how ``das
winzige Paradox'' lacks positive force, and therefore cannot hold its own against ``positive
theology''; remarkably enough, the same thing --- let me say it without self-praise --- the very
same thing my consciousness also saw clearly until shortly before the year 1848. I therefore need
--- what in such a retrospect is besides an easy matter --- only alter a single date, and I have
the fairest agreement between the theological consciousness and my own humble consciousness, so
closely are they akin.

What in later times may have so clouded my vision that I can no longer, as before, see the
paradox's ``evident one-sidedness'' those theologians who have still preserved clear sight will
surely know better than I. The only explanation I am able to give here is merely that from the year
1845 until now I have spent most of my time and diligence in studying --- besides the German
freethinkers and the theology dissolving itself in Germany with German consistency --- the writings
here alluded to. Whether I have in every respect understood these writings quite as the author
himself will have them understood I cannot possibly say --- one understands an existential author
only insofar as one understands oneself --- but studied them I have, and that with exceptional
interest. Who knows? perhaps this is precisely the fault! who knows? had I not studied them at all,
not guessed at these riddles, not lost myself in these regions, not stared
% --- p. 127 ---
at these sphinxes: who knows? perhaps my consciousness would then still have been, if not quite as
clear, then at any rate nearly as clear-eyed as the theological consciousness; they do say that
sharp light and much studying harm the eyes.

It is not unknown to me that the theological consciousness is dissatisfied with the paradox.
Doubtless the ``good forces'' that are now beginning to stir (especially in the very youngest
theology) would be able to knock ``the paradox'' over if it really had to be; but such a thing would
after all always have to happen in a conflict, and ``Christian wisdom'' --- prefers the blessing of
peace. Doubtless the theological consciousness is inwardly dissatisfied with the paradox, but ---
the paradox has once for all been put in; the man who put it in observed the precaution of putting
an existence in with it, and doing it so resolutely, so quickly, so deftly that theology noticed
nothing at all until it was done. Now it is done, and the theological consciousness sees very well
that what has happened has happened. Theology therefore lets the paradox stand for what it is,
untouchable and untouched, and is then firmly resolved to chastise everyone who might otherwise
presume to touch the untouchable.

The paradox is theology's stone of stumbling; if theology stumbles over this stone of stumbling, it
must fall and smash its dogmatics to pieces.

I have been reproached for changing my standpoint so suddenly --- for the paradox's sake; this
reproach is undeserved. If I were to reproach myself for anything, it would surely have to be that I
hesitated too long, and was not resolute enough to take my side at once and declare it plainly.

I have been reproached for attacking ``so personally'' one of our best theologians --- for the
paradox's sake; as though I thereby wished partly to show the world what a sacrifice I was making
for truth in thus changing my standpoint, partly to seize
% --- p. 128 ---
the opportunity of winning --- even if only at second hand --- a little victory over an otherwise
talented man whom we must all nevertheless highly esteem as a distinguished and deserving teacher.

The first of these accusations I pass lightly by, since I know how groundless it is. I have in this
matter made truth no sacrifice worth mentioning. I stood, in accordance with my position, outside
theology, and could therefore break with theology without any particular self-denial; I bore no
opinion, was not borne by any opinion, and consequently was not made anxious by any finite regard
for an opinion.

As regards the last surmise, on the other hand, it would grieve me deeply if it should really spread
and win credence; for I can in this matter win no victory whatever, either at first or at second
hand; nor have I ever felt any urge to win any victory whatsoever, either great or small, if I
except only the victory I, \textit{volente Deo}, am to win over myself. Before I began my so-called
``conflict,'' I first bound my hands: they are bound, I bound them myself, and cannot loose them
again. So I contend as best I can, with bound hands; but he who contends with bound hands is surely
a fool if he contends for an outward victory. I do not intend to stand in anyone's way, and least
of all to offend a man for whom I have personal goodwill; I wish only to make use of a permission
that I give myself, and that surely no one can deny me: a permission to express my conviction, a
permission to shed light on theological science, a permission to repeat my \textit{cæterum censeo},
\emph{that Christianity is higher than science, and the faith of the Gospel heterogeneous with
theology.}

%% ============================================================
%%  TENTH LECTURE (pp. 129--143)
%%  NB: head p. 129 reads "giældende" / "Theorien, thi"; the Indhold reads
%%  "gjældende" / "Theorien; thi". Real variant --- both witnesses kept.
%% ============================================================
\chapter*{Tenth Lecture}
\addcontentsline{toc}{chapter}{Tenth Lecture}
\label{ch:forel10}
\markboth{Tenth Lecture}{}

\begin{center}\itshape
According to the method now in force, theoretical theology finds itself in
an evident disproportion to religious practice; in order to maintain the
principle of faith, the formation of pastors must emancipate itself from
theory, for the faith of the Gospel is heterogeneous with theology.
\end{center}

% --- p. 129 ---
It is well known how the Greek philosophers asked whether virtue was something that could be
taught or something that could not be taught. When by exact examination they had found out that
virtue could indeed in one respect not be taught, but in another respect could be, they then
asked in turn how and under what conditions that in virtue which could be taught ought also
reasonably to be taught.

If it is already a great question whether virtue can be taught at all, and how it ought to be
taught, then the question whether and how far Christianity can be taught by human instruction,
and how it then properly ought to be taught, is assuredly a still far greater and more difficult
question. Recent theology is now so occupied with Easter reckonings and critical introductions,
and on that account has doubtless only little time for a closer discussion of this question; but
the question nevertheless retains its significance when what is at stake is the relation of
principle between Christianity and theology.

% --- p. 130 ---
Theology divides itself into the theoretical and the practical. Theoretical theology first forms
the pastors-to-be into scientific theologians; practical theology is then next to form the young
theologians into actual pastors.

If we now place our question in a definite relation to this two-sided task of theology, then
practical theology undeniably seems to give us a quite different answer from the theoretical.

If Christianity were something that could be taught by human instruction, if it were possible by
human art to form and educate human beings to understand Christianity --- that is, to understand
how under all life's circumstances they are to feel and think, be silent and speak, suffer and
act, as Christians; indeed to understand how they, after having themselves learned Christianity,
ought now also to teach it in turn to others: if it were possible, I say, that Christianity could
thus be taught to the last dot and tittle, who then was nearer to learning it thus than precisely
the congregation's future teachers, the pastors-to-be? and if Christianity in this way let itself
be taught, who then was nearer to taking charge of the teaching and the education than precisely
that theology which has after all set itself the practical task of forming and educating the
congregation's teachers.

If practical theology now really supposed in earnest that Christianity was something one could
learn by thorough studies and methodical exercises: should one then not expect that this theology
with all possible force and zeal took charge of the pastors' formation, and in every possible way
endeavored to maintain its standing and answer to its fair name: pastoral theology.

But that practical theology should really have such high thoughts of itself is hardly probable
--- at least not here in this country, where it is content with so exceedingly modest a role. For
what, pray, are the practical exercises at the pastoral seminary? even
% --- p. 131 ---
if one would count in the trial sermon, the catechetical test, and the bishop's examination, how
utterly trifling it all is beside the great scientific forces that theoretical theology is able
to sacrifice on its ``Easter reckonings'' and ``images of Christ'' and ``critical introductions.''

When, then, pastoral theology, which after all has most immediately to do with the Christian,
withdraws so modestly, and on that account, far from plaguing its candidates with many studies
and methodical exercises, rather commends them to their own subjective inventiveness, so that
they must now go about and try themselves practically on their own to see whether they may here
or there hit upon the Christian: must we not then draw from this the conclusion that Christianity
is something that in pastoral theology's opinion really cannot be taught at all --- just as
Socrates in his day drew the conclusion that virtue must surely be something that could not be
taught at all, when he saw how Pericles, who was himself so virtuous, neglected his sons'
instruction and let them ``go about like those animals that were consecrated to the gods, and
graze where they would, until they might perhaps of themselves hit upon virtue''?

If we turn with our question to theoretical theology, do we not then come to the entirely
opposite result --- namely that Christianity is at bottom a learned matter, a matter that both
can and ought to be taught, and taught thoroughly? --- In the Protestant Church, where ordination
has ceased to be a sacrament, one has then particularly felt a need to inculcate the ecclesial
necessity of ``Christian science,'' of theological learning, because theological learning is
really the ideal mark that distinguishes the pastor from the layman. The pastor is a learned
Christian, the layman an unlearned Christian --- behold, that is the difference, a difference that
holds as a rule, although the rule has its exceptions.

% --- p. 132 ---
In the state church's proper period of authority, when there was zeal for objective orthodoxy,
the difference between learned and unlearned Christians was in an ecclesial respect a very
considerable difference. What mattered was to get a subtle concept of doctrine, articulated at
all points, developed in scholastic forms with scholastic consequences, and to purge out heresy.
In order rightly to distinguish between true and false faith, the pastor necessarily had to be a
theologian, and learned theology had to have the first voice in all ecclesial matters of faith.
In later times theology has doubtless abandoned that strict orthodoxy in order to engage with a
more sensible wisdom and keep step with advancing enlightenment; but insofar as theoretical
theology has to that extent lost a part of its ecclesial standing, it feels itself all the more
vigorously called upon to convince the state of its scientific standing, and to hold fast to the
``fundamental view'' that Christianity is a theory.

Theoretical theology doubtless grants that Christianity is a life just as much as a theory; but
since the theoretical is and remains the most important thing in a theory, it is quite natural
that theoretical theology holds to the theory. The proof of this we have in the great, decisive
--- indeed for a whole life decisive --- emphasis that is to this very day laid on the so-called
theological attestation, the pastors-to-be's theoretical examination. Or how should one explain to
oneself the singular sight that the whole yield of several years' laborious studies could be
tested so reliably in a few but significant moments, if theology were not in this very examination
best able to gather the theoretical rays of the Christian light as into a single theoretical focal
point? And how should one explain to oneself the singular sight that a scientific examination for
the ministry, confined to a few hours, could still preserve --- not indeed a sacramental, but
assuredly so fundamental a significance in the state's eyes that it became decisive for the
candidate's
% --- p. 133 ---
whole ecclesial-civil future: how, I say, was this possible if the state itself did not
theoretically see that Christianity is a theory?

Let one judge of this as one will: so much is certain, that the formation of pastors among us is
essentially a theoretical formation, in that it is exclusively a theological formation. One
proceeds without further ado from the presupposition that a theologian can form a pastor; with
this presupposition one then restricts practical theology and expands the theoretical: it is a
remarkable procedure. Can a theologian, then, form a pastor? Practical theology shakes its head
doubtfully and withdraws; theoretical theology steps boldly forward and takes over the whole
formation. It looks somewhat strange! Can a theologian, then, form a pastor?

Had Socrates lived, he would have seized on this question; he would surely have gone about and
inquired --- now of theologians, now of pastors, now of students, now of candidates (as occasion
offered) --- whether there was no one who would kindly tell him how it really comes about that a
theologian can form a pastor.

That like can in general be formed by like, the homogeneous by the homogeneous --- that, for
example, a mathematician can be formed by a mathematician, a physician by a physician, a jurist
by a jurist, a theologian by a theologian --- would at once have been quite evident to Socrates,
at least as evident as that like can be born of like, a human being of a human being, a sheep of
a sheep, and so on. But how like can form unlike, and namely how a theologian can form a pastor,
that would surely have been as troublesome for that questioning dialectician to get into his head
as how the entirely like can be born of its unlike, the homogeneous of its heterogeneous, a human
being of a sheep, a sheep of a human being.

% --- p. 134 ---
It was in a way well that the old questioner died when he died! had he still lived, he would have
caused ``Christian science'' much bother; had he still lived, he would have caused ``Christian
wisdom'' many an anxious hour! for he would after all not have been able to refrain from
instituting a hearing, he would after all not have been able to keep himself from plaguing and
harassing the theologians, the candidates, the pastors, indeed perhaps the deans, until he had at
length got the secret out of them, the great secret,
% Letterspaced in the Danish print (Sperrsatz), beginning with „hvorledes“ at
% the end of the preceding line and running to the full stop.
\emph{how a theologian can form a pastor.}

It is not necessary, however, that a Socrates rise from the dead to tell us how much lies in this
question; everyone can tell it to himself once he has looked about him a little in a theological
lecture room.

That the relation between teachers and hearers cannot be the same in theological science as, for
example, in medical science strikes the eye at once.

In medical science theory and practice are most closely bound together. He who is skilled in
medicine and communicates the theory is as a rule himself a practicing physician; the hearer who
grasps the theory soon notices that in order himself one day to be able to practice, he must
first of all form himself by the theory. Nor does the medically skilled teacher meet his hearers
with the explanation that the art of medicine is something that properly speaking cannot be
taught at all, but that one will nevertheless lecture both on the theories and on the critical
introductions to the theories, because one must after all hold to medical science's necessity and
wishes to have scientifically formed physicians. No, the medically skilled teacher lays out his
lecture from the start in such a way that the hearer at every point notices that what is here
communicated in theory is also something that can and shall be used in practice. The student who
thus becomes a partaker in medical-scientific formation will
% --- p. 135 ---
then also soon notice in himself whether he is at all suited to be a physician, or whether he has
perhaps mistaken his vocation --- in which case he can withdraw in time and choose himself another
course of study. Finally, the intimate connection of theory and practice that distinguishes
medical science has this good in it too, that the young physician's civil future is not absolutely
settled by academic examinations. The young physician for whom these examinations have in one or
another respect had a less happy outcome can, if he feels strength and courage, still procure
himself redress; by continued study and continued practice he can develop his talent and open a
way for himself that tolerably answers to his individual ability.

How quite otherwise does it not look with theological science's relation to pastoral formation.
The theories lectured in the theological halls are now so scientific, so learned, so objective
that on that account they must of course keep as far from practice as is well possible. What is
the consequence? That the formation begins in ambiguity, and that this ambiguity then in turn
prevents the young hearers from understanding themselves. Here now, for example, sit two hearers
side by side: they have both chosen the same course of study, and yet --- how different! The one
has brought with him a living sense --- not indeed for learned science (for his scientific
abilities are perhaps somewhat weak) --- but he came with a living sense for the religious; he
chose theology, he chose it with a good hope, for he so much wanted to be a pastor. He has
reverence for theology, he listens attentively to the theologian's lecture, and as he listens ---
the learned material grows; he gives himself over with a certain love to theology, and now expects
that the many explanations will at length explain to him what revelation is, what faith is, what
love is: but it seems to him almost
% --- p. 136 ---
as if the many explanations were driving clarity away, for the theories contradict one another.
Still he will not lose courage, but will work at the material and study his theology, for --- he
is studying to be a pastor. Learned criticism weighs on his memory and confuses his concepts; yet
he will not yet give up hope, but will diligently study his theology, for he is studying --- to be
a pastor. But --- if the theories nevertheless turn against him, humble him, reject him, and judge
him on the day of judgment; if criticism makes it manifest that despite all his theological
studying he will nevertheless never become a proper theologian, alas! then he has wasted his time
and diligence, then he will surely never become a pastor either!

The other, on the other hand, brings with him no particular sense for religion as religion, but he
brings a quick, sound, and vigorous head, a faithful memory, a sharp understanding, an eye for the
objective idea! --- So he has presumably not come in order to study for the ministry, but in order
to form himself into a scientific theologian? Do not disturb him with irrelevant questions! Of
science's ideals he has not yet thought, any more than of the clergy's duties; but --- he listens
with pleasure, he reads with pleasure, he will in 3 years take his attestation with great
pleasure. When that gust is over, the future stands open to him, and only then can there be any
talk of recognizing his calling and choosing his vocation. If there is then a prospect of a
theological teaching post, he can easily go the theoretical way; but if it turns out otherwise, if
he is called to a practical-clerical activity, he will have nothing against going over into the
Church and settling into --- a good living.

We will for the rest readily grant that the majority of theology's young devotees combine a living
religious sense with a quick scientific interest; we will at the same time grant that the extent of
the mass of knowledge they have to work through is not
% --- p. 137 ---
greater either than what a tolerably gifted student can manage when he uses his time as he ought;
but with this the disproportion between means and end is by no means justified.

When one rightly considers what it costs before a young man brings it so far that he passes his
theological examination, and passes with honor: then one is surely also entitled to ask what yield
these preliminary studies really give, and how far an academic preparation that has cost the family
and the state so considerable a sacrifice now really benefits both the candidate and the state.

Our candidate is then a cultivated and knowledgeable young man, not merely cultivated in general,
but in possession of a quite particular cultivation. When we now therefore ask about the yield that
this cultivation gives, we are not asking in general about the general advantage that knowledge and
cultivation always give; but we ask with the candidate, and he with us: how shall the so dearly
bought theory now find in actuality its corresponding practice?

Our candidate is formed in the theologians' school. In the theologians' school he heard, then, the
exegete's lectures on exegesis, the church historian's on church history, the dogmatician's on
dogmatics, the critic's on criticism, and so on. And the exegete formed him after his own image
into an exegete in miniature, the church historian after his own image into a church historian in
miniature, the dogmatician after his own image into a dogmatician in miniature --- in short, the
theologians joined in forming him into a whole theologian in miniature. For it is certain that a
theologian can form a theologian.

The practice that directly corresponds to the theory thus won would now directly be this: that the
young theologian, the theologian in miniature, should be further advanced in theological insight
and theological learning, and then, when he had himself become
% --- p. 138 ---
a proper theologian, a theologian on the larger scale, should in turn work at forming others into
theologians. Behold, that is the peculiar station in life that corresponds wholly and entirely to
our candidate's peculiar formation.

But --- we ask --- is it then for such a station in life that theology's young devotees as a rule
prepare themselves? By no means. Theological teaching posts are with us but few. Of 400 students
who all study theology, there is perhaps one who goes through and becomes a theologian; the 399
seek other appointment.

But when our candidate cannot in the stricter sense become a theologian, what then? Well, then he
shall assuredly learn how difficult it is to find a practice exactly corresponding to his theory. In
the Church he does not find it; for the congregation has no need whatever to hear popular lectures
either on church history in miniature, or dogmatics in miniature, or the critical introduction in
miniature. If the candidate is to speak in the congregation, he must begin over again; if he is to
speak the Church's language out of the Church's own train of thought, he must begin over again, and
in beginning over again he has --- much to forget, but little to remember.

Nor does his peculiar theory find any peculiar practice anywhere in the cultivated world. The
questions in whose answering the young candidate has his real and true strength are not the
cultivated world's questions; the leading thoughts with which he is in a particular degree familiar
are not the cultivated world's leading thoughts. Our young candidate soon notices that he must
constantly convert his peculiar theory into the age's general cultivation,
% The Danish here sets „Dannelse,og“ with no space after the comma --- printer's
% error, normalised in the transcription and set correctly here.
and that to that extent he is wronged, since at no point can he come to operate with the powers
peculiar to him. He then easily comes into a precarious position; for he comes only too easily into
the unpleasant position that he
% --- p. 139 ---
must straightforwardly throw away his theory in order to find his practice: for the poet's words:

% The two German lines are set in antiqua against the surrounding Fraktur in the
% Danish, but upright, not italic; set as a plain quotation, as with the Schiller
% couplet on printed p. 90.
\begin{quote}
Was man nicht weiß, das eben brauchte man,\\
Und was man weiß, kann man nicht brauchen,
\end{quote}

fit as though they were written for a theological candidate.

Many will perhaps think this depiction as one-sided as it is rash, since experience daily shows us
that theological formation benefits the candidate in many ways. I imagine, then, that in favor of
this view one might put forward the following:

``In order to work as a teacher of religion, whether it be in the church or in the school, the
candidate must be in possession of the particular formation that enables him to master his material
and arrange his discourse according to the hearers' need. The teacher of religion ought, like every
other teacher, always to possess far more than he can immediately communicate --- if for no other
reason, then because only on this condition can he possess what he communicates in a higher and
freer way. It is an entirely unfair demand that the spiritual yield the candidate has from his
learned formation should be demonstrable piece by piece in such a way that it could, like a material
yield, a pecuniary advantage, be pettily reckoned up in rixdollars, marks, and shillings. One need
only note what difference there is between the religious instruction a good seminarist can furnish
and the instruction the confirmands receive from the able pastor; and one will without doubt grant
that the higher theological formation has its significance.''

``Doubtless a theological student learns much that he must afterward forget; but even if he can
perhaps forget a part of the many particulars, he nevertheless preserves a scientific impression of
the whole, and with this impression at the same time a freedom of mind with which he can all the
more surely judge the religious
% --- p. 140 ---
life's various phenomena and find his way in practical circumstances.''

``Doubtless one must grant that theoretical theology among us has at present a certain preponderance
over the practical. But partly it is a one-sidedness that can with time be remedied; partly it is on
the whole a questionable matter to restrict scientific formation too much. For one will not be able
to deny that the practical formation communicated in clerical seminaries only too easily acquires a
subjective, sickly, pietistic stamp when it lacks a firm scientific foundation.''

``Doubtless scientific theology cannot form the clergyman so immediately and directly as the
practical school; but it nevertheless forms him mediately and indirectly, it forms him objectively,
and judges his formation by a fixed objective rule. This is especially of importance with respect to
the state. For the state, which neither can nor may engage in any subjective assessment of
individuals' subjective zeal for faith, can and ought on the other hand to have a fixed objective
rule for filling the clerical offices.''

I have thought it right not to overlook such considerations, which always look best when they are
put forward in pure generality, and can therefore easily take it into their heads to present
themselves as serious admonitions. If, however, we go into the matter itself, such considerations
will hardly be able to uphold the theory's standing. One appeals to the necessity of a higher
scientific formation. The significance of higher formation is here not at all called in question;
but the question is whether such a higher formation could not be set in a far more intimate
connection with the corresponding station in life than has hitherto happened. One appeals to the
necessity of an objective measure of the candidate's ability. The significance of such a measure is
here not at all denied; but the question is whether the practical
% --- p. 141 ---
school were not able to furnish a fully as just a measure, if only one could get it laid out
according to a plan answering to the idea, and become conscious of wherein the higher practical
formation then really consists.

That the state has a measure --- indeed a scale divided with mathematical exactness --- in the
qualifying examination now existing shall never be denied. But it is not merely a matter of having
an exactly divided measure; it is at the same time a matter of having a measure that tolerably
answers to the nature of the things that are to be measured with it. --- If what is to be measured
is far too heterogeneous with the measure, then it is of little use that the measure itself is
exactly divided, like a graduated scale in inches and lines.

As regards the state, it can with respect to ecclesial affairs reasonably will nothing other than
that the Church's servants shall in the best way look after the Church's needs. The question is
therefore this: how far the theoretical preparation can now really make the candidate fit to take
over a clerical office and in this his office look after the true needs of the Church and the
congregation?

To discuss this question is a prolix matter; only a couple of words more before I conclude!

As a defense of theological dogmatics we have already heard that it is
% The Danish here reads „allerde“ (a letter dropped from „allerede“) ---
% printer's error, kept as printed there. The English has the correct word.
to stand in the most intimate connection with evangelical preaching: ``to such a sermon corresponds
such a dogmatics, to such a dogmatics such a sermon.''

If we now hold to this theological presupposition, and add what is likewise undeniable, that
theological science has many different, mutually contradictory dogmatics, then it follows that the
congregation must gradually also get many different, mutually contradictory sermons to hear.

This is then also in practice a well-known matter. Every preach%
% --- p. 142 ---
er as a rule preaches, if not exactly according to a definite dogmatics, then at any rate according
to the theological tendency to which his academic formation belongs: the orthodox orthodoxly, the
rationalist rationalistically, and so on. --- In this way the ecclesial original Word is exposed to
all the alterations that human knowing and human theories undergo; in this way the age's spirits
force their way into the congregations, and make themselves little by little
% The Danish here sets „Menighedeine“ --- a wrong sort (a dotted i where an r
% belongs) for „Menighederne“ --- printer's error, kept as printed there. The
% English has the correct word.
masters over revelation's Word.

We praise Luther and the Reformers for having cleansed the Gospel of the Middle Ages' manifold
human, arbitrary additions (\textit{traditiones humanæ}); but what shall we now judge of the many
mutually overturning dogmatic theories? do these not also contain \textit{traditiones humanæ}? But
if the evangelical Church has been freed from superstition's medieval corruptions only to fall
forfeit to the scientific theories' manifold modern corruptions: how much is then, in a religious
sense, gained by the exchange?

One does indeed now and then see how theoretical theology begins to pull itself together, and we
must really acknowledge that in many ways, among us too, it strives to appropriate the Christian
content and make the doctrine fruitful for the Church and the congregation. Yet the endeavor will
not rightly succeed; ``Christian science'' is lawfully excused: one cannot demand of the theory what
exceeds the theory's powers. --- Every Christian theory that keeps within objective knowing will
further be dissolved into a half-Christian one, and this in turn into an un-Christian one; and that
for the simple reason that Christian theory is a contradiction.

Insofar as preaching goes on conforming to theory and theology, it cannot possibly withstand the
spirit of the age. The weakening already showed itself in the last half of the previous century,
when the newer enlightenment broke in over the old state-church dams. Instead of defying the modern
consciousness
% --- p. 143 ---
and maintaining Christianity's Christian character, the pastors seized on the expedient of abating
something in original Christianity's strict demand, so that
% The „at“ before the quotation is set with a worn t in the Danish: the crossbar
% has not printed, so the sort reads „ai“ on the page. Read as „at“ there.
``the enlightened'' should not entirely disdain religion and wholly abandon the Church. With this
accommodation Christianity lost its force, in losing its sting; but the theory of accommodation was
itself an invention of theology.

In a religious sense Christianity is the absolute; where it cannot and may not be the absolute, it is
not at all. In a religious sense faith is the absolute; where it cannot and may not be the absolute,
it is not at all. Here, then, we stand again at the result
% Letterspaced in the Danish print (Sperrsatz) to the end of the lecture.
\emph{that Christianity is higher than science, and the faith of the Gospel heterogeneous with
theology.}

%% ============================================================
%%  ELEVENTH LECTURE (pp. 144--158)
%% ============================================================
\chapter*{Eleventh Lecture}
\addcontentsline{toc}{chapter}{Eleventh Lecture}
\label{ch:forel11}
\markboth{Eleventh Lecture}{}

\begin{center}\itshape
Theoretical theology must relate itself skeptically to Christianity's
revealed truth; practical theology, which in imagined existence labors to
solve the Christian tasks, must end by dissolving itself: the faith of the
Gospel is heterogeneous with theology.
\end{center}

% --- p. 144 ---
That a religious science is a contradiction, that a science of revelation is a contradiction; that
therefore any theology whatsoever (be it Muhammadan, Jewish, or Christian) is in this respect a
contradiction, I have in the foregoing striven to show, in that from several sides I have endeavored
to illuminate the opposition between religious knowledge and objective knowing.

% The scan carries a former reader's pencil underlining under the corresponding
% Danish clause here. It is a modern hand, not letterspacing; the type
% underneath is ordinary Fraktur. No emphasis in either file.
But although in these lectures I have thus expressly protested against theology's Christian
scientific rigor and scientific Christianity, I have nevertheless --- be it noted --- protested
neither against an objective science about Christianity nor against a religious knowledge of
Christianity. If one will now comprehend these two different kinds of knowledge under one
designation, and use for it the ancient name hitherto employed, theology, I shall have nothing
against it: about mere names we will not quarrel.

Theology will therefore, according to the view here set up, not be abolished; but it will alter its
nature and its
% --- p. 145 ---
relation to Christianity. --- Theology thus altered will then again divide itself into two main
parts: the theoretical and the practical.

Theoretical theology we then characterize as science in the sense that it is an objective knowing
about Christianity's facticity. This knowing about Christianity's facticity must be carefully
distinguished from the so-called objective knowing of Christianity's truth. The fundamental error in
the prevailing theology is, in my conviction, precisely this: that one throughout confuses the
latter with the former.

With the false presupposition that in faith's name there is to be an objective knowing of
Christianity's objective truth, theology has hitherto placed itself in a positive, affirming
relation to Christianity itself, and undertaken its defense by objective grounds. The insufficiency
of this defense emancipated reflection has discovered, and therewith the modern prejudice has forced
its way in that the freethinker could make an end of Christianity because he can in a certain sense
make an end of theology.

If, on the other hand, we conceive theoretical theology as an objective knowing about Christianity's
facticity, the point of view will instantly alter. Theological science will then no longer place
itself in a directly affirming relation to Christian truth, but only in a direct relation to the
Christian phenomenon. The theoretical task will therefore become a double one. In the first place
science must strive to make clear to itself what the Christian is, whereby it is distinguished from
the non-Christian --- in general, how original Christianity appears to objective contemplation. Next,
science must in relation to the Christian phenomenon thus clarified give an account of what proofs
and objective grounds there are for the phenomenon's validity. --- As science now comes out with
% --- p. 146 ---
% „imod“ and „for“ are letterspaced in the Danish print (Sperrsatz); verified at
% 400 dpi against the unspaced „imod“/„for“ two lines further down.
all its grounds, the result will show itself that there are exactly as many grounds that speak
\emph{against} as there are grounds that speak \emph{for} Christianity. Only when science sees this
is it impartial, and only when it is impartial is it what it ought to be, a true objective science.

Accordingly: the science that from its investigations gets the result that there are more grounds
speaking for Christianity than against it is \emph{one-sided}; and likewise the science that supposes
it can exhibit more and weightier grounds against Christianity than for it is \emph{one-sided}. True
science's true
% „non liquet“ is set in antiqua against the surrounding Fraktur.
result is here a negative result, a \textit{non liquet}.

Objective science is here to be presented under the image of objective Justice holding the balance in
her hand: in the one scale are all the positive, in the other all the negative grounds: so the scales
are held in equilibrium. In this equilibrium there is no decision; for the decision to come about,
something else must be added: this something else is --- either faith or unbelief. If I now take
unbelief and lay it on the negative scale, this at once sinks, and the positive grounds lose their
weight; if conversely I lay faith's infinite weight in the positive scale, then all the negative
grounds are outweighed, and as easily as if they were not there at all. As Justice, who holds the
balance, has her eyes bound so that she shall not be partial to either of the contending parties, so
science too has in this respect its eyes bound, so that it shall not be partial to faith or unbelief.

That this is now not merely an image but at the same time a truth I believe I can establish by a
consideration of theology's various disciplines. To avoid prolixity I confine myself to three:
isagogics, hermeneutics, dogmatics.

% The name of each discipline is letterspaced in the Danish where it heads its
% paragraph (here „Isagogiken“, p. 147 „Hermeneutiken“, p. 149 „Dogmatiken“),
% but not in the list just above --- verified at 400 dpi.
\emph{Isagogics}, free scientific biblical criticism, theology ought not to give up, but should
rather work at carrying it
% --- p. 147 ---
through, so that it may reach the final and purest result. But everything now depends on what notion
one is to form of the final and purest result. --- So long as one persists in the prejudice that the
result can in any way become decisive for faith or unbelief, biblical criticism will never be able to
raise itself to a scientific impartiality: the religious man will then always be too sparing, the
freethinker too unsparing. There is surely nothing that in recent times has so inflamed freethinking
as the ambiguity that biblical criticism, according to theology's own confession, was to be a free
science, and yet could never become one so long as every theologian went on mixing his bit of faith
into the purely scientific proceedings. If, on the other hand, we proceed from the presupposition
that biblical criticism neither can nor shall settle anything for faith, but rather has the task of
bringing it to the highest possible equilibrium of proofs of probability \emph{for and against}, then
science becomes impartial and therewith at the same time thorough.

\emph{Hermeneutics}, the science of free interpretation, cannot be halted and ought not to be either.
At every new stage that the science of interpretation traverses, the human spirit wins ever new
scientific clarity about itself and its general relation to Holy Scripture. But everything again
depends on what notion one forms of biblical interpretation's final and purest result. --- If one
imagines that scientific biblical interpretation will at length get so far that it will be able to
settle how far the inspired Word is divine and how far it is human, then this science will never
become impartial, and to that extent never thorough either. The religious man will always be too
sparing, the freethinker too unsparing: the impersonal relation will at one stroke become personal,
the scientific dispute unscientific.

A scientific theory of inspiration is a contradiction. Every%
% --- p. 148 ---
one who attempts to develop such a theory is by the force of the consequences continually thrown from
the one extreme to the other, and thus tossed about as between Scylla and Charybdis.

Inspiration in the whole is an empty turn of phrase; inspiration in the particular is intolerable to
science. If I begin by striking out the particulars, then it ends with there being --- in the whole
--- nothing at all left. If, for example, I cannot believe the miracle of the Gadarene demoniacs,
then neither can I believe the miracle of Peter's draught of fishes; then neither Christ's walking on
the sea; then neither Christ's transfiguration on the mountain; then neither the angels' appearance
at Christ's birth; then neither the angels' appearance at Christ's grave; then neither Christ's
visible ascension; but let us now be a little consistent --- then surely neither Christ's resurrection
either; and then finally --- for consistency's sake --- neither Christianity.

If conversely I begin by positing a single narrative as inspired and infallible, then the one draws
the other after it, and in such a way that the whole Bible with all its wonders at last becomes
inspired. If, for example, I assume that Christ is risen by a miracle, then I must surely also assume
that the angels may have appeared to the women by a miracle; then also that they may have appeared to
the shepherds by a miracle; then also Christ's transfiguration on the mountain by a miracle; then
also his walking on the sea by a miracle; then also the miracle of Peter's draught of fishes; then
also the miracle of the Gadarene demoniacs; then also --- to be a little consistent --- the miracle
of Balaam's ass; then finally also --- to be entirely consistent --- every single miracle narrative
in the whole Old Testament.

Here in these consequences of the theory of inspiration there is a fluctuation, a surf, a breaking of
waves more terrible than
% --- p. 149 ---
the breaking of waves at the stormy headland. In such a breaking of waves it is entirely impossible to
steer scientifically by science's compass; for the expedient of accepting some miracles, construing
others differently, and rejecting others is nothing but an unscientific evasion, an evasion by which
one nevertheless at last makes everything depend on a subjective judgment; but in subjective judgment
science is dissolved.

As Holy Scripture the Bible sets itself in the most definite terms against every objective theory of
inspiration; for as Holy Scripture it is only for the believing subjectivity. Scientific
interpretation must consequently end with there being just as many grounds speaking against
inspiration as there are grounds speaking for it; that is, in other words: science ends in
skepticism, and can settle nothing whatever about the Bible's inspiration.

\emph{Dogmatics}, as objective theological science, ought not to be abolished, but should rather be
developed consistently to the point where the true relation between faith and knowing can show itself
clearly and definitely. If one now conceives the relation in such a way that faith can in any way be
clarified, grounded, and confirmed in knowing, then one is also compelled to grant that faith can
likewise be darkened, confused, and weakened in knowing; if one grants that knowing can in any way be
clarified, confirmed, and grounded in faith, then one must surely also grant that knowing itself can
be darkened, confused, and weakened in faith. The believer then gets in this way a believing knowing,
the unbeliever an unbelieving knowing, the half-believer a half-believing knowing. Knowing has then in
this way lost its objective character; objective knowing courts subjective feeling, knowing enters
second childhood: science is dissolved.

If science is to treat the dogmas scientifically, then it will surely show itself that every dogma has
just as many repelling as attracting grounds, that at every dogma there stands
% --- p. 150 ---
``a cross for thought,'' that scientific reflection relates itself skeptically to faith.

The manner in which the scientific testing of the dogmas is undertaken can for the rest be various;
but this stands fast, that the question of the dogma's comprehensibility and incomprehensibility must
be settled impartially --- that is, without any subjective regard whatever to subjective conviction.
--- I am permitted to believe the dogma although I do not comprehend it --- indeed that is precisely
% „skal“ and „kan“ are letterspaced here in the Danish; the same words are
% unspaced in the next sentence, which settles it.
religion's demand, that I \emph{shall} believe because I \emph{cannot} comprehend ---; on the other
hand I am not at all permitted to act as though I could nearly comprehend a dogma because I should so
much like to believe it.

When dogmatics develops with scientific consistency, it passes over into the philosophy of religion;
but the philosophy of religion's final and purest result is, in relation to revelation's fact, to be
expressed thus:
% Letterspaced to the end of the sentence.
\emph{Science comprehends that it cannot comprehend.}

Thus all theoretical theology's disciplines now stand in the same negative, skeptical relation to
Christianity's truth. All the same, this negative relation cannot be settled once and for all. Science
does not let itself be dismissed by immediate protest; the not-knowing in which science is here to end
must each time be attained by knowing's way. Theoretical theology can therefore to that extent be said
to find itself in a constant advance, in that each time it runs anew the road from presupposed knowing
to not-knowing, it at the same time rejuvenates itself and wins renewed clarity about itself and its
relation to revelation's phenomenon.

As regards next practical theology, it is at once seen that it lies in quite another sphere than the
theoretical. Theoretical theology is a free science independent of faith; practical theology, on the
other hand, is no science;
% --- p. 151 ---
it is a school for Christian upbringing, for existential knowledge, knowledge in faith and out of
faith.

As we now expressly remark that practical theology is no proper science, this remark will perhaps
make on many the impression that this theology in its unscientific character must therefore
necessarily stand at a far lower stage than the theoretical. That is a prejudice, a modern prejudice,
that hangs together quite naturally with the presupposition that science as science must everywhere in
knowledge's kingdom be the highest. But in relation to Christian upbringing, science as science is not
the highest, and cannot possibly be so, insofar as Christianity itself is higher than science. --- As
a school for the higher Christian formation, practical theology ought as a rule to have the
theoretical for its presupposition, i.e. for its negative foundation. Theoretical theology can then in
this respect be regarded as a learned preparatory school, a scientific school of preparation for the
practical.

That the matter must stand thus one can best convince oneself by setting the practical disciplines
alongside the corresponding theoretical ones. --- Insofar as theoretical theology occupies itself with
Christianity's facticity and the Christian phenomena, it must necessarily also go on occupying itself
with the Church and the ecclesial phenomena. But since all its investigations are hypothetical, its
results too will dissolve into nothing but hypothetical judgments.

% In the Danish print the closing „ mark of the second hypothetical judgment
% stands BEFORE its full stop: „... Bestemthed“. . . .  --- kept as printed
% there; the English sets it in the ordinary place.
``If Christianity has reality, the Church must also have reality.'' --- ``If the Church has reality,
it must be possible to show what it is and where it is, i.e. then the concept of the Church must be
capable of being fixed with decisive definiteness.'' . . . Conversely: ``If the concept of the Church
cannot be fixed with decisive definiteness, then the Church has no reality.'' --- ``If the Church has
no reality, then Christianity has no reality either.''

% --- p. 152 ---
Thus theoretical theology moves among nothing but hypothetical judgments; but from these judgments it
can nevertheless not come to draw a single definite conclusion, since it, as objective science, has no
power over faith and unbelief, and therefore everywhere must lack the categorical minor premise. Only
in practical theology is faith's categorical minor premise put in, and the hypothetical judgments can
then here furnish premises for hypothetical conclusions.

In practical theology one now infers thus:

``If Christianity has reality, then the Christian Church must also have reality;

Now Christianity has --- for us, the believers --- the highest, absolute reality;

% „Ergo“ is letterspaced in both syllogisms in the Danish; compare the unspaced
% „Endvidere:“ standing between them (verified at 400 dpi).
\emph{Ergo}: the Christian Church also has its true, fully valid reality.''

Further:

``If the Church has reality, then it must also be possible to show with definiteness what the Church
is and where it is, i.e. then the concept of the Church must be capable of being fixed and held fast
with decisive definiteness.

Now the Church has --- for us, the believers --- a true, decisive reality;

\emph{Ergo}: the concept of the Church must also let itself be fixed and held fast with decisive
definiteness.''

From this one sees how naturally theoretical theology introduces the practical, and therewith at the
same time how different they are. By the hypothetical major premise both schools are connected; but
the categorical minor premise fixes a yawning gulf between them. For theoretical theology the Church
is a hovering magnitude whose reality is hypothetical; for practical theology the Church stands fast
as a categorical authority, an objective power whose authority no one may doubt.

% --- p. 153 ---
The Church is then now, according to its concept, to be determined as Christianity factually present,
as that communion of the Spirit which in its Apostles' Creed, in its sacraments, and in its Holy
Scripture has preserved revelation's inspired original Word.

But although the Church in this way presents itself objectively, it is nevertheless just as impossible
to relate oneself objectively to the Church as it is, according to our foregoing, impossible to relate
oneself objectively to Christian truth. If one now again attempts to derive objective proofs for
Christianity from the Church's factual existence, then one falls back again into theoretical
theology's hypotheses, and must begin over. What the Church proclaims objectively, each individual
believer must appropriate subjectively: the Church lives only in the congregation's believing
subjectivity.

Practical theology therefore gets the great task of making clear the fundamental existential relation
that thus obtains between the Church's objective essence and the congregation's subjective faith. To
elucidate this by an example I refer again to Holy Scripture.

When with respect to the inspired Scripture we place the Church's objective authority in relation to
subjective faith, the religious tension at once sets in. This tension is in its immediate appearance so
strong that at first glance it looks as if all religious bonds of communion must at one stroke burst.

It is as though by the force of the consequences one were again thrown from the one extreme to the
other, from the extreme of a Catholic compelling objectivity that cows freedom with ``blind faith,'' to
the extreme of a Protestant unbridled subjectivity that chooses and rejects in subjective
arbitrariness. Superficially regarded, the relation must show itself thus; for we have in actuality set
the very strictest objectivity against the highest subjective freedom.

% --- p. 154 ---
What follows from this? By no means that the objective should now be softened in scientific mediation
with the subjective, or the subjective enfeebled in speculative mediation with the objective. By such a
mediation one assuredly attains no ``organic union''; one attains only a neutralization in which one
has neither the objective nor the subjective, but an unclear mixture in a ``third and higher.'' The
relation between the objective Church and subjective faith is existential, not scientific.

In the existential relation free subjectivity must strive its way forward in order to win faith; in this
strife of faith the subject must, so to speak, wrestle with the Church's inspired Scripture, as Jacob in
the valley wrestled with God. The subjectivity that in the wrestling is kindled with offense and rebels
against God's Word becomes an enemy of Christ, a freethinker who mocks Christianity and hates the
Church. The subjectivity that wrestles with Scripture and yet holds out in faith receives a new name and
a blessing when the morning dawns; but the strong wrestling at the same time sets its mark on him: he
limps his life through. --- He who will hear a human word and listen to a human counsel must first be
still; but he who is to hear God's Word and listen to God's counsel must be inwardly still.

In inward stillness, then, appropriation begins --- not of the whole Scripture, for there is surely no
one who can begin with the whole at once --- but of a single original word, a single passage, a single
Gospel. If there is no passage at all, no original word, no event in all Holy Scripture that finds entry
with subjectivity, then the individual still stands outside the Church, and here there can be no talk of
any understanding. If, on the other hand, the Church and Scripture have a word that makes its way through
the human being's inner life, but a single word, an original word that like ``a two-edged sword pierces
through, even to the dividing asunder of soul
% --- p. 155 ---
% DEFECT (Danish pp. 154--155): the quotation opened with „et tveegget Sværd is
% never closed --- there is no “ after „(Hebr. 4, 12)“ on the page, nor anywhere
% later in the paragraph. Kept as printed there; the English closes it after the
% reference, where the sense ends.
and spirit, and of the joints and marrow, and judges the thoughts and intents of the heart'' (Heb.\ 4:12)
--- if the Church and Scripture have but a single such word: then faith has begun, and with it
understanding.

Existential understanding is now not aimed at building a new theory of the Bible either on the Gadarene
demons or on Balaam's ass or on Noah's rainbow or on the six days of creation in the beginning; such
theories are now readily left to ``the scientific,'' who always have time enough to dispute about the
theories. The believer who is to appropriate the inspired Word has, as a believer, in faith no time to
dispute: every moment he thus disputes away in time he has lost for eternity. The believer who is to
appropriate revelation's Word has, as a believer, in faith no boldness to reason (the boldness to reason
with the Church about the separation of the divine and the human in Scripture's word is now readily
conceded to ``the scientific,'' who in witty theories seek witty diversion); the believing subjectivity
does not need theoretical diversion, but what he needs is --- to save his soul, to collect his mind on
the one thing needful.

As the many subjectivities now begin appropriation, each from its own point of view, there is here indeed
a diversity and therewith at the same time a subjective freedom; but this diversity is not dissipating,
this freedom is not arbitrary.

The believing subjectivities that, each in its own way, strive to appropriate Scripture's word all find
themselves under the same presupposition; they have all bowed before the confession of faith and the
sacraments; they have all humbled themselves under the Church's objective authority. The Christian spirit
of communion with its holy ideals then rests over all these subjectivities, so bound in faith, so set free
in faith. Each individual begins originally
% --- p. 156 ---
with his own particular, but each individual has in his particular at the same time the whole. Each
individual has begun the appropriation from his own peculiar point of departure; but each individual
strives nevertheless at the same time from his own peculiar point of departure toward the same goal, i.e.
the full appropriation of the order of salvation. Each individual is a one-sidedness in his limited
subjectivity, but the limited subjectivities have at the same time, each in itself, the original center,
and therefore move together about the one common center, that is: about the congregation's ideal
subjectivity.

On these conditions it now becomes possible to lay the foundation of a biblical interpretation that
exactly answers to the principle of the faith of the Gospel. Every attempt to explain away the wondrous and
to efface the distinctive is here rejected as un-Christian. The Christian interpretation of the Bible's
facts becomes recognizable in the particular by the Word's revealing power: here is the foundation for
practical hermeneutics.

Here in the Church, in Christ's kingdom, are the Spirit's various gifts; here are the means of grace; here
are possibilities and conditions for mutual instruction, admonition, and edification, as here are
possibilities and conditions for a true development of freedom in existential appropriation of faith's
Gospel.

In this communion of the Spirit, simplicity is compatible with the freest reflection. The relation between
the congregation and the Church is a relation of faith; the relation between the congregation's members
among themselves is a relation of faith: faith is in simplicity, as simplicity is in faith. But when
reflection enters in, there is here at the same time required the highest intelligence, the strictest
dialectic. The relation between the Church and the world is dialectical, in inversion: the Church is in
this world, and is yet not of this world; the Church is visible in the world and yet not visible to the
world; the Church is seen and yet not seen; this is the dialectical. The relation between the Church and
the congregation is dialectical, in inversion.
% --- p. 157 ---
The congregation's subjectivity builds on the Church as on its firm rock, its objective foundation, but the
Church is objective only for the congregation's faith. The believing subjectivity itself bears the
foundation by which it is borne: subjectivity bears the Church insofar as it is borne by the Church; this
is the dialectical. The relation between the individual and the congregation is dialectical, in continued
inversion. The brothers' and sisters' faith strengthens my faith, their hope my hope, their love my love: I
thank God, then, that I do not stand alone; for if I stood alone --- in the Church, eternity's powers would
terrify me, and holiness's ideals cast me down in the dust. But the brothers' and sisters' faith is
nevertheless there for me only by means of my own faith; their hope is there for me only insofar as it is my
hope; their love is there for me only insofar as I understand it in my love --- therefore: the objective is
there for me only insofar as I understand it subjectively, and I understand it subjectively only insofar as
I understand it in inversion, and precisely thereby understand it as faith's objective. So I am in inversion
everywhere referred back to myself, and, so to speak, driven into myself; so I shall myself bear the whole
congregation's faith in my faith, the whole congregation's hope in my hope, the whole congregation's love in
my love: so I thank God, then, that I am to stand alone, quite alone, although I stand in the midst of the
congregation; for only he who stands alone --- before God, stands on freedom's ground; only he who stands
alone before God becomes terrified before eternity's powers; only he who is terrified before eternity's
powers and bowed down with terror is raised up and strengthened with confidence.

Regarded from this point of view, practical theology will have many and great tasks to solve; everywhere it
comes upon the inverted relations of faith's existence, and yet these inverted relations are at the same time
to be explained in such a way that they can be understood in all
% --- p. 158 ---
simplicity, for faith is the simple. However much dialectical art practical theology may for the rest summon
up in order to solve these tasks: the last and most important thing in every solution is nevertheless always
this, that it at the same time dissolves itself. Only in dissolving itself in the solution of the tasks does
practical theology point to existence, actual existence, where everyone in particular is to solve his task on
his own responsibility. If practical theology neglects this thorough self-dissolution, then the theory comes
again, and with the theory scientific rigor, and with scientific rigor objective theology, and with objective
theology the weakening of the faith of the Gospel, and then everything is in confusion. This confusion must
then further be halted by a decisive protest; this decisive protest will then further sound still as before,
% Letterspaced in the Danish print (Sperrsatz) to the end of the lecture.
\emph{that Christianity is higher than science, and the faith of the Gospel heterogeneous with theology.}

%% ============================================================
%%  TWELFTH LECTURE (pp. 159--174)
%% ============================================================
\chapter*{Twelfth Lecture}
\addcontentsline{toc}{chapter}{Twelfth Lecture}
\label{ch:forel12}
\markboth{Twelfth Lecture}{}

\begin{center}\itshape
If practical theology does not understand how to remove its theoretical
semblance, this semblance will mislead religious formation still far more
than open theory does. Theoretical theology, which cannot deny its essential
difference from religious practice, will, by furthering its scientific
interests in the spirit of science, go on working fruitfully and formatively
upon all who in the inwardness of faith recognize that Christianity is
higher than science, and the faith of the Gospel heterogeneous with theology.
\end{center}

% --- p. 159 ---
The present age asserts, and rightly, that there is a difference between learning and cultivation, a
great difference between being in the strictest sense a learned man and being in the highest sense a
cultivated man. This distinction between learning and cultivation ought then also to come into
consideration when the talk is of theology's two-sided character and purpose. Theoretical theology has
its strength chiefly in learning; the practical, on the other hand, must use learning as a means for
personal cultivation.

As a learned science, theoretical theology then gets a double labor: partly to propagate and perfect
itself, partly to prepare practical cultivation. In the first respect theological science, like every
other science, has its purpose in itself, and ought to be cultivated for its own sake; in the last
respect it is to be used as a means for a purpose
% --- p. 160 ---
lying outside it; in the first respect it furnishes a high school for theologians, in the last a
preparatory school for pastors.

These two heterogeneous tasks may not be confused with one another; every such confusion occasions
disorder. If, for example, at the Academy of Art, where lectures are indeed given on the classical
works of antiquity, one were to go so thoroughly into learned investigations that out of sheer
thoroughness one forgot the hearers' interest, so that one at last lectured for archaeologists instead
of lecturing for artists, would the mistake not be conspicuous? Or if at the military academy, where
among other things lectures are indeed also given on chemistry, one were now to hold these lectures so
extraordinarily thoroughly that out of sheer thoroughness one forgot the aim, so that one was at bottom
lecturing not for warriors-to-be but for learned natural scientists: would the confusion not be
striking? In every endeavor it is a matter of holding fast to the aim, the ultimate οὑ ἕνεκα, without
% NOTE (from the transcription): the print sets the breathing alone over the
% upsilon (οὑ), without the circumflex of the usual classical „οὗ ἕνεκα“. Kept
% as printed in both files.
which the endeavor must necessarily fail. Just as a philologist or a historian must arrange his lecture
differently when his science is merely to be used as a higher means of cultivation than when he
lectures for such as cultivate the science for its own sake, so the theologian too ought surely to give
his lecture a different stamp when he lectures for students who use the theory only as a preparatory
school than when he lectures for theologians by profession.

That, for example, isagogics, exegesis, hermeneutics, church history, and so on are lectured with the
greatest possible fullness and the most thorough use of the whole literature belonging to them is in
order when learned lectures are given for theologians-to-be; but it is by no means in order when
propaedeutic lectures are given for pastors-to-be.

Propaedeutic lectures overloaded with too much material fail of their purpose: they take up too much
time, they
% --- p. 161 ---
weaken the hearers' interest, they occasion a hidden discontent on both sides. It must constantly seem
to the hearers as though the lecturer silently wished to force them against their interest to become
learned, and that in such subjects as ought really only to serve their cultivation; and it seems in turn
to the lecturer as though the hearers at the theoretical examination seldom or never furnished what he,
after the rich compass of his thorough lecture, should think himself entitled to demand. The consequence
is that the theoretical examination must be at one and the same time both too strict and too lax: too
strict, in that one at bottom demands more than one can demand in accordance with the real aim; too lax,
in that one must then on the other hand turn a blind eye to the candidate's weaknesses. Fate, chance,
and blind luck are then also apt to play a part in all such qualifying examinations.

These defects fall away when the propaedeutic lectures are laid out so as to correspond exactly to their
definite aim. The lectures given in the theological preparatory school on isagogics, exegesis,
hermeneutics, church history, and so on ought by no means to be superficial or to confine themselves to
insipid surveys and aphoristic remarks. On the contrary! precisely for that reason such lectures ought to
be given by learned and thorough theologians, because the learned and thorough theologian best
understands how to master the material, give the presentation its due stamp of character, and place the
detail in the right spot, without wearying the hearers with fumbling remarks and rambling investigations.
In this way the hearer would receive a fruitful impression of science's course and significance; in this
way the student who really studies would be enabled to use science as a means to a freer cultivation; in
this way the mass of knowledge would on the whole not become greater than that at the ensuing
% --- p. 162 ---
qualifying examination one could demand, and demand with inviolable strictness, that the candidate now
% NOTE: the Danish print has no comma after „Strenghed“; kept as printed there.
also had both the idea and the material in his power, and was able to develop each single point with the
greatest exactness.

To such courses in the learned disciplines philosophical studies would have to attach themselves in a
free and natural way. Suitable lectures on the history of philosophy, on metaphysics, and then
especially on rational ethics and the philosophy of religion would surely in this connection be very
fruitful.

Thus prepared in the theoretical preparatory school, the candidate now passes over into the practical
school, the institution of formation that shows him his vocation, his real purpose, his life's ideal
task --- the institution of formation that must therefore for him become the true and proper ``high
school.''

In this practical high school the various courses can again be laid out and arranged so that the
transition from the preparatory theory to the consummating practice takes place as naturally as
possible.

Practical theology must of course begin with the doctrine of the Church. In the doctrine of the Church,
the Church's history is at once to be conceived from a new principal side, namely from faith's point of
view. From this religious-ethical point of view the characters and events of church history appear
otherwise than they do to objective scientific contemplation. For the believing conception of the
Church's essence, symbolics and liturgics, the various types of doctrine with their corresponding rites,
appear otherwise than in the theoretical presentation. The entering hearer meets here everywhere the old
well-known facts; but these facts now show themselves to him at the same time in a new light; he sees
the old become new again; he feels that he must begin over: he feels himself rejuvenated and refreshed,
and begins with rising interest.

% --- p. 163 ---
On the doctrine of the Church follows the religious-ethical interpretation of Holy Scripture. Here too
the preliminary knowledge will come into application, here too the old will return: scientific
interpretation together with the philosophy of religion and rational ethics will now at many points
touch religious-ethical biblical interpretation; but the old becomes new in showing itself from a new
side: religious-ethical interpretation will now first in earnest maintain its distinctive character, and
Christian ethics be illuminated by rational ethics according to the well-known rule: \textit{opposita
juxta se posita magis illucescunt}.

To the interpretation of Scripture attach themselves catechetics and homiletics. One will already see
from this how imagined existence strives out beyond itself toward the actual. In the religious-ethical
interpretation of Holy Scripture the Christian concepts are still developed in merely imagined
existence, that is: without particular regard to the various human beings who are at once to translate
these concepts into life. Catechesis and preaching, on the other hand, presuppose straightforwardly the
presence of religious individuals, of individuals who must each for himself lay the revealed truth to
heart, understand it for his own use, understand it in practical appropriation.

So the tasks rise as one approaches actual life. As pastoral theology at last concludes with an
existential development of the fellowship of souls, the intimate relation that shall and ought to obtain
between the pastor and the congregation --- i.e. with everything that in a stricter sense belongs under
the cure of souls --- it at the same time points beyond itself, dissolves itself, and then sends the
young preacher out from the school into life, where he must then in turn begin over with caring for his
own soul, in becoming a curate of souls for others.

From this fleeting sketch it will surely be possible to discern how different practical education is
from theoretical
% --- p. 164 ---
pre-formation. This difference presses forward from yet another side.

In the theoretical school the hearer is really too passive, insofar as he is merely contemplative, and he
is contemplative insofar as he most immediately has only to think through and study up the disciplines
that are immediately lectured to him; in the practical school the candidate is active: here he can be
taught only in teaching himself as well; here he can appropriate the lecture only in trying himself as
well in the practice corresponding to the lecture's direction. He learns to preach only in alternately
hearing lectures on homiletics, hearing another preach, and therewith exercising himself in preaching; he
learns to catechize only in alternately hearing catechetics' direction, hearing another catechize, and
therewith exercising himself in catechizing.

The intimate union of theory and practice hereby presupposed brings it about at the same time that the
actual qualifying examination in the practical school must be undergone in quite another way than in the
theoretical. The theoretical examination settles everything at one stroke. After the hearer has for
several years sat in contemplative concealment as an intelligence half unknown to the teacher --- perhaps
even to himself --- there comes at last the day of examination, and this day of examination is then at one
stroke to unveil to him the yield he has won, and make him manifest to himself and the world. It is
otherwise with the practical examination. The material of learning here employed is not of so considerable
a compass that it gives a hearer much to do with mere learning by heart for an examination. The candidate
in the practical school is really, if you will, every day --- indeed every moment --- under examination.
The quality by which he is to know himself is then neither a strong memory nor a quick understanding ---
still less some single distinctive talent --- for it is the religious sense.
% --- p. 165 ---
If he lacks this sense --- that is, if the religious-ethical interest is not the fundamental interest
governing his whole life --- he will soon notice in himself that he does not belong in the circle of
teachers of religion, and will then also be able to withdraw in time in order to strike out on another
course. If, on the other hand, the candidate understands himself in the religious-ethical interest, if he
works at his personal formation with all diligence and vigilance, then he will also recognize that the
studies and exercises here laid down are precisely calculated to help him to the spiritual development he
so inwardly desires. The proof of this is not postponed to a total reckoning at the end, but expresses
itself in part every time he succeeds by a single test in displaying his diligence and ability. The
candidate then does not sit in contemplative concealment as a doubtful quantity whose hidden qualities
first come to light when the final experiment is undertaken; no, the candidate becomes from the beginning
manifest in relation to his striving: he gets from the beginning the opportunity to test himself piece by
piece, in that in his fellow students' presence he is piece by piece tested by his teachers.

If one now arranges these preliminary tests so that their result has due influence on the assessment of
the final concluding examination, then one has a measure of the candidate's ability, a measure that is
doubtless as objectively qualified as is well possible. ---

Thus far have I thought it right to go in making practical theology's method intuitable, in order to
determine its purpose and indicate its direction. But further I ought not to go in this connection either.
I have set myself the task of shedding light on a principle, but by no means the task of putting forward
any practical proposal or working at any outward alteration in the established order.

Why, in view of my attempts, there can in general be no
% --- p. 166 ---
thought of any outward alteration of the established order I must nevertheless explain a little more
closely in conclusion, lest the whole should at last dissolve into a great misunderstanding.

If the religious-ethical idea is to acknowledge an outward alteration, this alteration may not at all be
hastened by any outward consideration whatsoever, but must proceed from the idea itself alone. It is a
difficult task to realize this idea's demand; if the execution then failed, the blunder would only be the
more painful. It is almost unavoidable that practical formation, as I have here described it, should also
in its own way, at least provisionally, take on the semblance of a theory. This theoretical semblance can
be removed and done away with only by the most vigilant thinking, the most strenuous dialectic. If the
thinking that develops the theory does not at the same time have all the acuteness, maturity, and
superiority necessarily required to dissolve the theory again in all directions, practical theory will
retain its theoretical semblance, and this theoretical semblance will then work far more perniciously and
be far more dangerous to religious practice than theoretical theory.

Theoretical theory stands, in accordance with its scientific character, in a merely mediate relation to
actual existence. As this relation is now recognized for what it is, a merely mediate relation, the
semblance is at the same time removed, and existence \textit{eo ipso} freed from the theory. The
theoretically formed man soon notices that it is something quite other to work in life than to read in
school. Existence involuntarily compels him to begin over, and in beginning over he begins primitively
with discovering the ethical laws and the conditions for practical formation. If, on the other hand, he is
at the same time formed by a directly practical theory, a theory that gives itself the semblance of
explaining existence itself to him without nevertheless being able to free him from the theoretical that
clings to the explanation, then he is much
% --- p. 167 ---
exposed to confusing practical theory with actual practice, and finds his way with difficulty in
existence's freedom. All the one-sided methods, all the mannerisms and bad habits that are usually carried
over from the practical schools into practical life are as a rule explained by the fact that practical
theory has lacked dialectical force to remove its theoretical semblance and thrust the disciple away from
itself --- out into existence.

An outward alteration really calculated on the religious-ethical idea may never be hastened by selfish
impatience; such
% The Danish line ends „Utaalmo=“ and the next begins „modighed“, so the print
% doubles the „mo“ --- resolved in the transcription to the intended
% „Utaalmodighed“, as with the doubled „er“ on p. 111.
impatience is in evident contradiction with the idea itself. The religious-ethical principle is precisely
inwardness's principle, and that in the very highest potency. True inwardness knows itself in this world
always heterogeneous with any outward condition whatsoever. This heterogeneity is inwardness's burden; but
inwardness is strengthened precisely by bearing its burden. There is no danger at all that true inwardness,
where it has once been awakened, should be stifled and die under the pressure of the outward; on the
contrary, it is much to be feared that an all too happy inwardness should all too happily expire in an all
too happy harmony with its outward condition; such a happy harmony is always in this existence a more or
less unhappy delusion of sense: true harmony is only in the blessed life. Under every tension with the
outward, religious inwardness must know how to understand itself in its relation to eternity, and in this
its relation to eternity be sufficient to itself. Everywhere that inwardness has not yet learned to be
sufficient to itself, but restlessly and selfishly craves outward alterations, it is still only an immature
awakening and no true inwardness; but where true inwardness is not yet, there inwardness's true principle
cannot be recognized either.

If, then, anyone understands the religious-ethical principle in such a way that in this understanding he
now feels himself called and
% --- p. 168 ---
moved to prepare an outward objective defeat for objective theology, then he shows precisely thereby that
his outward-turned consciousness still lacks true inwardness; and in thus, in a lamentable lack of true
force of inwardness, setting all his hope on an outward alteration of conditions, he proves at the same
time that he has not understood the principle at all.

As regards these lectures in particular, I must expressly remark that I should have felt obliged to lay
them out on a quite different plan and to hold them in a different tone if it had been my intention to
propose any outward practical alteration. I should then have had to lay upon myself a threefold
obligation: partly to express myself at length on the particular manner in which theology is treated and
lectured at this University;
% NOTE: the Danish print reads „ved dette Universitetet“, with a double definite
% (this + the University), as kept there. English has no analogue; rendered
% plainly.
partly with respect to this to put forward my practical proposals in due detail; partly finally to treat
the whole with all the caution, circumspection, and forbearance that so highly important and difficult a
matter requires. None of these is here the case. I have not inveighed against any single theological
school in particular; what I have here uttered against the prevailing theology touches the individual
tendencies only insofar as it touches theology in general. Nor have I engaged in any final assessment of
the theological practice followed among us (the difficulty of such an assessment I fully acknowledge);
what I have said in the tenth lecture about the relation, or rather the disproportion, between theory and
practice is exclusively calculated to make intuitable the difference of principle between theoretical and
practical theology, but by no means calculated to hasten any outward, final alteration in the academic
order introduced among us. Finally, these lectures, which I shall now conclude, are not exactly held with
any particular regard to what one otherwise understands by caution, circumspection, and forbearance; on
the contrary!
% --- p. 169 ---
they are even in a certain sense --- I say it without regret --- they are in a certain sense held as
inconsiderately and unsparingly as possible.

Thus have I willed it with deliberate purpose. But I have nevertheless not willed it so because I for my
part was so inwardly dissatisfied with the outward established order that I then knew nothing else to do
than to see whether I might perhaps awaken the same dissatisfaction in others. Far from it! I acknowledge
with all readiness that in our theological institution, as it now stands, there stir both sound and able
forces, forces that can satisfy even a higher demand, if only the students study as they ought. I have, as
a university teacher, no calling to put forward any public judgment on my most reverend colleagues' official
activity; but on the other hand I have as a disciple of this University a right to express my gratitude to
my former teachers. The learned ability with which the various disciplines are lectured in our theological
faculty is assuredly so great that it will surely be able to enable the student to penetrate into science's
and cultivation's various directions, if only he himself seriously will, and does not --- what perhaps is
unhappily sometimes the case --- sink into critical peevishness and self-satisfied indolence.

It is no personal discontent with the outward established order that has occasioned my objections to the
prevailing theology. I have no ground for such discontent: no man has personally offended me; no theological
tendency stands hindering in the way of my endeavors. If I am to express a personal feeling, then I am for
my part inwardly grateful both to the University and to the state for my position, and find myself very
well in my academic relation to the studying youth. I strive to the best of my ability to do my duty; my
word has hitherto not lacked a well-disposed attention from the
% --- p. 170 ---
students' side: so I hope, then, on these terms to be able to attend to my own affairs, and feel not at all
called upon to meddle in others' --- which exceeds both my powers and my competence.

When I have nevertheless held these discourses as freely and inconsiderately as I hereby confess, then the
reason is simply this: that I have so much wished to forget all outward finite considerations in one single
exclusive consideration, in regard for the religious-ethical idea and the question of principle connected
with it.

With exclusive regard to the idea it is entirely necessary to set aside all other irrelevant considerations,
and that especially when it is a matter of laying bare the want of principle in existing theories. This
freedom and inconsiderateness is besides, at humanity's and cultivation's present stage, bound up with so
little danger that if anyone will do anything at all for the idea's sake, this is the very least he can do.

The inconsiderateness with which the principle disentangles itself from a prevailing delusion it then makes
good again by the self-limitation it so strictly imposes on every individual who might in any degree feel an
inclination to serve the idea: that he may not at all strive for any outward purpose or think of winning any
outward victory. --- But inwardness's law is nevertheless not that the inner should be indifferent to the
outer, or a conviction of truth indifferent to its utterance. There can and must be striving for the idea's
right, only that one always strive with the idea's weapons, and use for it nothing but ideal forces.

Should it ever come so far with this high school that both teachers and hearers at last became so weary of
the idea that they did not even care to contend either for or against a principle: would this not then be a
sorry sign that the spirit had
% --- p. 171 ---
departed, and the whole instruction sunk down into an empty mechanism?

When the question is of the relation of principle between faith-knowledge and objective knowing, then the
philosopher is surely as entitled to have his say as the theologian; and when a theological theory is charged
with want of principle, then there are, so far as I can judge, only two ways out that satisfy the idea: the
one is to come out with a vigorous defense, if the attack is unwarranted; the other is to admit one's fault,
if the attack is warranted; for the third way out --- namely to hush up the whole investigation and let the
attack blow over, as a hailstorm blows over --- is only worldly prudence's evasion, an evasion that can
easily end with one's at last having to evade the idea.

In order now not to become an accomplice in such an evasion, I have here expressed myself as plainly as
possible, feeling convinced that at this University one could still set a question of principle in motion
without this inner ideal motion's occasioning any outward disturbance whatsoever in any way.

So I conclude these lectures with the confidence that, even if on one point or another they may have offended
some, they shall nevertheless not confuse anyone or seduce anyone into untimely despondency over the existing
theories.

It goes with the dazzling theories as with the insubstantial semblance: they infatuate only him who cannot
see through them. He who sees through the semblance is not deceived by the semblance; he who sees through the
theory is not infatuated by the theory. But when the theories have ceased to infatuate us, they first begin in
a higher sense to form us. We then learn to study the theory with inner freedom and therewith at the same time
with earnestness and thoroughness. When the theory has then done its part, existence comes to our aid with our
personal formation, and does its part.

% --- p. 172 ---
Theology's young devotees can then still, as before, study their theological sciences with pleasure and zeal;
they can study --- the freethinker's attack, the apologist's defense, the interpreter's rules, the
dogmatician's speculations --- they can study it all, and it shall all do them good, if only under all these
studies they will at the same time remember what lies in this thought, this one simple thought, \emph{that
Christianity is higher than science, and the faith of the Gospel heterogeneous with theology.}

% Printed here: a short centred double rule (two thin tapered lines) closing
% the twelfth lecture, roughly a third of the way down the otherwise blank
% remainder of p. 172.
\begin{center}
\rule{0.30\textwidth}{0.4pt}\\[2pt]
\rule{0.30\textwidth}{0.4pt}
\end{center}

%% ============================================================
%%  EFTERSKRIFT / POSTSCRIPT (pp. 173--174)
%%  Set smaller than the lectures in the print; heading "Efterskrift."
%% ============================================================
% --- p. 173 ---
\chapter*{Postscript}
\addcontentsline{toc}{chapter}{Postscript}
\label{ch:efterskrift}
\markboth{Postscript}{}

\begingroup\small

When I began to work out and deliver these lectures, I had (as the Preface indeed shows) a kind of
feeling that a prolix, protracted, and momentous conflict surely lay before us about ``the great problem
of the age,'' about ``faith's relation to knowledge,'' ``knowledge's relation to faith''; but ---
strangely enough! --- just as I have now finished the revision of the work and see my Preface in print,
the thought suddenly dawns on me that here at bottom there is nothing at all to dispute about.

The question of faith's relation to knowledge has by two masters preeminent in dialectic, Johannes
\textit{de silentio} and Johannes Climacus, already been sufficiently answered in the one proposition:
\emph{Faith is a paradox}. To sit down and wait for new contributions, further explanations, personal
declarations, and other such things before one delivers one's final verdict on this point may indeed be
as tempting to the spectators' indolence as it accords with their caution --- they who attend not so much
to the matter as to the circumstances; but it is nevertheless true that he who cannot now see from the
documents already before him what the principle is will hardly ever come to see it.

It stays, then, in short, with the proposition: \emph{Faith is a paradox}. But however contentious this
proposition may look at first glance, it is least of all a literary proposition of dispute. If one
accepts the proposition, there is nothing here to dispute about. If one refuses to accept the proposition,
it must surely be either because one does not understand it rightly or because one understands it as
incorrect: in neither of these cases is there here material for dispute.

If one dares not accept the proposition because one does not understand it rightly, it is assuredly far
more sensible to turn to the dialecticians named above and in all quietness make use of their guidance
than at once to take the field with unassailable opinions and doubtful objections, and enter into a
dispute about something one does not yet rightly understand.

If one rejects the proposition because one understands (or at any rate supposes one understands) it as
incorrect, the endeavor will of course be directed at
% --- p. 174 ---
explaining away ``the paradox'' and conserving faith; but neither shall this endeavor be able to occasion
dispute. One can really, if only one seriously will, quite readily explain away the paradox without much
bother with dialectic; the supervising dialectic makes in that case no objections, but merely sees to it
that the undertaking really succeeds --- that is to say, that the explaining-away is done in such a way
that both Christianity and Judaism, both the apostles' and Abraham's faith, are exactly and consistently
--- explained away at the same time.

The faith that then remains when the paradox has been explained away is therefore a paradox-free faith, a
pure, general, essential faith of reason common to all human beings, a faith of reason that every
reasonable, nobly thinking, and cultivated man can be known to defend without exposing himself to
``science's'' taking it ill and accusing ``the knight of faith'' of superstition. --- In what
% NOTE (from the transcription): the opening mark before „Troesridderen“ is set
% in the Danish with a pair of upright low commas rather than the book's usual
% low-high sort; the same word takes the ordinary sort three lines further down.
% Transcribed uniformly there; nothing to reproduce here.
category the faith of reason, or the paradox-free reasonable faith, is for the rest to be conceived I
shall however not be able to say, since, frankly spoken, I do not know it. There is, if I remember
rightly, a German novel entitled ``Der Bräutigam ohne Braut''; better still would it be if, on the
occasion of the pure paradox-free faith of reason accessible to ``science,'' one could at last get a
% The Danish print has a full stop after „paradoxfrie“ where the sense requires
% a comma --- printer's error, kept as printed there. The English has the comma.
Danish novel entitled ``The Knight of Faith without Faith.'' But suppose even that so remarkable a result
% Nielsen's imagined Danish title is „Troesridderen uden Tro“. It is Englished
% here, unlike the real German title above, because the whole point of it is the
% sense; „Troesridder“ is Kierkegaard's knight of faith, from Frygt og Bæven.
were attained --- a knight of faith without faith would after all cut a strange, almost paradoxical figure
in a literary dispute about faith.

If one finally, on at last seeing oneself compelled to concede the proposition's unconditional validity,
were so resolute as to give up the religious and cast away all faith, in order at least in that way to be
rid of the paradox, the matter would at one stroke take so serious a turn that one could then least of all
think of accomplishing anything by learned disputes about faith.

To dispute about faith is human; but to strive for faith --- yes, that is in truth divine. Happy is he who
then knows how to strive so that he does not dispute; blessed is he to whom the Spirit granted to strive
as one ought --- the true divine strife in faith for faith!

\endgroup

% The book ends here. Printed below the last line of p. 174 is a single centred
% typographic fleuron --- a horizontal stroke flaring into two pointed leaves
% that flank a small open ring. There is no colophon, no printer's imprint and
% no „Ende“; PDF 190--191 are the Indhold.
% ⚠ OPEN ITEM, carried from the transcription: the fleuron is stood in for by a
% plain rule in BOTH files. If it is ever replaced by a real glyph (pifont, or a
% ❦-class sort), change it in transcription.tex and here in the same pass, or
% the two editions will stop matching.
\begin{center}
\rule{0.22\textwidth}{0.4pt}
\end{center}

%% ============================================================
%%  INDHOLD / CONTENTS (PDF 190--191)
%%  Carries the Indhold's OWN readings where they differ from the lecture
%%  heads (VIII "half-speculative", X "gjældende"/"; thi"). Do not normalise.
%% ============================================================
\backmatter
\chapter*{Contents}
\addcontentsline{toc}{chapter}{Contents}
\markboth{Contents}{}

\begin{center}
\rule{0.12\textwidth}{0.4pt}
\end{center}

\newcommand{\indpost}[2]{%
  \par\medskip
  \begin{center}\bfseries #1\end{center}
  \par\nobreak\noindent #2\par}

%  ⚠ THIS SECTION CARRIES THE INDHOLD'S OWN READINGS, which in two places
%  genuinely differ from the lecture heads in the 1850 printing. Verified at
%  600 dpi on both witnesses. DO NOT normalise them against the chapter blocks
%  above:
%    - VIII: Indhold „Det halvspeculative Christusbillede“  -> "half-speculative"
%            head p. 97 „Det halvphilosophiske Christusbillede“ -> "half-philosophical"
%      This variant IS visible in the English, and is the one to check first if
%      anyone "tidies" the two into agreement.
%    - X:    Indhold „gjældende“ … „fra Theorien; thi“
%            head p. 129 „giældende“ … „fra Theorien, thi“
%      Only HALF of this variant survives translation: the semicolon/comma is
%      reproduced (semicolon here, comma in the chapter block above), but
%      „gjældende“/„giældende“ is a purely orthographic difference with no
%      English analogue --- both render "in force". So the English witnesses
%      cannot carry it, and the Danish is the only place it can be seen.
%
%  The page ranges are the Indhold's own. NOTE for anyone working from older
%  notes: lecture VI's range here reads 66--80, NOT 66--86. The "66--86" was an
%  error in the scan's ABBYY text layer, corrected during Phase 1; 66--80 agrees
%  with lecture VII opening on p. 81. There is nothing to reproduce faithfully
%  here --- the print is simply right.

\indpost{First Lecture.}{Introductory considerations. Christianity is higher
than science, and the faith of the Gospel heterogeneous with theology. Pp.
1--13.}

\indpost{Second Lecture.}{In its conception of the relation between historical
and religious faith, theology must separate what in religion is inseparable, and
therefore the faith of the Gospel is heterogeneous with theology. Pp. 14--25.}

\indpost{Third Lecture.}{In its conception of the confession of faith and of
Holy Scripture, scientific theology must consistently go on working against
faith; for the faith of the Gospel is heterogeneous with theology. Pp. 26--37.}

\indpost{Fourth Lecture.}{For scientific theology the Bible is a historical
document of antiquity whose heterogeneous constituents must submit to
ratiocinating criticism; for the believing conception the Bible's wondrous
content is exalted above criticism, since it is impossible for the observer to
relate himself critically and ratiocinatingly to the miracle: the faith of the
Gospel is heterogeneous with theology. Pp. 38--51.}

\indpost{Fifth Lecture.}{Theological interpretation of Scripture misconstrues
the peculiar character of revelation, and dissolves the content of faith into
science's explanation; religious interpretation of Scripture reduces science to
a means for a believing appropriation of the revealed Word: the faith of the
Gospel is heterogeneous with theology. Pp. 52--65.}

% The page range here settles the old query: 66--80, not 66--86.
\indpost{Sixth Lecture.}{The theological interpreter of Scripture appeals to the
Bible's inspiration in the whole, but may on no account acknowledge its
inspiration in the particular. Since there is thus doubt about every particular,
there must of course also be doubt about the whole: for the theologian the
Bible's inspiration is an empty phrase. The evangelical interpreter of Scripture,
who in accordance with the law of existence everywhere sees the whole
concentrated in the particular, conceives every single passage of Scripture as
inspired, under the ecclesial presupposition that the whole Bible is inspired:
the faith of the Gospel is heterogeneous with theology. Pp. 66--80.}

\indpost{Seventh Lecture.}{In the treatment of sacred history --- ``the Life of
Jesus'' --- the theologian will work against the freethinker; but the scientific
criticism on which he proceeds shows plainly enough that in principle he is at
one with the freethinker. The believer, who in his conception of the Gospels'
Christ holds exclusively to religious-ethical criticism, is in principle at one
with the evangelists: the faith of the Gospel is heterogeneous with theology.
Pp. 81--96.}

% VARIANT: „halvspeculative“ here -> "half-speculative"; the head on p. 97 reads
% „halvphilosophiske“ -> "half-philosophical". Both stand. Do not reconcile.
\indpost{Eighth Lecture.}{The half-speculative image of Christ that the
speculating dogmatician conjures up consists in a mystical fusion of the
scientific and the religious, and is on that account so infinitely different
from the evangelists' Christ that it falls even outside the religious-ethical
horizon: the faith of the Gospel is heterogeneous with theology. Pp. 97--111.}

\indpost{Ninth Lecture.}{Theology's half-speculative dogmatics cannot justify
its ecclesial significance from the history of the Church, since the history of
the Church proves, on the contrary, how dogmatic prejudices have in this respect
harmed the Church, and what a thoroughgoing opposition there is here between
theological and ecclesial dogmatics. Theological dogmatics must fall before the
proposition that Christianity is a paradox: the faith of the Gospel is
heterogeneous with theology. Pp. 112--128.}

% VARIANT: the Indhold's semicolon before „thi“ is reproduced here; the head on
% p. 129 has a comma, and is set that way in the chapter block above. The
% Indhold's „gjældende“ against the head's „giældende“ cannot be carried into
% English --- both are "in force" --- so that half of the variant lives only in
% the Danish.
\indpost{Tenth Lecture.}{According to the method now in force, theoretical
theology finds itself in an evident disproportion to religious practice; in
order to maintain the principle of faith, the formation of pastors must
emancipate itself from theory; for the faith of the Gospel is heterogeneous with
theology. Pp. 129--143.}

\indpost{Eleventh Lecture.}{Theoretical theology must relate itself skeptically
to Christianity's revealed truth; practical theology, which in imagined
existence labors to solve the Christian tasks, must end by dissolving itself:
the faith of the Gospel is heterogeneous with theology. Pp. 144--158.}

\indpost{Twelfth Lecture.}{If practical theology does not understand how to
remove its theoretical semblance, this semblance will mislead religious
formation still far more than open theory does. Theoretical theology, which
cannot deny its essential difference from religious practice, will, by
furthering its scientific interests in the spirit of science, go on working
fruitfully and formatively upon all who in the inwardness of faith recognize
that Christianity is higher than science, and the faith of the Gospel
heterogeneous with theology. Pp. 159--174.}

\medskip
\begin{center}
\rule{0.30\textwidth}{0.7pt}
\end{center}

\end{document}
